% \iffalse 元注释
% 
% 本文件为 `caption.dtx'。
% 
% 版权所有 (C) 1994-2023 Axel Sommerfeldt (axel.sommerfeldt@f-m.fm)
% 
% --------------------------------------------------------------------------
% 
% 本作品可在 LaTeX 项目公共许可证（LaTeX Project Public License）
% 1.3 版或其后续版本所规定的条件下分发与修改。该许可证的最新版本见：
%   http://www.latex-project.org/lppl.txt
% 1.3 或更高版本属于 2003/12/01 或更高版本 LaTeX 发行版的一部分。
% 
% 本作品在 LPPL 下的维护状态为 “maintained（受维护）”。
% 
% 本作品当前维护者为 Axel Sommerfeldt。
% 
% 本作品包含以下源文件：
%   caption.ins, caption.dtx, caption-light.dtx, caption2.dtx, caption3.dtx,
%   caption-ams-smf.dtx, caption-beamer.dtx, caption-elsarticle.dtx,
%   caption-koma.dtx, caption-memoir.dtx, caption-ntg.dtx,
%   caption-thesis.dtx, bicaption.dtx, ltcaption.dtx, subcaption.dtx,
% 以及派生文件：
%   caption.sty, caption-light.sty, caption2.sty, caption3.sty,
%   caption-ams-smf.sto, caption-beamer.sto, caption-elsarticle.sto,
%   caption-koma.sto, caption-memoir.sto, caption-ntg.sto,
%   caption-thesis.sto, bicaption.sty, ltcaption.sty, subcaption.sty.
% 
% \fi
%
% \CheckSum{2925}
%
% \iffalse
%<*driver>
\NeedsTeXFormat{LaTeX2e}[1994/12/01]
\ProvidesFile{caption.drv}[2023/09/07 v3.6 宏包 caption（中文译本）]
%\errorcontextlines=3
%
\documentclass{captiondoc-cn}
\hypersetup{pdfkeywords={LaTeX, 宏包, caption, 中文}}
%
\usepackage[listof=false,hypcap=false]{caption}[2022/01/06]
%
\usepackage{graphicx,longtable}
%
% \ContinuedFloat
\DeclareCaptionLabelFormat{continued1}{续 #1~#2}
\DeclareCaptionLabelFormat{continued2}{#1~#2\alph{ContinuedFloat}}
% \DeclareCaptionLabelFormat
\DeclareCaptionLabelFormat{bf-parens}{(\textbf{#2})}
% \DeclareCaptionStyle
\DeclareCaptionStyle{mystyle}[margin=5mm,justification=centering]%
                    {font=footnotesize,labelfont=sc,margin={10mm,0mm}}
% 示例 1
\DeclareCaptionFormat{myformat1}{#1#2\\#3}
\newlength\myindention
\DeclareCaptionFormat{myformat2}{#1#2\\\hspace*\myindention#3}
\DeclareCaptionOption{myindention}{\setlength\myindention{#1}}
% 示例 2
\DeclareCaptionFormat{reverse}{#3#2#1}
\DeclareCaptionLabelFormat{fullparens}{(\bothIfFirst{#1}{~}#2)}
\DeclareCaptionLabelSeparator{fill}{\hfill}
% 示例 3
\DeclareCaptionFormat{llap}{\llap{#1#2}#3\par}
\DeclareCaptionFormat{llapx}{\llap{\makebox[5em][l]{#1}}#3\par}
% 示例 4
\DeclareCaptionLabelFormat{andtable}{#1~#2 \& \tablename~\thetable}
%
\begin{document}
  \DocInput{caption.dtx}
\end{document}
%</driver>
% \fi
%
% \iffalse
% ----------------------------------------------------------------------------------------------- %
% \fi
%
% \def\thispackage{宏包 \pkg{caption}}
% \def\Thispackage{宏包 \pkg{caption}}
%
% \def\x{\-}
%
% \newenvironment{Annotation}%
%   {\ifvmode\else\unskip\par\fi\pagebreak[2]%
%    \smallskip
%    \small(\ignorespaces}{\unskip)\par}
% \newenvironment{Annotation*}%
%   {\ifvmode\else\unskip\par\fi\pagebreak[2]%
%    \small(\ignorespaces}{\unskip)\par}
%
% \newenvironment{Example}%
%   {\ifvmode\else\unskip\par\fi\pagebreak[2]%
%    \minipage{\linewidth}\smallskip}%
%   {\smallskip\endminipage\par}
%
% \makeatletter
% \newcommand\example{\@ifstar
%   {\@example{belowskip=\abovecaptionskip}}%
%   {\@example{}}}
% \def\@example#1{%
%   \@testopt{\@@example{#1}}{figure}}
% \long\def\@@example#1[#2]#3#4{%
%   \begingroup
%     \captionsetup{#1,size=small,margin={\leftmargini,10pt},#3}%
%     \captionof{#2}[]{#4}%
%   \endgroup}
% \makeatother
%
% \newenvironment{Expert}%
%   {\ifvmode\else\unskip\par\fi\pagebreak[3]%
%    \smallskip
%    \footnotesize\ignorespaces}{\par}
% \newenvironment{Expert*}%
%   {\ifvmode\else\unskip\par\fi\pagebreak[3]%
%    \footnotesize\ignorespaces}{\par}
%
% \newenvironment{Note}[1][\DefaultNoteText]%
%   {\ifvmode\else\unskip\par\fi
%    \smallskip
%    \small\emph{#1:}~\ignorespaces}{\par}
% \newenvironment{Note*}[1][\DefaultNoteText]%
%   {\ifvmode\else\unskip\par\fi
%    \small\emph{#1:}~\ignorespaces}{\par}
% \providecommand*\DefaultNoteText{注}
%
% \makeatletter
% \def\see{\@ifstar{\SEE\relax{参见}}{\SEE\small{参见}}}
% \def\See{\@ifstar{\SEE\relax{参见}}{\SEE\small{参见}}}
% \newcommand*\SEE[3]{\nopagebreak{#1(#2 #3)}}
% \newcommand*\SeeUserDefined[1][]{\See{\longref{declare}#1}}
% \makeatother
%
% \iffalse
% ----------------------------------------------------------------------------------------------- %
% \fi
%
% \GetFileInfo{caption.drv}
% \let\docdate\filedate
% \let\docversion\fileversion
% \GetFileInfo{caption.sty}
%
% \title{浮动环境题注的自定义排版%^^A
%         \texorpdfstring{\thanks{%^^A
%          本宏包版本号为 \docversion。}}%^^A
%        {}}
% \author{Axel Sommerfeldt\\
%         \url{https://gitlab.com/axelsommerfeldt/caption}}
% \date{\docdate}
% \maketitle
%
% \begin{abstract}
% \Thispackage\ 用于自定义 |figure|、|table| 等浮动环境中的题注，
% 并能与众多其他宏包协同工作。
% \par\smallskip
% \textit{请注意：}许多文档类本身已内置了用于自定义题注的选项与命令。
% 如果这些功能对你已经足够，通常完全没有必要再使用 \pkg{caption} 宏包。
% 若你只需要使用 \cs{captionof} 命令，加载体积极小的 \pkg{capt-of} 宏包
% 通常就已足够。
% \end{abstract}
%
% \newcommand\shortexampletext{%
%   洁白的沙滩，粉色光滑的海螺贝壳，一片充满可能性的大海。免税店里摆满了
%   欧洲最精美的礼品与香水。}
% \newcommand\exampletext{\shortexampletext\space
%   在众多锦标赛级球场之一的热带绿茵上，挥出你最喜爱的那一杆高尔夫。}
%
% \iffalse
% --------------------------------------------------------------------------- %
% \fi
%
% \section*{简介}
%
% 在标准 \LaTeX\ 文档类中，题注并没有得到应有的重视：它只是像普通段落
% 一样排版，与正文其余部分在视觉上没有明显区别，就像下面这样：
%
% \example*{size=normalsize,margin=0pt}{\shortexampletext}
%
% 我们应当有办法改变这种状况；例如，可以让题注文字比正文略小一些、
% 增加额外边距、让题注标签采用与标题相同的字体族和字形，等等，就像
% 下面这样：
%
% \example*{size=small,margin=10pt,labelfont=bf,labelsep=endash}{\shortexampletext}
%
% 借助本宏包即可轻松做到这一点：它提供了大量开箱即用的题注排版选项，
% 同时你也完全可以自定义属于自己的样式。
%
% \bigskip
%
% \INFO\begin{minipage}[t]{\textwidth}
% \small
% 请注意，\thispackage\ 只控制题注的外观与风格，\emph{不}控制题注的
% 放置位置。（若要控制位置，可以使用 \pkg{floatrow} 等其他宏包\cite{floatrow}。）
% \end{minipage}
%
% \medskip
%
% \INFO\begin{minipage}[t]{\textwidth}
% \small
% 本文档已显陈旧，目前正在修订之中。修订进度可在此查看：
% \url{https://gitlab.com/axelsommerfeldt/caption/-/commits/doc/caption}
% \end{minipage}
%
% \iffalse
% --------------------------------------------------------------------------- %
% \fi
%
% \clearpage
% \tableofcontents
%
% \iffalse
% --------------------------------------------------------------------------- %
% \fi
%
% \clearpage
% \section{使用本宏包}
% \label{usage}
%
% 在文档导言区（即 |\document|\-|class| 与 |\begin{doc|\-|u|\-|ment}|
% 之间的部分）插入
% \begin{quote}
%   |\usepackage|\oarg{选项}|{caption}[|\texttt{\filedate}|]|
% \end{quote}
% 这些选项用于控制题注的外观；例如：
% \begin{quote}
%   |\usepackage[margin=10pt,font=small,labelfont=bf,|\\
%   |            labelsep=endash]{caption}|%
%   %|[|\texttt{\filedate}|]|
% \end{quote}
% 得到的题注将与简介中第二个示例的样式相同。
%
% \DescribeMacro{\captionsetup}
% 若要在之后修改选项，\thispackage\ 提供了命令
% \begin{quote}
%   |\captionsetup|\oarg{浮动体类型}\marg{选项}\quad。
% \end{quote}
% 因此
% \begin{quote}
%   |\usepackage[margin=10pt,font=small,labelfont=bf]{caption}|
% \end{quote}
% 与
% \begin{quote}
%   |\usepackage{caption}|\\
%   |\captionsetup{margin=10pt,font=small,labelfont=bf}|
% \end{quote}
% 效果完全相同。
%
% 需要了解的是，|\caption|\-|setup| 只对当前环境生效。因此，若只想修改
% 当前 |figure| 或 |table| 的设置，只需把 |\caption|\-|setup| 命令放在
% 该 |figure| 或 |table| 内部、|\caption| 命令之前即可。例如：
% \begin{quote}
%   |\begin{figure}|\\
%   |  |\ldots\\
%   |  \captionsetup{singlelinecheck=off}|\\
%   |  \caption{|\ldots|}|\\
%   |\end{figure}|
% \end{quote}
% 仅关闭了这一个 |figure| 的单行检查，其余题注不受影响。
%
% \begin{Annotation}
% 有关 |\caption|\-|setup| 的详细说明参见 \longref{captionsetup}。
% \end{Annotation}
%
% \iffalse
% --------------------------------------------------------------------------- %
% \fi
%
% \clearpage
% \section{选项}
% \label{options}
%
% \def\OptionLabel{RaggedRight}
% \def\UserDefined{\ldots}
%
% \iffalse
% --------------------------------------------------------------------------- %
% \fi
%
% \subsection{排版格式}
% \label{formats}
%
% \DescribeMacro{format=}
% 插图或表格的题注主要由三部分组成：题注标签（标明该对象是“插图”还是
% “表格”以及对应的编号）、题注文字（通常是对内容的简短描述），以及
% 分隔标签与文字的题注分隔符。
%
% \textit{题注格式}决定上述信息的呈现方式，由选项
% \begin{quote}
%   |format=|\meta{格式名称}\quad,
% \end{quote}
% 指定，其参数为题注格式的名称。
%
% 内置两种标准题注格式：
% \begin{Options}{\OptionLabel}
%   \item[plain]%\NEWdescription{v3.0h}
%   将题注排版为普通段落。
%
%   \item[hang]
%   缩进题注文字，使文字“悬挂”在第一行之下。
%
%   \item[\UserDefined]
%   也可使用 |\Declare|\-|Caption|\-|Format| 自定义格式。\LineBreak
%   \SeeUserDefined
% \end{Options}
%
% \begin{Example}
%   示例：指定选项
%   \begin{quote}
%     |format=hang|
%   \end{quote}
%   得到如下题注：
%   \example{format=hang}{\exampletext}
% \end{Example}
%
% \pagebreak[3]
% \DescribeMacro{indention=}
% 对两种格式（\texttt{plain} 与 \texttt{hang}）都可以设置从题注第二行
% 起附加缩进量，使用选项
% \begin{quote}
%   |indention=|\meta{长度}\quad。
% \end{quote}
%
% 下面是两个示例：
%
% \begin{Example}
%   \begin{quote}
%     |format=plain,indention=.5cm|
%   \end{quote}
%   \captionsetup{skip=0pt}
%   \example{format=plain,indention=.5cm}{\exampletext}
% \end{Example}
%
% \begin{Example}
%   \begin{quote}
%     |format=hang,indention=-0.5cm|
%   \end{quote}
%   \captionsetup{skip=0pt}
%   \example{format=hang,indention=-0.5cm}{\exampletext}
% \end{Example}
%
% \bigskip
%
% \pagebreak[3]
% \DescribeMacro{labelformat=}
% 使用选项
% \nopagebreak[3]
% \begin{quote}
%   |labelformat=|\meta{标签格式名称}
% \end{quote}
% \nopagebreak[3]
% 可以指定题注标签的排版方式。内置五种标准题注标签格式：
%
% \begin{Options}{\OptionLabel}
%   \item[original]
%   题注标签按文档类规定的方式排版，通常为名称加编号（类似 \texttt{simple}）。
%
%   在 \thispackage\ \version{3.4} 之前，该标签格式名为 `default'；
%   自 \version{3.5} 起更名为 `original'，通常也是默认标签格式。因此
%   |label|\-|format=|\-|default| 仍会选中该格式，除非默认值已被文档类
%   或 babel 语言宏包更改。
%
%   \item[empty]
%   题注标签留空。
% \iffalse
%   （该选项与其他选项（如 \texttt{labelsep=none}）搭配使用时才有意义。）
% \fi
%
%   \item[simple]
%   题注标签排版为名称加编号。
%
%   \item[brace]\NEWfeature{v3.1f}
%   题注标签末尾加一个右花括号。
%
%   \item[parens]
%   题注标签中的编号放在圆括号内。
%
%   \item[\UserDefined]
%   也可使用 |\Declare|\-|Caption|\-|Label|\-|Format| 自定义标签格式。
%   \SeeUserDefined
% \end{Options}
%
% \begin{Example}
%   示例：使用选项
%   \begin{quote}
%     |format=plain,labelformat=parens,labelsep=quad|
%   \end{quote}
%   得到如下题注：
%   \example{format=plain,labelformat=parens,labelsep=quad}{\exampletext}
% \end{Example}
%
% \medskip
%
% \begin{Note*}
% 某些环境（例如 \pkg{algorithm2e} 宏包提供的 |algorithm| 环境）
% 可能对题注标签格式的更改较为敏感。
% \end{Note*}
%
% \bigskip
%
% \pagebreak[3]
% \DescribeMacro{labelsep=}
% 使用选项
% \begin{quote}
%   |labelsep=|\meta{分隔符名称}
% \end{quote}
% 指定所使用的题注分隔符。\footnote{%
%   若题注标签或题注文字为空，则不使用分隔符。}
% 可在以下值中选择：
%
% \begin{Options}{\OptionLabel}
%   \item[none]
%   不使用题注分隔符。
% \iffalse
%   （该选项与其他选项（如 \texttt{labelformat=empty}）搭配使用时才有意义。）
% \fi
%
%   \item[colon]
%   题注标签与文字之间以冒号加空格分隔。
%
%   \item[period]
%   题注标签与文字之间以句点加空格分隔。
%
%   \item[space]
%   题注标签与文字之间以单个空格分隔。
%
%   \item[quad]
%   题注标签与文字之间以 |\quad| 分隔。
%
%   \item[newline]
%   题注标签与文字之间以换行（|\\|）分隔。请注意，该分隔符并非与所有
%   题注格式都能搭配（例如 |format=|\x|hang|）；强行组合这样的选项会
%   产生错误信息。
%
%   \item[endash]\NEWfeature{v3.0h}
%   题注标签与文字之间以两端带空格的短破折号分隔（\verb*| -- |）。
%
%   \item[\UserDefined]
%   也可使用 |\Declare|\-|Caption|\-|Label|\-|Sep|\-|a|\-|ra|\-|tor| 自定义分隔符。
%   \SeeUserDefined
% \end{Options}
%
% 下面是三个示例：
% \begin{Example}
%   \begin{quote}
%     |format=plain,labelsep=period|
%   \end{quote}
%   \captionsetup{skip=0pt}
%   \example{format=plain,labelsep=period}{\exampletext}
% \end{Example}
%
% \begin{Example}
%   \begin{quote}
%     |format=plain,labelsep=newline,singlelinecheck=false|
%   \end{quote}
%   \captionsetup{skip=0pt}
%   \example{format=plain,labelsep=newline,singlelinecheck=false}{\exampletext}
% \end{Example}
%
% \begin{Example}
%   \begin{quote}
%     |format=plain,labelsep=endash|
%   \end{quote}
%   \captionsetup{skip=0pt}
%   \example{format=plain,labelsep=endash}{\exampletext}
% \end{Example}
%
% \bigskip
%
% \pagebreak[3]
% \DescribeMacro{textformat=}
% \NEWfeature*{v3.1}
% 使用选项
% \nopagebreak[3]
% \begin{quote}
%   |textformat=|\meta{文字格式名称}
% \end{quote}
% \nopagebreak[3]
% 可以指定题注文字的排版方式。内置三种标准题注文字格式：
%
% \begin{Options}{\OptionLabel}
%   \item[empty]
%   不显示题注文字。
%
%   \item[simple]
%   题注文字按原样排版。
%
%   \item[period]
%   题注文字末尾追加句点。
%
%   \item[\UserDefined]
%   也可使用 |\Declare|\-|Caption|\-|Text|\-|Format| 自定义文字格式。
%   \SeeUserDefined
% \end{Options}
%
% \iffalse
% --------------------------------------------------------------------------- %
% \fi
%
% \PageBreak
% \subsection{对齐方式}
% \label{justification}
%
% \DescribeMacro{justification=}
% 除题注格式外，还可以指定 \emph{题注对齐方式}，由选项
% \begin{quote}
%   |justification=|\meta{对齐方式名称}\quad。
% \end{quote}
% 指定。
%
% 可在以下值中选择：
%
% \begin{Options}{\OptionLabel}
%   \item[justified]
%   将题注排版为普通两端对齐段落。
%
%   \item[centering]
%   题注每一行都居中排列。
%
% \iffalse
%  \item[Centering]
% \fi
%
%   \item[centerlast]
%   题注各段落的最后一行居中排列。
%
%   \item[centerfirst]
%   仅题注第一行居中排列。
%
%   \item[raggedright]
%   题注每行均靠左边界排列。
%
% \iffalse
%   \item[RaggedRight]
%   每行均靠左边界排列，使用 \pkg{ragged2e} 宏包的 |\RaggedRight| 命令。
% \else
%   \item[RaggedRight]
%   题注每行同样靠左边界排列，但此时使用 \pkg{ragged2e} 宏包的
%   |\RaggedRight| 命令实现。与 |raggedright| 的主要区别在于：题注内部
%   仍可使用 \TeX\ 的断词算法。\footnote{%
%     是否需要 \pkg{ragged2e} 宏包在运行时检测，因此首次使用本选项时
%     可能需要再编译一次 \LaTeX。}
% \fi
%
%   \item[raggedleft]
%   题注每行均靠右边界排列。
%
% \iffalse
%  \item[RaggedLeft]
%  每行均靠右边界排列，使用 \pkg{ragged2e} 宏包的 |\RaggedLeft| 命令。
% \fi
%
%   \item[\UserDefined]
%   也可使用 |\Declare|\-|Caption|\-|Jus|\-|ti|\-|fi|\-|ca|\-|tion| 自定义对齐方式。
%   \SeeUserDefined
% \end{Options}
%
% 下面是三个示例：
% \begin{Example}
%   \begin{quote}
%     |format=hang,justification=raggedright|
%   \end{quote}
%   \captionsetup{skip=0pt}
%   \example{format=hang,justification=raggedright}{\exampletext}
% \end{Example}
%
% \begin{Example}
%   \begin{quote}
%     |format=hang,justification=centerlast|
%   \end{quote}
%   \captionsetup{skip=0pt}
%   \example{format=hang,justification=centerlast}{\exampletext}
% \end{Example}
%
% \begin{Example}
%   \begin{quote}
%     |format=plain,labelsep=newline,justification=centering|
%   \end{quote}
%   \captionsetup{skip=0pt}
%   \example*{format=plain,labelsep=newline,justification=centering}{\exampletext}
% \end{Example}
%
% \smallskip
%
% \PageBreak
% \DescribeMacro{singlelinecheck=}
% 标准 \LaTeX\ 文档类（\cls{article}、\cls{report} 与 \cls{book}）会在
% 题注恰好能排成一行时自动将其居中：
%
% \example*{}{一则简短的题注。}
%
% \INFO
% \Thispackage\ 沿用了这一行为，因此在上述情况下通常会忽略你通过
% |justification=| 与 |indention=| 设置的对齐方式与缩进。但可以使用选项
% \begin{quote}
%   |singlelinecheck=|\meta{布尔值}\quad
% \end{quote}
% 关闭这种针对短题注的特殊处理。\meta{布尔值} 取 |false|、|no|、|off|
% 或 |0| 时关闭额外居中：
% \begin{quote}
%   |singlelinecheck=false|
% \end{quote}
% 这样上面那则短题注就会变成如下样式：
%
% \begingroup
%   \captionsetup{type=figure}
%   \ContinuedFloat
%   \example*{singlelinecheck=false}{一则简短的题注。}
% \endgroup
%
% \meta{布尔值} 取 |true|、|yes|、|on| 或 |1| 时重新开启额外居中
% （默认开启）。
%
% \iffalse
% --------------------------------------------------------------------------- %
% \fi
%
% \subsection{字体}
% \label{fonts}
%
% \DescribeMacro{font=}
% \DescribeMacro{labelfont=}
% \DescribeMacro{textfont=}
% 共有三个字体选项，分别作用于题注的不同部分：作用于整个题注的
% |font|、仅作用于题注标签与分隔符的 |label|\-|font|，以及仅作用于
% 题注文字的 |text|\-|font|。可使用以下选项设置：
% \begin{quote}\begin{tabular}{@{}r@{}ll}
%   |font=|      & \marg{字体选项} & ,\\
%   |labelfont=| & \marg{字体选项} & ,以及\\
%   |textfont=|  & \marg{字体选项} & ,\\
% \end{tabular}\end{quote}
% 其中 \meta{字体选项} 是以逗号分隔的字体选项列表。
%
% 可用的字体选项如下：
%
% \begin{Options}{stretch=\meta{系数}}
%   \item[scriptsize]   {\scriptsize 极小字号}
%   \item[footnotesize] {\footnotesize 通常用于脚注的字号}
%   \item[small]        {\small 小字号}
%   \item[normalsize]   {\normalsize 常规字号}
%   \item[large]        {\large 大字号}
%   \item[Large]        {\Large 更大字号}
% \end{Options}
% \vspace{0pt}\pagebreak[3]
% \begin{Options}{stretch=\meta{系数}}
%   \item[normalfont]   {\normalfont 常规字形、字重与字体族}
%
%   \item[up]           {\upshape 直立字形}
%   \item[it]           {\itshape 斜体字形}
%   \item[sl]           {\slshape 倾斜字形}
%   \item[sc]           {\scshape 小型大写字形}
%
%   \item[md]           {\mdseries 常规字重}
%   \item[bf]           {\bfseries 粗体字重}
%
%   \item[rm]           {\rmfamily 罗马字体族}
%   \item[sf]           {\sffamily 无衬线字体族}
%   \item[tt]           {\ttfamily 打字机字体族}
% \end{Options}
% \vspace{0pt}\pagebreak[3]
% \begin{Options}{stretch=\meta{系数}}
%   \item[singlespacing] 单倍行距 \See{\longref{setspace}}
%   \item[onehalfspacing] 一点五倍行距 \See{\longref{setspace}}
%   \item[doublespacing]  双倍行距 \See{\longref{setspace}}
%   \item[stretch=\meta{系数}] |\setstretch|\marg{系数} \See{\longref{setspace}}
% \end{Options}
% \vspace{0pt}\pagebreak[3]
% \begin{Options}{stretch=\meta{系数}}
%   \item[normalcolor]          |\normalcolor|
%   \item[color=\meta{颜色}]    |\color|\marg{颜色}
%        {\small（若已加载 \pkg{color} 或 \pkg{xcolor} 宏包，
%         示例见 \longref{declare}）}
% \end{Options}
% \vspace{0pt}\pagebreak[3]
% \begin{Options}{stretch=\meta{amount}}
%   \item[normal] |normal|\-|color|、|normal|\-|font|、|normal|\-|size|
%                 与 |single|\-|spacing| 等选项的组合
%
%   \item[\UserDefined]
%   也可使用 |\Declare|\-|Caption|\-|Font| 自定义字体选项。
%   \SeeUserDefined
% \end{Options}
%
% 若只使用其中一个选项，可以省略花括号；例如：
%   |font={small}|
% 与
%   |font=small|
% 结果相同。
%
% 下面是三个示例：
% \begin{Example}
%   \begin{quote}
%     |font=it,labelfont=bf|
%   \end{quote}
%   \captionsetup{skip=0pt}
%   \example{font=it,labelfont=bf}{\exampletext}
% \end{Example}
%
% \begin{Example}
%   \begin{quote}
%     |labelfont=sc,textfont=it|
%   \end{quote}
%   \captionsetup{skip=0pt}
%   \example{labelfont=sc,textfont=it}{\exampletext}
% \end{Example}
%
% \begin{Example}
%   \begin{quote}
%     |font={small,stretch=0.80}|
%   \end{quote}
%   \captionsetup{skip=0pt}
%   \example{font={small,stretch=0.80}}{\exampletext}
% \end{Example}
%
% \medskip
%
% \pagebreak[3]
% \DescribeMacro{font+=}
% \DescribeMacro{labelfont+=}
% \DescribeMacro{textfont+=}
% \NEWfeature*{v3.1}
% 也可以在现有字体选项的基础上追加字体选项，例如：
% \begin{quote}
%   |\captionsetup{font=small}|\\
%   |\captionsetup{font+=it}|
% \end{quote}
% 等价于
% \begin{quote}
%   |\captionsetup{font={small,it}}|
% \end{quote}
%
% \iffalse
% --------------------------------------------------------------------------- %
% \fi
%
% \subsection{边距与其他段落选项}
% \label{margins}
%
% \DescribeMacro{margin=}
% \DescribeMacro{width=}
% 对所有题注，可以指定 \emph{要么} 额外边距、\emph{要么} 固定宽度：
% ~\footnote{此处仅支持固定宽度；若希望把题注宽度限制为与插图或表格
% 同宽，请参阅 \pkg{floatrow}\cite{floatrow} 或
% \pkg{threeparttable}\cite{threeparttable} 宏包。}
% \begin{quote}\begin{tabular}{@{}r@{}ll}
%   |margin=| & \meta{长度} & \emph{——或者——}\\
%   |margin=| & |{|\meta{左长度}|,|\meta{右长度}|}| & \emph{——或者——}\\
%   |width=|  & \meta{长度} & \\
% \end{tabular}\end{quote}
% \NEWfeature{v3.1}
% 若边距只给出一个 \meta{长度}，则左右两边采用相同边距，例如
% |margin=|\x|10pt| 等价于 |margin=|\x|{10pt,10pt}|。
% 在双面排版文档中，偶数页上左右边距会互换。
% \DescribeMacro{oneside}
% \DescribeMacro{twoside}
% 若不希望互换，可额外指定选项 |oneside|，例如
% |\caption|\-|setup{margin=|\x|{0pt,10pt},|\x|oneside}|。\par
% 若指定的是宽度，则左右两边距取相同数值。
%
% 下面用三个示例加以说明：
% \begin{Example}
%   \begin{quote}
%     |margin=10pt|
%   \end{quote}
%   \captionsetup{skip=0pt}
%   \example{margin=10pt}{\exampletext}
% \end{Example}
%
% \begin{Example}
%   \begin{quote}
%     |margin={1cm,0cm}|
%   \end{quote}
%   \captionsetup{skip=0pt}
%   \example{margin={1cm,0cm}}{\exampletext}
% \end{Example}
%
% \begin{Example}
%   \begin{quote}
%     |width=.75\textwidth|
%   \end{quote}
%   \captionsetup{skip=0pt}
%   \example{width=.75\textwidth}{\exampletext}
% \end{Example}
%
% \begin{Note}
% 当题注置于内容旁侧时（例如 \pkg{sidecap} 宏包\cite{sidecap} 提供的
% |SC|\-|figure| 环境），或插图为文内环绕式插图时（例如 \pkg{wrapfig}
% 宏包\cite{wrapfig} 提供的 |wrap|\-|figure| 环境），环境一开始就会把
% 边距设置自动重置为 |0pt|。若确实需要为这些环境设置额外边距，可以在
% 环境内部设置，或针对该特定环境指定边距，例如
% |\caption|\-|setup[SC|\-|figure]{margin=|\x|10pt}|。
% \end{Note}
%
% \begin{Expert}
% \DescribeMacro{margin*=}
% \NEWfeature*{v3.1}
% |margin=| 选项还有一个带星号的变体 |margin*=|：仅当未设置宽度时，
% 它才修改边距。
% \end{Expert}
%
% \begin{Expert}
% \DescribeMacro{minmargin=}
% \DescribeMacro{maxmargin=}
% \NEWfeature*{v3.1}
% 还可以设置边距的最小或最大值。这在较小的环境（例如 |minipage|）中
% 限制边距时很有用。例如 \SmF 文档类会把边距限制为
% |maxmargin=|\x|0.1\linewidth|。（参见 \longref{SMF}）
% \end{Expert}
%
% \medskip
%
% \pagebreak[3]
% \DescribeMacro{parskip=}
% 当题注包含多个段落时，本选项用于设置段落之间插入的额外垂直间距：
% \begin{quote}
%   |parskip=|\meta{长度}
% \end{quote}
% 示例如下：
% \begin{Example}
%   \begin{quote}
%     |margin=10pt,parskip=5pt|
%   \end{quote}
%   \captionsetup{skip=0pt}
%   \example{margin=10pt,parskip=5pt}{%
%     题注第一段。这里写一些文字，只是为了展示这些选项如何影响题注的
%     排版外观。\par
%     题注第二段。这里同样写一些文字，以展示这些选项如何影响题注的
%     排版外观。}
% \end{Example}
%
% \pagebreak[3]
% \DescribeMacro{hangindent=}
% 选项
% \begin{quote}
%   |hangindent=|\meta{长度}
% \end{quote}
% 用于设置每个段落从第二行起的悬挂缩进。若题注只有一个段落，该选项
% 与已知的 |indention=|\meta{长度} 效果相同；但题注包含多个段落时，
% 就能看出差别：
%
% \begin{Example}
%   \begin{quote}
%     |format=hang,indention=-.5cm|
%   \end{quote}
%   \captionsetup{skip=0pt}
%   \example{format=hang,indention=-.5cm}{%
%     题注第一段。这里写一些文字，只是为了展示这些选项如何影响题注的
%     排版外观。\par
%     题注第二段。这里同样写一些文字，以展示这些选项如何影响题注的
%     排版外观。}
% \end{Example}
%
% \begin{Example}
%   \begin{quote}
%     |format=hang,hangindent=-.5cm|
%   \end{quote}
%   \captionsetup{skip=0pt}
%   \example{format=hang,hangindent=-.5cm}{%
%     题注第一段。这里写一些文字，只是为了展示这些选项如何影响题注的
%     排版外观。\par
%     题注第二段。这里同样写一些文字，以展示这些选项如何影响题注的
%     排版外观。}
% \end{Example}
%
% \begin{Note}
% 若题注包含多个段落，必须通过 |\caption| 或 |\caption|\-|of| 的可选参数
% 为图表目录指定另一则题注，否则会得到错误信息。
% \end{Note}
%
% \iffalse
% --------------------------------------------------------------------------- %
% \fi
%
% \PageBreak
% \subsection{样式}
% \label{style}
%
% \DescribeMacro{style=}
% 一组搭配得当的题注选项称为\emph{题注样式}。它大致可与用
% |\page|\-|style| 设置的页面样式相比：题注样式为整套题注布局提供
% 全部设置。
%
% 使用选项
% \begin{quote}
%   |style=|\meta{样式名称}\quad
% \end{quote}
% 即可切换到已定义的题注样式。\Thispackage\ 预定义了两种样式：
% |base| 与 |default|。
%
% \NEWfeature{v3.1}
% |base| 样式把前述各选项设置为反映基础 \LaTeX\ 文档类
% \cls{article}、\cls{report}、\cls{book} 题注外观的取值。也就是说，
% 指定选项
% \begin{quote}
%   |style=base|
% \end{quote}
% 与指定下列全部选项效果相同：
% \begin{quote}
%   |format=plain,labelformat=default,labelsep=colon,|\\
%   |justification=justified,font={},labelfont={},|\\
%   |textfont={},margin=0pt,indention=0pt|\\
%   |parindent=0pt,hangindent=0pt,singlelinecheck=true|
% \end{quote}
% \begin{Annotation*}
% 但当题注恰好排成一行时，仍会设置 |justification=centering,indention=0pt|。
% \end{Annotation*}
%
% 与之相对，|default| 样式遵循默认值，反映你实际使用的文档类所给出的
% 题注外观。该样式为默认选中，对应以下选项：
% \begin{quote}
%   |format=default,labelformat=default,labelsep=default,|\\
%   |justification=default,font=default,labelfont=default,|\\
%   |textfont=default,margin=0pt,indention=0pt|\\
%   |parindent=0pt,hangindent=0pt,singlelinecheck=true|
% \end{quote}
% \begin{Annotation*}
% 同样，当题注恰好排成一行时，仍会设置 |justification=centering,indention=0pt|。
% \end{Annotation*}
%
% 因此，若使用基础 \LaTeX\ 文档类 \cls{article}、\cls{report} 或
% \cls{book}，|base| 与 |default| 两种题注样式指向（几乎）相同的设置。
%
% \begin{Note}
% 也可使用 |\Declare|\-|Caption|\-|Style| 自定义题注样式。\LineBreak
% \SeeUserDefined
% \end{Note}
%
% \iffalse
% --------------------------------------------------------------------------- %
% \fi
%
% \subsection{间距}
% \label{skips}
%
% \DescribeMacro{skip=}
% 题注与插图或表格内容之间的垂直间距由选项
% \begin{quote}
%   |skip=|\meta{长度}\quad
% \end{quote}
% 控制。标准 \LaTeX\ 文档类 \cls{article}、\cls{report} 与 \cls{book}
% 将其预设为 |skip=|\x|10pt|，但其他文档类可能采用不同数值。
%
% \bigskip
%
% \pagebreak[3]
% \DescribeMacro{position=}
% \LaTeX{} 提供的 |\caption| 命令存在一处设计缺陷：它无法判断自己位于
% 插图或表格的开头还是末尾，因此也不知道该把题注与内容之间的间距放在
% 哪一侧。标准实现在浮动环境中总是把间距放在题注上方（而在
% |long|\-|table| 中又不一致地放在题注下方）；本宏包提供的实现则更为
% 灵活。
% 给出选项
% \begin{quote}
%   |position=top|\quad 或\quad |position=above|
% \end{quote}
% 时，假定题注位于环境的 \emph{顶部}，因此由 |skip=|\x\meta{长度}
% 设置的间距放在题注下方。
% （请注意，|position=|\x|top| 并\emph{不}表示题注实际就排在插图或表格
% 顶部；题注通常仍排在你放置 |\caption| 命令的位置。）
% 而使用
% \begin{quote}
%   |position=bottom|\quad 或\quad |position=below|
% \end{quote}
% 时，假定题注位于环境的 \emph{底部}，间距因此放在题注上方。
% 最后，使用
% \begin{quote}
%   |position=auto|\quad {\small（默认设置）}
% \end{quote}
% 时，\thispackage\ 会尽力自行判断题注的实际位置。请注意，虽然多数
% 情况下判断正确，但极少数情形下仍可能得出错误结果。
%
% \medskip
%
% \pagebreak[3]
% \DescribeMacro{figureposition=}
% \DescribeMacro{tableposition=}
% |position| 选项与 |\caption|\-|setup| 命令的可选参数搭配时尤其有用。
% \See{另见 \longref{captionsetup}}\par
% 例如
% \begin{quote}
%   |\captionsetup[table]{position=above}|
% \end{quote}
% 会把表格内所有题注都视为 \emph{位于表格上方} 的题注（就其周围间距
% 而言）。由于这是非常常见的设置，\thispackage\ 还提供了简写选项
% |figure|\-|position=|\x\meta{位置} 与 |table|\-|position=|\x\meta{位置}，
% 例如
% \begin{quote}
%   |\usepackage[|\ldots|,tableposition=top]{caption}|
% \end{quote}
% 等价于
% \begin{quote}
%   |\usepackage[|\ldots|]{caption}|\\
%   |\captionsetup[table]{position=top}|
% \end{quote}
%
% \medskip
%
% \INFO
% 请注意，|skip=|、|position=|、|figure|\-|position=| 与
% |table|\-|position=| 选项并非总是生效。由于 |figure| 与 |table| 环境
% 由文档类提供，文档类可能采用自己的间距，并自行决定题注按“顶部”还是
% “底部”题注排版。例如 \KOMAScript\ 文档类支持 |skip=| 设置，但始终把
% |figure| 题注按“底部”题注排版，而 |table| 题注取决于全局选项
% |table|\-|captions|\-|above| 或 |table|\-|captions|\-|below|。
% \See{\longref{KOMA}}
%
% 此外，有些宏包自行控制题注上下间距的行为，例如 \pkg{float}、
% \pkg{floatrow} 与 \pkg{supertabular} 宏包。
%
% \bigskip
%
% \begin{Expert*}
% 在内部，题注与内容之间的间距由 |\above|\-|caption|\-|skip| 表示
% （在 \LaTeX{} 的实现中它始终位于题注上方）。此外还有第二个值
% |\below|\-|caption|\-|skip|（默认为 |0pt|），在 \LaTeX{} 实现中位于
% 题注下方；而在本宏包中它位于与 |\above|\-|caption|\-|skip| 相反的一侧。
% 因此严格来说，当设置 |position=|\x|top| 时，本宏包会交换这两个间距的
% 含义。请注意，另有若干宏包采用相同手法（例如 \pkg{ftcap}、\pkg{nonfloat}
% 与 \pkg{topcap} 宏包）；若同时使用其中任一宏包，则不支持 \pkg{caption}
% 的 |position=| 选项。
% \end{Expert*}
%
% \iffalse
% --------------------------------------------------------------------------- %
% \fi
%
% \subsection{目录条目}
% \label{lists}
%
% \DescribeMacro{list=}
% \NEWfeature*{v3.1}
% |\caption| 命令通常会在插图目录或表格目录中添加一条目。既可以通过给
% |\caption| 传入空可选参数来逐条取消{\small（见 \longref{caption}）}，
% 也可以程序化地设置
% \begin{quote}
%   |list=no|\quad（或用 |no| 以外的其他布尔值代替）\quad。~\footnote{%
%   请注意，\pkg{subfig} 宏包\cite{subfig} 不支持本选项；它改用计数器
%   \texttt{lofdepth} 与 \texttt{lotdepth} 实现相应功能。}
% \end{quote}
%
% \pagebreak[3]
% \DescribeMacro{listformat=}
% \NEWfeature*{v3.1}
% 使用选项
% \nopagebreak[3]
% \begin{quote}
%   |listformat=|\meta{目录格式名称}
% \end{quote}
% \nopagebreak[3]
% 可以指定插图目录或表格目录中编号的排版方式。内置五种标准目录格式：
%
% \begin{Options}{\OptionLabel}
%   \item[empty]
%   不排版编号。
%
%   \item[simple]
%   编号连同标签前缀一起排版。
%
%   \item[parens]
%   编号放在圆括号内，连同标签前缀一起排版。
%
%   \item[subsimple]
%   与 |simple| 相同，但不带标签前缀。（默认）
%
%   \item[subparens]
%   与 |parens| 相同，但不带标签前缀。
%
%   \item[\UserDefined]
%   也可使用 |\Declare|\-|Caption|\-|List|\-|Format| 自定义目录格式。
%   \SeeUserDefined
% \end{Options}
%
% 编号（即 |\the|\-|figure| 或 |\the|\-|table|）之前的前缀
% （即 |\p@figure| 或 |\p@table|）在目录（如插图/表格目录）和引用中
% 通常为空，因此目录格式 |simple| 与 |subsimple| 结果相同，|parens|
% 与 |subparens| 亦然。但当子插图、子表格列入插图目录或表格目录时，
% 或标签前缀被另作他用时，情况可能不同。\footnote{%
%   子插图与子表格可使用 \pkg{subcaption} 或 \pkg{subfig} 宏包排版。}
%
% \iffalse
% --------------------------------------------------------------------------- %
% \fi
%
% \subsection{名称}
% \label{names}
%
% \DescribeMacro{name=}
% \NEWfeature*{v3.1}
% 选项
% \begin{quote}
%   |name=|\meta{名称}
% \end{quote}
% 用于更改 \emph{当前} 环境的名称。它与 |\caption|\-|setup| 的可选参数
% 搭配时会很有用，例如
% \begin{quote}
%   |\captionsetup[wrapfigure]{name=Fig.}|
% \end{quote}
% 把所有 |wrap|\-|figure| 的名称改为“Fig.”（而其他插图题注仍以
% “Figure” 作为名称，即英文文档中的默认插图名）。
%
% \iffalse
% --------------------------------------------------------------------------- %
% \fi
%
% \subsection{类型}
% \label{types}
%
% \DescribeMacro{type=}
% |\caption| 命令可以为不同类型（例如 |figure| 与 |table|）排版题注。
% 若在这些环境之外使用 |\caption| 命令，会得到错误信息，因为它不知道
% 此处要排哪种题注。但在这种情形下，可以在使用 |\caption| 命令之前
% 手动设置题注类型：
% \begin{quote}
%   |type=|\meta{浮动体类型}
% \end{quote}
% （其他命令同样适用，例如 |\Continued|\-|Float|、\pkg{subcaption}
% 宏包提供的 |\sub|\-|caption|\-|box|，或 \pkg{subfig}
% 宏包\cite{subfig} 提供的 |\sub|\-|float|。）例如在 |minipage|
% 这类非浮动环境中：
% \begin{quote}
%   |\noindent\begin{minipage}{\textwidth}|\\
%   |  \captionsetup{type=figure}|\\
%   |  \subfloat{|\ldots|}|\\
%   |  |\ldots\\
%   |  \caption{|\ldots|}|\\
%   |\end{minipage}|
% \end{quote}
%
% \begin{Expert}
% 本选项还有带星号的变体 |type*=|\meta{浮动体类型}：加载
% \pkg{hyperref} 宏包\cite{hyperref} 时，其行为与 |type=|\meta{浮动体类型}
% 不同。|type=| 会设置超链接锚点（当 |hypcap=|\x|true| 时），而
% |type*=| 不会。（另见 \longref{hyperref}\,）
% \end{Expert}
% \begin{Expert}
% \emph{\DefaultNoteText:} 请不要像某些文档建议的那样自行重定义内部宏
% |\@captype|；应始终改用 |\caption|\-|setup{type=|\x\ldots|}|。
% \end{Expert}
%
% \medskip
%
% 自定义浮动体类型可通过以下方式定义：\pkg{newfloat} 宏包提供的
% |\Declare|\-|Floating|\-|Environment|、\pkg{float}
% 宏包\cite{float} 提供的 |\new|\-|float|，或 \pkg{floatrow}
% 宏包\cite{float} 提供的 |\Declare|\-|New|\-|Float|\-|Type|。
%
% \medskip
%
% \INFO % \NEWdescription{v3.1}
% 请注意，|type=| 选项只应在盒子或环境 \emph{内部} 使用（例如
% |\par|\-|box| 或 |mini|\-|page|），最好是内容与题注之间不可能发生
% 分页的盒子或环境。此外，若不使用盒子或环境，可能出现一些视觉副作用
% （例如插图与表格的题注设置互相混淆）；因此强行这样做时会发出警告。
% \footnote{%
% 仅当底层 \TeX\ 引擎为 \eTeX\ 时才会给出此警告。}
%
% \iffalse
% --------------------------------------------------------------------------- %
% \fi
%
% \clearpage
% \section{命令}
%
% \subsection{题注排版}
% \label{caption}
% \label{captionlistentry}
%
% \DescribeMacro{\caption}
% 命令
% \begin{quote}
%   |\caption|\oarg{目录条目}\marg{标题文字}
% \end{quote}
% 在 |figure| 或 |table| 等浮动环境中排版题注。这一点你已经知道了，
% 而 \thispackage\ 还提供了一项扩展：若将参数 \meta{目录条目} 留空，
% 则不向插图目录或表格目录添加条目。例如：
% \begin{quote}
%   |\caption[]{一则不含目录条目的插图。}|
% \end{quote}
%
% \begin{Expert}
% 请记住：当未给出 \meta{目录条目} 时，\meta{标题文字} 是所谓的
% \emph{移动}参数；若给出了 \meta{目录条目}，则该参数变为移动参数。
% 所谓\emph{移动}参数，是指会写入目录文件、从而出现在“插图目录”或
% “表格目录”中的参数。\emph{移动}参数中不允许包含\emph{脆弱}命令，
% 所有内容都必须是\emph{稳固}的，否则参数可能被\emph{破坏}，在下一次
% \LaTeX\ 编译时产生奇怪的错误。部分\emph{脆弱}命令可在前面加
% |\protect| 加以\emph{保护}；自定义定义则可改用
% |\Declare|\-|Robust|\-|Command| 代替 |\new|\-|command|，使其
% 变得\emph{稳固}。
%
% 例如：|\caption{${}^{137}_{\phantom{1}55}$Cs}| 会报错，因为
% |\phantom| 是\emph{脆弱}命令。此时要么使用可选参数 \meta{目录条目}
% （例如 |\caption[${}^|\x|{137}_|\x|{55}$|\x|Cs]|\x|{${}^|\x|{137}_|\x|{\phantom{1}55}$|\x|Cs}|），
% 要么加上 |\protect|，避免 |\phantom| 命令被\emph{破坏}：
% |\caption{${}^|\x|{137}_|\x|{|\textcolor{blue}{\cs{protect}}|\phantom{1}55}$|\x|Cs}|。
%
% 但有时这样仍不够。其原因在于所谓的单行检查：它会把 \meta{标题文字}
% 放入水平盒子以测定题注宽度，这同样可能引发错误。例如：
% |\caption{一则| |示意图。| |\[V_{C}| |\sim| |\left| |\{| |\begin{array}{cc}|
% |E_{g}| |&| |\textrm{p-n}| |\\| |e\phi_{B}| |&| |\textrm{M-S}| |\end{array}|
% |\right.| |\]}|.
% 这里仅用 |\caption[一则| |示意图]{|\ldots|}| 仍不足以避免错误（仍会出现
% “\texttt{Missing \$ inserted.}”）。因此还必须在 |\caption| 命令正前方
% 加上 |\caption|\-|setup{single|\-|line|\-|check=|\x|off}|。
%
% 有关\emph{移动}参数以及\emph{脆弱}与\emph{稳固}命令的更多信息，请参阅
% 你的 \LaTeX\ 手册，或访问
% \url{http://www-h.eng.cam.ac.uk/help/tpl/textprocessing/teTeX/latex/latex2e-html/fragile.html}。
% \end{Expert}
%
% \medskip
%
% \pagebreak[3]
% \DescribeMacro{\caption*}
% \pkg{longtable} 宏包定义了命令 |\caption*|，用于排版不带标签、也不进入
% 表格目录的题注。示例：
% \begin{quote}
%   |\begin{longtable}{|\ldots|}|\\
%   |  \caption*{一个表格}\\|\\
%   |  |\ldots\\
%   |\end{longtable}|
% \end{quote}
% 外观如下：
% \DeleteShortVerb{\|}
% \vskip\intextsep
% \begin{minipage}{\linewidth}
%   \captionsetup{type=table,position=t}
%   \caption*{一个表格}
%   \centering
%   \begin{tabular}{r|rr}
%       & x & y \\\hline
%     a & 1 & 2 \\
%     b & 3 & 4 \\
%   \end{tabular}
% \end{minipage}
% \vskip\intextsep
% \MakeShortVerb{\|}
% \Thispackage\ 同样提供了这一功能，因此现在在 |figure|、|table| 等
% 各类浮动环境中都可使用该命令，例如：
% \begin{quote}
%   |\begin{table}|\\
%   |  \caption*{一个表格}|\\
%   |  |\ldots\\
%   |\end{table}|
% \end{quote}
%
% \pagebreak[3]
% \DescribeMacro{\captionof}
% \DescribeMacro{\captionof*}
% 有时需要在浮动环境\emph{之外}排版题注，例如把插图放入非浮动的
% |mini|\-|page| 中。为此 \thispackage\ 提供命令
% \begin{quote}
%   |\captionof|\marg{浮动体类型}\oarg{目录条目}\marg{标题文字}\quad。
% \end{quote}
% 请注意，第一个参数 \meta{浮动体类型} 在这里是必填的，因为
% |\captionof| 需要知道题注标签中使用什么名称（例如“插图”还是
% “表格”），以及把目录条目放入哪一个目录。示例：
% \begin{quote}
%   |\captionof{table}{一个表格}|
% \end{quote}
% 排版出的题注如下：
% \begin{Example}
%   \captionsetup{position=top}
%   \captionof{table}{一个表格}
% \end{Example}
%
% 带星号的变体 |\caption|\-|of*| 行为与 |\caption*| 命令相同：
% 排版不带标签、也不进入插图目录或表格目录的题注。
%
% \INFO % \NEWdescription{v3.1}
% 由于 |\caption|\-|of| 内部使用 |type| 选项，因而同样适用于
% |type| 选项的限制：|\caption|\-|of| 与 |\caption|\-|of*| 都只应在
% 盒子或环境\emph{内部}使用。
% \See{\longref{types}}
%
% \bigskip
%
% \pagebreak[3]
% \DescribeMacro{\captionlistentry}
% \NEWfeature*{v3.1}
% 在某些情形下，单独添加一条插图目录（或表格目录）条目会很有用。
% 这可以通过
% \begin{quote}
%   |\captionlistentry|\oarg{浮动体类型}\marg{目录条目}\quad
% \end{quote}
% 实现。
%
% 举例来说：在 |long|\-|table| 中，题注置于内容上方、且目录中只保留
% 一条指向表格首页的条目，是很容易做到的：
% \begin{quote}
%   |\begin{longtable}{|\ldots|}|\\
%   |  \caption{|\ldots|}\\|\\
%   |\endfirsthead|\\
%   |  \caption[]{|\ldots|}\\|\\
%   |\endhead|\\
%   |  |\ldots
% \end{quote}
% 但由于 \pkg{longtable} 宏包没有提供 |\end|\-|first|\-|foot| 命令，
% 要让题注置于表格内容\emph{下方}、同时目录中只保留一条指向表格
% 首页的条目就不那么容易了。这时就可以使用 |\caption|\-|list|\-|entry|
% 命令：
% \begin{quote}
%   |\begin{longtable}{|\ldots|}|\\
%   |  \caption[]{|\ldots|}\\|\\
%   |\endfoot|\\
%   |  \captionlistentry{|\ldots|}|\\
%   |  |\ldots
% \end{quote}
%
% \begin{Annotation*}
% 另一个示例见 \longref{examples}。
% \end{Annotation*}
%
% \begin{Expert}
% 还有带星号的变体 |\caption|\-|list|\-|entry*|，它不会使
% \meta{浮动体类型} 计数器加一。
% （请注意，在 |long|\-|table| 环境内部，|\caption|\-|list|\-|entry|
% 从不会使 |table| 计数器加一。另见 \longref{longtable}。）
% \end{Expert}
%
% \begin{Expert}
% 请注意，\meta{目录条目} 是\emph{移动}参数，其中所有内容都必须是
% \emph{稳固}的。（另见 |\caption| 的说明。）
% \end{Expert}
%
% \subsection{设置选项}
% \label{captionsetup}
%
% \DescribeMacro{\captionsetup}
% 我们已经认识了 |\caption|\-|setup| 命令（见 \longref{usage}），
% 这里进一步说明其可选参数 \meta{浮动体类型}。
%
% 回忆一下，该命令的语法为
% \begin{quote}
%   |\captionsetup|\oarg{浮动体类型}\marg{选项}\quad。
% \end{quote}
% 若指定了 \meta{浮动体类型}，所有 \meta{选项} 此时不会立即生效；
% 它们只是被标记为供以后使用——当在该特定类型
% \meta{浮动体类型} 的浮动环境中排版题注时才起作用。例如
% \begin{quote}
%   |\captionsetup[figure]|\marg{选项}
% \end{quote}
% 会强制 |figure| 环境中的题注使用给定的 \meta{选项}。
%
% 下面的示例加以说明：
% \begin{quote}
%   |\captionsetup{font=small}|\\
%   |\captionsetup[figure]{labelfont=bf,textfont=it}|
% \end{quote}
% 得到如下题注：
% \begin{Example}
%   \captionsetup{font=small}
%   \captionsetup[figure]{labelfont=bf,textfont=it}
%   \begingroup
%     \captionof{figure}{一个插图}
%   \endgroup
%   \captionsetup{position=top}
%   \begingroup
%     \captionof{table}{一个表格}
%   \endgroup
% \end{Example}
% 可以看到，命令 |\caption|\-|setup[figure]{|\ldots|}| 只改变
% |figure| 题注标签的外观，不影响其他题注。
%
% \meta{浮动体类型} 通常只能取 |figure| 或 |table| 之一。但如后文所述，
% 有些 \LaTeX\ 宏包（例如 \pkg{floatrow}、\pkg{longtable}、\pkg{sidecap}）
% 以及本宏包还提供了其他带题注的环境，这两个命令同样可用于这些环境。
% \See{\longref{declare} 和 \longref{packages}}
%
% \begin{Expert}
% \NEWfeature{v3.1}
% |\caption|\-|setup| 还有带星号的变体：
% \begin{quote}
%   |\captionsetup*|\oarg{浮动体类型}\marg{选项}
% \end{quote}
% 不带星号的变体可能给出警告——例如 \meta{选项} 在整个文档中实际
% 未被使用（如没有 |table| 却写了 |\caption|\-|setup[table]{font=|\x|sf}|）；
% 带星号的变体则不会发出此类警告。
% \iffalse
% 当设置与具体文档无关的通用选项时，这一点会相当有用。
% \fi
% \end{Expert}
%
% \medskip
%
% \pagebreak[3]
% \DescribeMacro{\clearcaptionsetup}
% 若要清除为某个特定环境标记、供自动使用的这些参数，可使用命令
% \begin{quote}
%   |\clearcaptionsetup|\oarg{选项}\marg{浮动体类型}\quad。
% \end{quote}
%
% 例如 |\clearcaptionsetup{figure}| 会清除上例中为插图设置的全部附加处理：
% \begin{quote}
%   |\captionsetup{font=small}|\\
%   |\captionsetup[figure]{labelfont=bf,textfont=it}|\\
%   \ldots\\
%   |  \caption{一个插图}|\\
%   \ldots\\
%   |\clearcaptionsetup{figure}|\\
%   \ldots\\
%   |  \caption{一个插图}|\\
%   \ldots
% \end{quote}
% \begin{Example}
%   \captionsetup{font=small}
%   \captionsetup[figure]{labelfont=bf,textfont=it}
%   \begingroup
%     \captionof{figure}{一个插图}
%   \endgroup
%   \captionsetup{position=top}
%   \clearcaptionsetup{figure}
%   \begingroup
%     \captionof{figure}{一个插图}
%   \endgroup
% \end{Example}
%
% \pagebreak[3]
% \NEWfeature{v3.1}
% 若给出可选参数 \meta{选项}，则仅清除与该特定 \meta{选项} 相关的设置。
% \footnote{此处只能指定\emph{一个}选项，不能是选项列表。若要清除多个
% 选项，需要多次调用 \cs{clearcaptionsetup}。}
% 上例清除的是插图的全部选项（而不仅是 |labelfont=bf,|\x|textfont=it|）；
% 而下面这一条只清除 |labelfont=bf| 设置，插图的其他设置保持不变：
% \begin{quote}
%   |\captionsetup{font=small}|\\
%   |\captionsetup[figure]{labelfont=bf,textfont=it}|\\
%   \ldots\\
%   |  \caption{一个插图}|\\
%   \ldots\\
%   |\clearcaptionsetup[labelfont]{figure}|\\
%   \ldots\\
%   |  \caption{一个插图}|\\
%   \ldots
% \end{quote}
% \begin{Example}
%   \captionsetup{font=small}
%   \captionsetup[figure]{labelfont=bf,textfont=it}
%   \begingroup
%     \captionof{figure}{一个插图}
%   \endgroup
%   \captionsetup{position=top}
%   \clearcaptionsetup[labelfont]{figure}
%   \begingroup
%     \captionof{figure}{一个插图}
%   \endgroup
% \end{Example}
% \begin{Expert}
% 与 |\caption|\-|setup*| 类似，也有带星号的形式
% |\clear|\-|caption|\-|setup*|：当给定 \meta{选项} 并未为指定
% \meta{浮动体类型} 设置过时，它不发出警告。
% \end{Expert}
%
% \medskip
%
% \pagebreak[3]
% \DescribeMacro{\showcaptionsetup}
% 出于调试目的，本宏包提供命令
% \nopagebreak[3]
% \begin{quote}
%   |\showcaptionsetup|\marg{浮动体类型}
% \end{quote}
% \nopagebreak[3]
% 它会在日志文件中输出指定 \meta{浮动体类型} 所设置的选项。例如
% \begin{quote}
%   |\captionsetup[figure]{labelfont=bf,textfont=it}|\\
%   |\showcaptionsetup{figure}|
% \end{quote}
% 给出如下信息：
% \begin{quote}\small
%   |Caption Info: Option list on `figure'|\\
%   |Caption Data: {labelfont=bf,textfont=it} on input line 5.|
% \end{quote}
%
% \pagebreak[3]
% \subsection{续接浮动体}
% \label{ContinuedFloat}
%
% \DescribeMacro{\ContinuedFloat}
% 有时需要把插图或表格拆分到多个浮动体中，但不希望它们各自获得独立
% 的编号。命令
% \begin{quote}
%   |\ContinuedFloat|
% \end{quote}
% 正是为此而设；它应作为浮动环境中的第一条命令使用。它会阻止相应计数器
% （通常由 |\caption|）加一，因此内部含 |\Continued|\-|Float| 的插图或
% 表格，与前一个插图或表格使用相同的编号。示例：
% \begin{quote}
%   |\begin{table}|\\
%   |  \caption{一个表格}|\\
%   |  |\ldots\\
%   |\end{table}|\\
%   \ldots\\
%   |\begin{table}\ContinuedFloat|\\
%   |  \caption{一个表格（续）}|\\
%   |  |\ldots\\
%   |\end{table}|
% \end{quote}
% 结果如下：
% \begin{Example}
%   \captionsetup{type=table,position=b,skip=0pt}
%   \caption{一个表格}
%   \centerline{\ldots}
% \end{Example}
% \begin{Example}
%   \captionsetup{type=table,position=b,skip=0pt}
%   \ContinuedFloat
%   \caption{一个表格（续）}
%   \centerline{\ldots}
% \end{Example}
%
% \bigskip
%
% \NEWfeature{v3.1}
% 此外，|\Continued|\-|Float| 命令会执行与类型名 “|Continued|\-|Float|”
% 关联的选项。例如可借此为续接的插图或表格切换到不同的标签格式，如下所示：
% \begin{quote}
%   |\DeclareCaptionLabelFormat{continued}{Continued #1~#2}|\\
%   |\captionsetup[ContinuedFloat]{labelformat=continued}|\\
%   \ldots\\
%   |\begin{table}\ContinuedFloat|\\
%   |  \caption{一个表格}|\\
%   |  |\ldots\\
%   |\end{table}|
% \end{quote}
% \begin{Example}
%   \captionsetup{type=table,position=b,skip=0pt}
%   \captionsetup[ContinuedFloat]{labelformat=continued1}
%   \ContinuedFloat
%   \caption{一个表格}\label{continued1}
%   \centerline{\ldots}
% \end{Example}
% \smallskip
% \See{\longref{declare} 中对
%      \cs{Declare}\-\texttt{Caption}\-\texttt{Label}\-\texttt{Format} 的说明。}
%
% \medskip
%
% 另有一个名为 |Continued|\-|Float| 的 \LaTeX\ 计数器，可供自定义使用。
% 普通（浮动）环境中它被置为零；第一个续接浮动体中为一，第二个中为二，
% 依此类推。因此每条 |\Continued|\-|Float| 都会使该计数器加一，而不含
% |\Continued|\-|Float| 命令的浮动环境会把它重置为零。示例：
% \begin{quote}
%   |\DeclareCaptionLabelFormat{cont}{#1~#2\alph{ContinuedFloat}}|\\
%   |\captionsetup[ContinuedFloat]{labelformat=cont}|\\
%   \ldots\\
%   |\begin{table}\ContinuedFloat|\\
%   |  \caption{一个表格}|\\
%   |  |\ldots\\
%   |\end{table}|
% \end{quote}
% \begin{Example}
%   \captionsetup{type=table,position=b,skip=0pt}
%   \captionsetup[ContinuedFloat]{labelformat=continued2}
%   \ContinuedFloat
%   \caption{一个表格}\label{continued2}
%   \centerline{\ldots}
% \end{Example}
% \medskip
% 由于只更改了题注标签格式，引用该表时仍输出
% 表 3。若还希望在引用中使用 |Continued|\-|Float|
% 计数器，可改为重定义命令 |\the|\-|Continued|\-|Float|；在续接浮动体中，
% 该命令会自动附加到插图或表格计数器之后，默认为空。
% \begin{quote}
%   |\renewcommand\theContinuedFloat{\alph{ContinuedFloat}}|\\
%   \ldots\\
%   |\begin{table}\ContinuedFloat|\\
%   |  \caption{一个表格}|\\
%   |  |\ldots\\
%   |\end{table}|
% \end{quote}
% \begin{Example}
%   \captionsetup{type=table,position=b,skip=0pt}
%   \renewcommand\theContinuedFloat{\alph{ContinuedFloat}}
%   \ContinuedFloat
%   \caption{一个表格}\label{continued3}
%   \centerline{\ldots}
% \end{Example}
% \medskip
% 引用该表时将输出 表 3d。
%
% \pagebreak[3]
% \DescribeMacro{\ContinuedFloat*}
% 假定你希望这一系列的第一个插图或表格编号为“图 7a”，而不是
% “图 7”（第二个则为“图 7b”，而非“图 7a”）。
% 这同样可以做到：使用带星号的变体 |\Continued|\-|Float*|。它与
% |\Continued|\-|Float| 一样会执行与类型名 “|Continued|\-|Float|”
% 关联的选项，并使 \LaTeX\ 计数器 |Continued|\-|Float| 加一，\emph{但}
% 把它标记为系列中的第一个插图或表格：
% \begin{quote}
%   |\renewcommand\theContinuedFloat{\alph{ContinuedFloat}}|\\
%   \ldots\\
%   |\begin{figure}\ContinuedFloat*|\\
%   |  |\ldots\\
%   |  \caption{系列插图之一}|\\
%   |\end{figure}|\\
%   \ldots\\
%   |\begin{figure}\ContinuedFloat|\\
%   |  |\ldots\\
%   |  \caption{系列插图之二}|\\
%   |\end{figure}|\\
%   \ldots\\
%   |\begin{figure}\ContinuedFloat|\\
%   |  |\ldots\\
%   |  \caption{系列插图之三}|\\
%   |\end{figure}|
% \end{quote}
% \begin{Example}
%   \captionsetup{type=figure,position=b,skip=0pt}
%   \renewcommand\theContinuedFloat{\alph{ContinuedFloat}}
%   \begingroup
%     \ContinuedFloat*
%     \centerline{\ldots}
%     \caption{系列插图之一}
%   \endgroup
%   \begingroup
%     \ContinuedFloat
%     \centerline{\ldots}
%     \caption{系列插图之二}
%   \endgroup
%   \begingroup
%     \ContinuedFloat
%     \centerline{\ldots}
%     \caption{系列插图之三}
%   \endgroup
% \end{Example}
%
% \begin{Note}
%   遗憾的是，加载 \pkg{subfig} 宏包\cite{subfig} 时，
%   \cs{ContinuedFloat*} 不可用。
% \end{Note}
%
% \pagebreak[3]
% \subsubsection*{关于 longtable 宏包的说明}
% 若希望在 \pkg{longtable} 宏包\cite{longtable} 提供的 |long|\-|table|
% 跨页之后使用不同的题注标签，\emph{不能}借助 |\Continued|\-|Float|
% 实现；可以改为编写类似下面的代码：
% \begin{quote}
%   |\DeclareCaptionLabelFormat{continued}{Continued #1~#2}|\\
%   \ldots\\
%   |\begin{longtable}{|\ldots|}|\\
%   |  \caption{长表格}\\|\\
%   |\endfirsthead|\\
%   |  \captionsetup{labelformat=continued}|\\
%   |  \caption[]{长表格}\\|\\
%   |\endhead|\\
%   |  |\ldots\\
%   |\end{longtable}|
% \end{quote}
%
% \iffalse
% --------------------------------------------------------------------------- %
% \fi
%
% \clearpage
% \section{自定义增强}
% \label{declare}
%
% 本宏包提供了一组命令，允许用户定义自己的格式等内容。借助它们，可以把
% 分隔符、对齐方式、字体与样式等信息与一个名称关联起来，并集中保存。
% （这些命令必须出现在文档导言区，即 |\document|\-|class| 与
% |\begin{document}| 之间的部分。）
%
% \pagebreak[3]
% \DescribeMacro{\DeclareCaptionFormat}
% 可使用命令定义自己的题注格式：
% \begin{quote}
%   |\DeclareCaptionFormat|\marg{名称}\marg{使用 \#1、\#2、\#3 的代码}\mbox{\quad.}
% \end{quote}
% 使用时，系统将 \#1 替换为题注标签，\#2 替换为分隔符，\#3 替换为题注
% 文字。例如标准格式 |plain| 在 \thispackage\ 中预定义为
% \begin{quote}
%   |\DeclareCaptionFormat{plain}{#1#2#3\par}|\quad。
% \end{quote}
% \begin{Expert}
% 还有带星号的变体 |\Declare|\-|Caption|\-|Format*|，它使代码在 \TeX\
% 的垂直模式（而非水平模式）中排版，但不支持 |indention=| 选项。
% \end{Expert}
%
% \pagebreak[3]
% \DescribeMacro{\DeclareCaptionLabelFormat}
% 同样可以定义自己的题注标签格式：
% \begin{quote}
%   |\DeclareCaptionLabelFormat|\marg{名称}\marg{使用 \#1、\#2 的代码}
% \end{quote}
% 使用时，\#1 替换为名称（例如 “figure”），\#2 替换为编号（例如 “12”）。
% 示例：
% \begin{quote}
%   |\DeclareCaptionLabelFormat{bf-parens}{(\textbf{#2})}|\\
%   |\captionsetup{labelformat=bf-parens,labelsep=quad}|
% \end{quote}
% \example*{labelformat=bf-parens,labelsep=quad}{\exampletext}
%
% \pagebreak[3]
% \DescribeMacro{\bothIfFirst}
% \DescribeMacro{\bothIfSecond}
% 若自定义题注标签格式并使用 \pkg{subcaption} 或
% \pkg{subfig}\cite{subfig} 宏包，需要留意题注标签名称为空的情形。
% 为此本宏包提供命令
% \begin{quote}
%   |\bothIfFirst|\marg{第一参数}\marg{第二参数}\quad 以及\\
%   |\bothIfSecond|\marg{第一参数}\marg{第二参数}
% \end{quote}
% |\bothIfFirst| 检测第一参数是否存在（即非空），|\bothIfSecond| 检测
% 第二参数是否存在。存在时两个参数都排版，否则一个都不排版。
%
% \smallskip
%
% \pagebreak[3]
% 例如标准标签格式 |simple| \emph{并非} 定义为
% \begin{quote}
%   |\DeclareCaptionLabelFormat{simple}{#1~#2}|\mbox{\quad,}
% \end{quote}
% 因为当 \#1 为空时这会产生多余空格。|simple| 实际定义为
% \begin{quote}
%   |\DeclareCaptionLabelFormat{simple}%|\\
%   |                  {\bothIfFirst{#1}{~}#2}|\mbox{\quad,}
% \end{quote}
% 使得仅在标签名称存在时才出现该空格。
%
% \smallskip
%
% \pagebreak[3]
% \DescribeMacro{\DeclareCaptionTextFormat}
% \NEWfeature*{v3.1}
% 同样可以定义自己的题注文字格式：
% \begin{quote}
%   |\DeclareCaptionTextFormat|\marg{名称}\marg{使用 \#1 的代码}
% \end{quote}
% 使用时 \#1 被替换为题注文字。
%
% \bigskip
%
% \pagebreak[3]
% \DescribeMacro{\DeclareCaptionLabelSeparator}
% 可使用命令定义自己的题注标签分隔符：
% \begin{quote}
%   |\DeclareCaptionLabelSeparator|\marg{名称}\marg{代码}\mbox{\quad.}
% \end{quote}
% 再举一个取自 \thispackage\ 自身的简单示例：
% \begin{quote}
%   |\DeclareCaptionLabelSeparator{colon}{: }|
% \end{quote}
% \begin{Expert}
% 还有带星号的变体 |\Declare|\-|Caption|\-|Label|\-|Separator*|，
% 它使代码排版时不应用 |label|\-|font=| 设置。例如分隔符 |quad|、
% |new|\-|line|、|en|\-|dash| 就是以这种方式定义的。
% \end{Expert}
%
% \bigskip
%
% \pagebreak[3]
% \DescribeMacro{\DeclareCaptionJustification}
% 可使用命令定义自己的题注对齐方式：
% \begin{quote}
%   |\DeclareCaptionJustification|\marg{名称}\marg{代码}\mbox{\quad.}
% \end{quote}
% \meta{代码} 会直接排版在题注之前。例如对齐方式 |raggedright| 定义为
% \begin{quote}
%   |\DeclareCaptionJustification{raggedright}{\raggedright}|\mbox{\quad,}
% \end{quote}
% 可将题注各行靠左边界排列。
%
% \bigskip
%
% \pagebreak[3]
% \DescribeMacro{\DeclareCaptionFont}
% 可使用命令定义自己的题注字体：
% \begin{quote}
%   |\DeclareCaptionFont|\marg{名称}\marg{代码}\quad。
% \end{quote}
% 例如本宏包把选项 |small| 与 |bf| 定义为
% \begin{quote}
%   |\DeclareCaptionFont{small}{\small}|\quad 以及\\
%   |\DeclareCaptionFont{bf}{\bfseries}|\quad。
% \end{quote}
% \iffalse
% 行距可借助 \pkg{setspace} 宏包自定义，例如：%\NEWdescription{v3.0h}
% \begin{quote}
%   |\usepackage{setspace}|\\
%   |\DeclareCaptionFont{singlespacing}{\setstretch{1}}|~\footnote{%
%   \emph{\DefaultNoteText:} 此处不能使用 \cs{singlespacing}，因为它含
%   \cs{vskip} 命令。}\\
%   |\DeclareCaptionFont{onehalfspacing}{\onehalfspacing}|\\
%   |\DeclareCaptionFont{doublespacing}{\doublespacing}|\\
%   |\captionsetup{font={onehalfspacing,small},labelfont=bf}|
% \end{quote}
% \example{font={onehalfspacing,small},labelfont=bf,singlelinecheck=off}{\exampletext}
% \fi
% 下面是一个引入颜色的示例：
% \begin{quote}
%   |\usepackage{color}|\\
%   |\DeclareCaptionFont{red}{\color{red}}|\\
%   |\DeclareCaptionFont{green}{\color{green}}|\\
%   |\DeclareCaptionFont{blue}{\color{blue}}|\\
%   |\captionsetup{labelfont={blue,bf},textfont=green}|
% \end{quote}
% \example*{labelfont={color=blue,bf},textfont={color=green},singlelinecheck=off}{\exampletext}
% 但由于 \thispackage\ 已内置了精巧定义
% \begin{quote}
%   |\DeclareCaptionFont{color}{\color{#1}}|
% \end{quote}
% 只需下面的代码即可得到相同结果：
% \begin{quote}
%   |\usepackage{color}|\\
%   |\captionsetup{labelfont={color=blue,bf},|\\
%   |               textfont={color=green}}|\quad。
% \end{quote}
%
% \medskip
%
% \pagebreak[3]
% \DescribeMacro{\DeclareCaptionStyle}
% 可使用命令定义自己的题注样式：
% \begin{quote}
%   |\DeclareCaptionStyle|\marg{名称}\oarg{附加选项}\marg{选项}\mbox{\quad.}
% \end{quote}
% 回忆一下，题注样式只是一组搭配得当的选项，以给定名称保存起来。
% 随时可用选项 |style=|\meta{样式名称} 启用这些选项。
%
% 所有题注样式都以 |base| 选项集合为基础。（完整列表见 \longref{style}。）
% 因此只需指定与基础集合不同的选项。
%
% 若指定了 \meta{附加选项}，当题注恰好排成一行、且未用
% |single|\-|line|\-|check=off| 关闭单行检查时，这些选项会一并使用。
%
% 再举一个取自本宏包内核的简单示例：题注样式 |base| 预定义为
% \begin{quote}
%   |\DeclareCaptionStyle{base}%|\\
%   |        [justification=centering,indention=0pt]{}|\quad。
% \end{quote}
%
% 更精彩一些的例子：
% \begin{quote}
%   |\DeclareCaptionStyle{mystyle}%|\\
%   |        [margin=5mm,justification=centering]%|\\
%   |        {font=footnotesize,labelfont=sc,margin={10mm,0mm}}|\\
%   |\captionsetup{style=mystyle}|
% \end{quote}
% 得到如下题注：
% \begin{Example}
%   \captionsetup{type=figure,style=mystyle,position=b}
%   \caption{一则简短题注。}
%   \caption{一则很长很长很长很长很长很长很长很长很长
%              很长很长很长很长很长很长很长很长很长很长
%              很长很长很长很长很长很长很长很长的题注。}
% \end{Example}
%
% \bigskip
%
% \pagebreak[3]
% \DescribeMacro{\DeclareCaptionListFormat}
% \NEWfeature*{v3.1}
% 可使用命令定义自己的题注目录格式：
% \begin{quote}
%   |\DeclareCaptionListFormat|\marg{名称}\marg{使用 \#1、\#2 的代码}\mbox{\quad.}
% \end{quote}
% 使用时 \#1 替换为标签前缀（例如 |\p@figure|），\#2 替换为编号
% （例如 |\thefigure|）。
%
% \iffalse
% --------------------------------------------------------------------------- %
% \fi
%
% \subsection{更多示例}
% \label{examples}
%
% \subsubsection*{示例 1}
%
% 若希望题注标签与文字之间换行，可这样定义：
% \begin{quote}
%   |\DeclareCaptionFormat{myformat}{#1#2\\#3}|
% \end{quote}
% 用 |\caption|\-|setup{format=|\x|myformat}| 选中该格式后，得到如下题注：
%   \example{format=myformat1,labelfont=bf}{\exampletext}
% 该题注格式甚至还可以搭配缩进：
% \begin{quote}
%   |\captionsetup{format=myformat,indention=1cm}|
% \end{quote}
% 外观如下：
%   \example{format=myformat1,indention=1cm,labelfont=bf}{\exampletext}
% 若只希望题注文字第一行缩进呢？同样没问题，例如：
% \begin{quote}
%   |\newlength\myindention|\\
%   |\DeclareCaptionFormat{myformat}%|\\
%   |               {#1#2\\\hspace*{\myindention}#3}|\\
%   \ldots\\
%   |\setlength\myindention{1cm}|\\
%   |\captionsetup{format=myformat}|
% \end{quote}
% 得到如下题注：
%   \example{format=myformat2,myindention=1cm,labelfont=bf}{\exampletext}
% 若还想把这一缩进做成题注选项，从而能例如用
% |\caption|\-|setup|\x|[figure]|\x|{my|\-|indention=|\x\ldots|}|
% 来设置，也可以做到：
% \begin{quote}
%   |\newlength\myindention|\\
%   |\DeclareCaptionOption{myindention}%|\\
%   |               {\setlength\myindention{#1}}|\\
%   |\DeclareCaptionFormat{myformat}%|\\
%   |               {#1#2\\\hspace*{\myindention} #3}|\\
%   \ldots\\
%   |\captionsetup{format=myformat,myindention=1cm}|
% \end{quote}
% 结果与上例相同。
%
% \subsubsection*{示例 2}
%
% 希望题注外观如下：
%   \example{format=reverse,labelformat=fullparens,labelsep=fill,labelfont=it}{\exampletext}
%
% \pagebreak[2]
% 可这样实现：
% \nopagebreak[3]
% \begin{quote}
%   |\DeclareCaptionFormat{reverse}{#3#2#1}|\\
%   |\DeclareCaptionLabelFormat{fullparens}%|\\
%   |               {(\bothIfFirst{#1}{~}#2)}|\\
%   |\DeclareCaptionLabelSeparator{fill}{\hfill}|\\
%   |\captionsetup{format=reverse,labelformat=fullparens,|\\
%   |              labelsep=fill,font=small,labelfont=it}|
% \end{quote}
%
% \subsubsection*{示例 3}
%
% 希望题注文字伸入左页边距？一种可行方案如下：
% \begin{quote}
%   |\DeclareCaptionFormat{llap}{\llap{#1#2}#3\par}|\\
%   |\captionsetup{format=llap,labelsep=quad,singlelinecheck=no}|
% \end{quote}
% 结果得到如下题注：
% \example{format=llap,labelsep=quad,singlelinecheck=no,margin=0pt}{\exampletext}
%
% \medskip
%
% 若伸入页边距的缩进量需要固定，可组合使用 |\llap| 与 |\makebox|，例如：
% \begin{quote}
%   |\DeclareCaptionFormat{llapx}%|\\
%   |               {\llap{\makebox[5em][l]{#1}}#3\par}|\\
%   |\captionsetup{format=llapx,singlelinecheck=off}|
% \end{quote}
% \example{format=llapx,singlelinecheck=off,skip=0pt,margin=0pt}{\exampletext}
%
% \subsubsection*{示例 4}
%
% 本例把插图与表格并排放置，但两者共用一则题注。通过组合
% |\Declare|\-|Caption|\-|Label|\-|Format| 与 |\caption|\-|list|\-|entry|
% 实现：
% \begin{quote}
%   |\DeclareCaptionLabelFormat{andtable}%|\\
%   |               {#1~#2 \& \tablename~\thetable}|\\
%   \ldots\\
%   |\begin{figure}|\\
%   |  \centering|\\
%   |  \includegraphics{|\ldots|}%|\\
%   |  \qquad|\\
%   |  \begin{tabular}[b]{\ldots}|\\
%   |    |\ldots\\
%   |  \end{tabular}|\\
%   |  \captionlistentry[table]{|\ldots|}|\\
%   |  \captionsetup{labelformat=andtable}|\\
%   |  \caption{|\ldots|}|\\
%   |\end{figure}|
% \end{quote}
% \DeleteShortVerb{\|}%
% \noindent\begin{minipage}{\linewidth}
%   \captionsetup{type=figure}
%   \centering
%   \includegraphics[width=30pt]{cat}%
%   \qquad
%   \begin{tabular}[b]{r|rr}
%       & x & y \\\hline
%     a & 1 & 2 \\
%     b & 3 & 4 \\
%   \end{tabular}
%   \captionlistentry[table]{插图与表格共用一个题注}
%   \captionsetup{labelformat=andtable}
%   \caption{插图与表格共用一个题注~\footnotemark}
% \end{minipage}
% \footnotetext{本图片经许可取自 \LaTeX\ Companion\cite{TLC2} 示例。}
% \MakeShortVerb{\|}
%
% （请记住，|\caption|\-|list|\-|entry| 会使插图或表格计数器加一。）
%
% \iffalse
% --------------------------------------------------------------------------- %
% \fi
%
% \clearpage
% \section{支持的文档类}
% \label{classes}
%
% 本节概述 \pkg{caption} 支持哪些文档类：它们在题注方面已提供哪些功能、
% 搭配 \pkg{caption} 宏包使用时会产生哪些副作用，以及各自的“默认”设置
% 是什么。
%
% \NEWfeature{v3.1}
% “默认”设置取决于所使用的文档类，它代表文档类作者所期望的题注外观。
% 因此，例如设置 |format=|\x|default|，在不同文档类下可能得到不同的
% 视觉设计。
%
% \begingroup\setlength\leftmargini{0.3em}% default = 2.5em
% \INFO
% 若本节中找不到你所用的文档类，通常无需担心：许多文档类（例如
% \cls{octavo}）派生自标准文档类 \cls{article}、\cls{report} 或
% \cls{book} 之一，在题注方面行为相同。\Thispackage\ 会自动对照所用
% 文档类进行兼容性检查；若检测到不兼容，将给出明确警告
% \begin{quote}\footnotesize
%   |Package caption Warning: Unsupported document class (or package) detected,|\\
%   |(caption)                usage of the caption package is not recommended.|\\
%   |See the caption package documentation for explanation.|
% \end{quote}
% 未收到此类警告即一切正常；若收到，则不建议使用 \thispackage，
% 尤其不提供支持。
% \endgroup
%
% \begin{Expert}
% 若收到此类兼容性警告，但仍决定使用 \thispackage{}，应仔细观察出现的
% 副作用：通常仅加载不带选项的 \thispackage{}，题注外观就会改变，即不再
% 呈现文档类作者本意的样子。若可以接受，应先通过
% |\use|\-|package[style=|\x|base]{caption}| 或
% |\caption|\-|setup{style=|\x|base}| 指定选项 |style=|\x|base|，把
% \thispackage\ 置于定义明确的状态；之后再逐步添加自己的选项，并
% 祝好运。
% \end{Expert}
%
% \newcommand*\Option{选项}
% \newcommand*\defaultvalue{默认取值}
% \newcommand*\uses{\textit{采用}}
% \newcommand*\settings{\textit{设置}}
% \newcommand*\nofont{\textit{无}}
%
% \subsection{标准 \LaTeX：article、report 与 book}
%
% \begin{tabular}{ll}
% \Option          & \defaultvalue \\\hline
% |format=|        & |plain| \\
% |labelformat=|   & |simple| \\
% |labelsep=|      & |colon| \\
% |justification=| & |justified| \\
% |font=|          & \nofont \\
% |labelfont=|     & \nofont \\
% |textfont=|      & \nofont \\
% \end{tabular}
%
% \begin{Annotation}
% 派生自这些文档类的其他文档类同样适用。
% \end{Annotation}
%
% \subsection{\AmS：amsart、amsproc 与 amsbook}
% \label{AMS}
%
% \begin{tabular}{ll}
% \Option          & \defaultvalue \\\hline
% |format=|        & |plain| \\
% |labelformat=|   & |simple| \\
% |labelsep=|      & |.\enspace| \\
% |justification=| & |justified| \\
% |font=|          & |\@captionfont| \\
% |labelfont=|     & |\@captionheadfont| \\
% |textfont=|      & |\@captionfont\upshape| \\
% \end{tabular}
%
% \begin{Annotation*}
% \AmS\ 文档类会把 |\@caption|\-|font| 设为 |\normal|\-|font|，把
% |\@caption|\-|head|\-|font| 设为 |\sc|\-|shape|。
% \end{Annotation*}
%
% 此外，多行题注的边距设为 |\caption|\-|indent|（\AmS\ 类将其设为
% |3pc|），单行题注的边距则为它的一半。若希望两者使用相同边距，可在
% 加载 \thispackage 之后，于文档导言区插入
% |\clear|\-|caption|\-|setup[margin*]{single|\-|line}|。
%
% 另外还会设置选项 |figure|\-|position=b,|\x|table|\-|position=t|。
% 在加载 \thispackage 时，可在选项列表中为 |figure|\-|position=| 或
% |table|\-|position=| 指定其他取值来覆盖这些设置。
%
% \subsection{beamer 类}
% \label{beamer}
%
% \begin{tabular}{ll}
% \Option          & \defaultvalue \\\hline
% |format=|        & |plain| \\
% |labelformat=|   & \textit{不编号} \\
% |labelsep=|      & |colon| \\
% |justification=| & |raggedright| \\
% |font=|          & \cls{beamer} “题注” \settings \\
% |labelfont=|     & \cls{beamer} “题注名称” \settings \\
% |textfont=|      & \nofont \\
% \end{tabular}
%
% \subsubsection*{内置功能与副作用}
% 可用 |\set|\-|beamer|\-|font{caption}|\marg{选项} 与
% |\set|\-|beamer|\-|font{caption name}|\marg{选项} 设置字体与颜色；（其中 |caption name| 即 beamer 的“题注名称”模板。）
% 除非再用 |\caption|\-|setup{font=|\x\meta{选项}|}| 或
% |\caption|\-|setup{label|\-|font=|\x\meta{选项}|}| 设置了不同字体，否则
% 这些设置仍然有效。\par
% 此外，\cls{beamer} 类提供不同的题注模板，可用
% |\set|\-|beamer|\-|template|\x|{caption}|\x|[|\meta{模板}|]| 选择。
% 由于 \thispackage 取代了这一题注模板机制，使用 \thispackage 时，
% |\defbeamer|\-|template|*|\x|{caption}|\x\marg{模板代码}| 与
% |\set|\-|beamer|\-|template|\x|{caption}|\x|[|\meta{模板}|]| 均不再
% 生效。
% （例外：选择模板 |default|、|numbered|，或“题注名称独占一行”时，\thispackage 会识别并映射到相应选项。）
%
% \subsection{\KOMAScript：scrartcl、scrreprt 与 scrbook}
% \label{KOMA}
%
% \begin{tabular}{ll}
% \Option          & \defaultvalue \\\hline
% |format=|        & \uses\ |\setcapindent| \textit{\&} |\setcaphanging| \settings \\
% |labelformat=|   & \textit{类似 \purett{simple}，但支持自动句点 “autodot”}\\
% |labelsep=|      & |\captionformat| \\
% |justification=| & |justified| \\
% |font=|          & |\setkomafont{caption}| \settings \\
% |labelfont=|     & |\setkomafont{captionlabel}| \settings \\
% |textfont=|      & \nofont \\
% \end{tabular}
%
% \subsubsection*{内置功能}
% \KOMAScript\ 文档类提供了多种自定义题注外观的方式。概览与完整说明
% 请参阅 \KOMAScript\ 文档中 “表格与插图” 一节。
%
% \subsubsection*{副作用}
% |\set|\-|cap|\-|width| 的可选参数不受支持；与 \pkg{caption} 宏包
% 搭配使用时将被忽略。此外，\KOMAScript\ 选项
% |table|\-|caption|\-|above| 与 |table|\-|caption|\-|below|，以及
% 命令 |\caption|\-|above| 与 |\caption|\-|below|，优先级高于
% \pkg{caption} 宏包提供的 |position=| 设置。
%
% \subsection{\NTG：artikel、rapport 与 boek}
% \label{NTG}
%
% \begin{tabular}{ll}
% \Option          & \defaultvalue \\\hline
% |format=|        & |plain| \\
% |labelformat=|   & |simple| \\
% |labelsep=|      & |colon| \\
% |justification=| & |justified| \\
% |font=|          & \nofont \\
% |labelfont=|     & |\CaptionLabelFont| \\
% |textfont=|      & |\CaptionTextFont| \\
% \end{tabular}
%
% \subsubsection*{内置功能与副作用}
% |\Caption|\-|Label|\-|Font| 与 |\Caption|\-|Text|\-|Font| 既可直接设置，
% 也可通过 |\Caption|\-|Fonts| 设置。除非使用 \thispackage 提供的
% |label|\-|font=| 或 |text|\-|font=| 这两个选项之一，否则两种方式
% 都仍然受支持。
%
% \subsection{\SmF：smfart 与 smfbook}
% \label{SMF}
%
% 由于 \SmF\ 文档类派生自 \AmS\ 文档类，这里沿用相同的默认取值。
%
% 此外，边距被限制为 |\line|\-|width| 的十分之一以内。若不喜欢此限制，
% 可用选项 |max|\-|margin=|\x|off| 或 |max|\-|margin=|\x|false| 关闭
% （二者含义相同）。
%
% \subsection{thesis 类}
% \label{thesis}
%
% \begin{tabular}{ll}
% \Option          & \defaultvalue \\\hline
% |format=|        & |hang| \\
% |labelformat=|   & \textit{同 \purett{simple}，但使用短名称}\\
% |labelsep=|      & |colon| \\
% |justification=| & |justified| \\
% |font=|          & \nofont \\
% |labelfont=|     & |\captionheaderfont| \\
% |textfont=|      & |\captionbodyfont| \\
% \end{tabular}
%
% \subsubsection*{内置功能与副作用}
% 题注标签字体可用 |\caption|\-|header|\-|font| 设置，题注文字字体可用
% |\caption|\-|body|\-|font| 设置。除非使用 \thispackage 提供的
% |label|\-|font=| 或 |text|\-|font=| 选项之一，否则两种方式都仍受支持。
%
% \iffalse
% --------------------------------------------------------------------------- %
% \fi
%
% \clearpage
% \section{支持的 babel 选项}
% \label{babel}
%
% \NEWfeature{v3.5}
% 自 \pkg{caption} 宏包 \version{3.5} 起，相关支持不再依赖
% \pkg{caption} 与 \pkg{babel} 的加载先后顺序。
%
% \subsection{french 语言 babel 选项}
% \label{french}
%
% 若在三种标准 \LaTeX{} 文档类之一（或其派生类）中使用 \pkg{babel}
% 宏包的 \pkg{french} 选项，默认的 |label|\-|sep=| 将被设为
% |\Caption|\-|Sep|\-|a|\-|ra|\-|tor|（由 \pkg{french} 提供），覆盖
% 文档类设定的默认值。因此，只要不选择不同于默认值的 |label|\-|sep=|，
% 重定义 |\Caption|\-|Sep|\-|a|\-|ra|\-|tor| 仍然有效。
%
% \iffalse
% --------------------------------------------------------------------------- %
% \fi
%
% \clearpage
% \section{支持的宏包}
% \label{packages}
% \label{compatibility}
%
% \Thispackage\ 已与下列同样处理题注的宏包做了适配：
% \begin{quote}
%   \pkg{float}、\pkg{floatflt}、\pkg{fltpage}、
%   \pkg{hyperref}、\pkg{hypcap}、\pkg{listings}、
%   \pkg{longtable}、\pkg{picinpar}、\pkg{picins}、
%   \pkg{rotating}、\pkg{setspace}、\pkg{sidecap}、
%   \pkg{subfigure}、\pkg{supertabular}、\pkg{threeparttable}、
%   \pkg{wrapfig} 与 \pkg{xtab}
% \end{quote}
%
% 此外，\pkg{floatrow} 宏包\cite{floatrow}、\pkg{subcaption} 宏包
% （属于 \pkg{caption} 宏包套件的一部分）以及 \pkg{subfig}
% 宏包\cite{subfig} 支持 \thispackage{}，并使用其
% |\caption|\-|setup| 接口。
%
% \bigskip
%
% \begingroup\setlength\leftmargini{0.3em}% default = 2.5em
% \INFO
% 若 \thispackage 不认识的某个宏包（或文档类）也重定义了 |\caption|
% 命令，则该重定义优先于本宏包提供的版本，以实现最大兼容性、避免冲突。
% 若检测到这种潜在不兼容，你会看到如下警告：~\footnote{%
%   加载 \thispackage 时指定选项 \texttt{compatibility=true} 可关闭此警告。}
% \begin{quote}\footnotesize
%   |Package caption Warning: \caption will not be redefined since it's already|\\
%   |(caption)                redefined by a document class or package which is|\\
%   |(caption)                unknown to the caption package.|\\
%   |See the caption package documentation for explanation.|
% \end{quote}
% \endgroup
%
% 结果是，\thispackage 提供的下列功能将不可用：
% \begin{itemize}
%   \item 选项 |labelformat=|、|position=auto|、|list=| 与 |listformat=|
%   \item |\caption*|（生成不带标签的题注）
%   \item |\caption[]{|\ldots|}|（不在插图目录或表格目录中生成条目）
%   \item |\caption{}|（生成不带标签分隔符的空题注）
%   \item |\ContinuedFloat|
%   \item 在 \texttt{wide}、\texttt{addmargin} 等附加边距的环境中
%         正确排版对齐的题注
%   \item \textsf{hypcap} 功能 \See{\longref{hyperref}\,}
%   \item 子题注功能 \See{\pkg{subcaption} 宏包文档}
% \end{itemize}
%
% \DescribeMacro{compatibility=}
% 加载 \thispackage 时，可通过指定选项
% \begin{quote}
%   |compatibility=false|
% \end{quote}
% 覆盖这种兼容模式。但请注意，既不推荐、也不支持使用该选项，因为其后
% 可能出现意外副作用甚至错误。（为此会给出相应警告。）
%
% \newcommand\packagedescription[1]{%
%   \ifvmode\else\par\fi
%   \nopagebreak
%   \parbox[b]{\linewidth}{\footnotesize\leftskip=10pt\rightskip=10pt\relax#1}\par
%   \nopagebreak\smallskip\nopagebreak}
%
% \PageBreak
% \subsection{algorithms 宏包}
% \label{algorithms}
% \packagedescription{%
% \pkg{algorithms} 宏包套件\cite{algorithms}提供两个环境：\texttt{algorithmic}
% 环境用于描述算法，\texttt{algorithm} 环境则为算法提供“浮动体”包装。}
%
% 由于 \texttt{algorithm} 环境是借助 \pkg{float} 宏包\cite{float} 提供的
% |\new|\-|float| 实现的，请参见 \longref{float}。
%
% \iffalse
% \pagebreak[3]
% \subsection{algorithm2e 宏包}
% \label{algorithm2e}
% \packagedescription{%
% \pkg{algorithm2e} 宏包\cite{algorithm2e} 提供在 LaTeX2e 中编写算法的环境。}
%
% 由于 \pkg{algorithm2e} 宏包定义浮动环境、排版题注时并未采用某种标准
% 接口（例如 \pkg{float} 宏包的 |\newfloat|），\thispackage\ 对
%  \pkg{algorithm2e} 宏包\emph{不}提供支持。
%
% 因此，若希望用 \thispackage 定制题注，必须新建一个在浮动体方面采用受
% 支持接口（例如借助 \pkg{float} 或 \pkg{floatrow} 宏包）的环境，再把该
% 环境与 \pkg{algorithm2e} 的内部代码结合起来。
%
% 第一步可先定义一个非浮动环境 \texttt{algorithmic}，例如：
% \begin{quote}
%   |\usepackage{algorithm2e}|\\
%   |% 保存 algorithm2e 宏包的 "algorithm" 环境|\\
%   |\let\ORIGalgorithm\algorithm|\\
%   |\let\ORIGendalgorithm\endalgorithm|\\
%   |% 定义 algorithmic 环境|\\ % 基于保存的环境
%   |\newenvironment{algorithmic}%|\\
%   |  {\renewenvironment{algocf}[1][h]{}{}% 略过浮动相关内容|\\
%   |   \ORIGalgorithm}%|\\
%   |  {\ORIGendalgorithm}|\\
%   \ldots
% \end{quote}
% 定义该非浮动环境后，可用 \pkg{float} 宏包的 |\newfloat|（或
%  \pkg{floatrow} 宏包的 |\DeclareNewFloatType|）定义自己的新浮动环境，
% 并在其正文内组合使用该浮动环境与 \texttt{algorithmic}，就像使用
%  \pkg{algorithm}/\pkg{algorithmic} 这对宏包时那样。
%
% 甚至可借助 \pkg{algorithm} 宏包完成此目的，例如：
% \begin{quote}
%   \ldots\\
%   |% 加载 algorithm 宏包以重定义|\\
%   |% 浮动环境 "algorithm" 与 \listofalgorithms|\\
%   |\let\listofalgorithms\undefined|\\
%   |\usepackage{algorithm}|\\
%   \ldots\\
%   |% 用法示例：|\\
%   |\begin{algorithm}|\\
%   |\caption{一个算法}|\\
%   |\begin{algorithmic}|\\
%   |  \SetVline|\\
%   |  \eIf{cond1}{|\\
%   |    一行代码\;|\\
%   |  }{|\\
%   |    另一行代码\;|\\
%   |  }|\\
%   |\end{algorithmic}|
% \end{quote}
% 至此即可像其他用 \pkg{float} 宏包\cite{float} 定义的浮动环境那样定制
% 浮动体与题注版面；请参见 \longref{float}。
%
% \begin{Annotation}
% 另一种替代方案是使用 \pkg{algorithmicx} 宏包。
% \end{Annotation}
% \fi
%
% \pagebreak[3]
% \subsection{float 宏包}
% \label{float}
% \packagedescription{%
% \pkg{float} 宏包\cite{float} 提供命令 \cs{restylefloat}，用于赋予已有
% 浮动环境新的外观；以及 \cs{newfloat}，用于定义新的浮动环境。它还提供
% “\texttt{H}” 浮动体放置选项，把环境固定在“此处”，而不是让它四处
% 浮动。}
%
% 对于用 \cs{newfloat} 或 \cs{restylefloat} 定义的浮动环境，
% |po|\-|si|\-|tion| 选项对主题注不再生效，因为其放置位置与间距改由
% 所选浮动体样式控制。
%
% 以浮动体样式名称定义的题注样式与选项，会在常规选项之外一并执行。
% 借助这一机制，\thispackage 模拟了 |ruled| 题注的默认外观：它定义了
% 如下题注样式
% \begin{quote}
%   |\DeclareCaptionStyle{ruled}%|\\
%   |       {labelfont=bf,labelsep=space,strut=off}|\mbox{\quad.}
% \end{quote}
% 若要更改，要么定义自己的、名为 |ruled| 的题注样式，要么使用
% |\caption|\-|setup[ruled]|\marg{选项} 指定附加选项。
%
% 同一机制还用于指定 |boxed| 浮动体与其题注之间的间距，覆盖全局取值：
% \begin{quote}
%   |\captionsetup[boxed]{skip=2pt}|
% \end{quote}
% 要更改它，只需用 |\caption|\-|setup[boxed]{skip=|\x\meta{取值}|}| 给定
% 适当数值。若希望改用全局 |skip| 设置，可用
% |\clear|\-|caption|\-|setup[skip]{boxed}| 移除这些浮动体的局部设置。
%
% \begin{Note}
% 用 |\new|\-|float| 或 |\re|\-|style|\-|float| 定义的环境中只能排版
% 一则题注；而且这些环境的行为与预定义浮动体 |figure|、|table| 并不完全
% 相同。因此许多宏包与它们配合得并不好。此外，\pkg{float} 宏包存在一些
% 注意事项与限制：若只是想定义一个行为类似 |figure| 或 |table| 的简单
% 新浮动环境，建议改用 \pkg{newfloat} 宏包提供的
% |\Declare|\-|Floating|\-|Environment|。若要定义非简单的浮动环境并做
% 定制，建议使用 \pkg{floatrow} 宏包\cite{floatrow} 提供的
% |\Declare|\-|New|\-|Float|\-|Type|。
% \end{Note}
%
% \iffalse
% 另请注意，使用“|H|”浮动体放置符并不需要 |\restyle|\-|float|。有些
% 文档如此声称，但并不正确。而 |\restyle|\-|float{table}| 很容易弄巧成拙：
% 许多内部使用 |table| 的宏包在其后将无法正常工作。
% \fi
%
% \pagebreak[3]
% \subsection{floatflt 宏包}
% \label{floatflt}
% \packagedescription{%
% \pkg{floatflt} 宏包\cite{floatflt} 提供不必占满整页宽度、可被文字环绕
% 浮动的插图与表格。}
%
% \NEWfeature{v3.1}
% 若要为 |float|\-|ing|\-|figure| 与 |float|\-|ing|\-|table| 环境设置
% 专门选项，可使用
% \begin{quote}
%   |\captionsetup[floatingfigure]|\marg{选项}\quad\emph{以及}\\
%   |\captionsetup[floatingtable]|\marg{选项}\quad。
% \end{quote}
% 这些选项会在 |figure| 或 |table| 的常规选项之外一并执行。
%
% \begin{Note}
% 除非通过 |\caption|\-|setup[floating|\-|figure]{|\ldots|}| 或
%  |\caption|\-|setup[floating|\-|table]{|\ldots|}| 显式设置，否则这些插图
%  或表格不会采用 |margin| 或 |width| 设置。
% \end{Note}
%
% \pagebreak[3]
% \subsection{fltpage 宏包}
% \label{fltpage}
% \packagedescription{%
% \pkg{fltpage} 宏包\cite{fltpage} 为内容需要占满整页的 |figure| 或
% |table| 提供题注的外送机制：把题注放在前一页或后一页的底部。}
%
% \iffalse
% \NEWfeature{v3.1}
% 有两个选项控制指向 |FPfigure| 与 |FPtable| 环境的链接：
% \begin{description}
% \item{\texttt{FPlist=caption} or \texttt{FPlist=figure}}\\[\smallskipamount]
% 设为 |caption| 时，目录条目链接到题注；设为 |figure| 时，链接到插图内容。
% （默认取值为 |FP|\-|list=|\x|caption|。）
% \item{\texttt{FPref=caption} or \texttt{FPref=figure}}\\[\smallskipamount]
% 设为 |caption| 时，\cs{ref}、\cs{pageref}、\cs{autoref} 或 \cs{nameref}
% 链接到题注；设为 |figure| 时，链接到插图内容。
% （默认取值为 |FP|\-|ref=|\x|figure|。）
% \end{description}
% \fi
%
% 若要为 |FPfigure| 与 |FPtable| 环境设置专门选项，可使用
% \begin{quote}
%   |\captionsetup[FPfigure]|\marg{选项}\quad\emph{以及}\\
%   |\captionsetup[FPtable]|\marg{选项}\quad。
% \end{quote}
% 这些选项会在 |figure| 或 |table| 的常规选项之外一并执行。
%
% \pagebreak[3]
% \subsection{frenchle 与 frenchpro 宏包}
% \label{frenchpro}
%
% 若使用 \pkg{frenchle} 或 \pkg{frenchpro} 宏包，默认的 |label|\-|sep=|
% 将被设为 |\caption|\-|sep|\-|a|\-|ra|\-|tor|（由 \pkg{frenchle/pro}
% 提供）加上 \cs{space}，覆盖文档类设定的默认值。因此，只要不选择不同于
% 默认值的 |label|\-|sep=|，重定义 |\caption|\-|sep|\-|a|\-|ra|\-|tor|
% 仍然有效。
%
% 此外，默认 |text|\-|font=| 将被设为 |text|\-|font=|\x|it|，因为这
% 模拟了 \pkg{frenchle} 或 \pkg{frenchpro} 宏包定义的 |\caption|\-|font|
% 默认设置。请注意，|\caption|\-|font| 命令在 \pkg{caption} 宏包内部
% 另作他用，因此不应再更改它。
%
% 命令 |\un|\-|numbered|\-|captions|\marg{插图 \emph{或} 表格} 仍然有效，
% 但前提是不选择不同于默认值的 |label|\-|format=|。
%
% \NEWfeature{v3.5}
% 自 \pkg{caption} 宏包 \version{3.5} 起，相关支持不再依赖
% \pkg{caption} 与 \pkg{frenchle}（或 \pkg{frenchpro}）的加载先后顺序。
%
% \pagebreak[3]
% \subsection{hyperref 宏包}
% \label{hyperref}
% \packagedescription{%
% \pkg{hyperref} 宏包\cite{hyperref} 用于处理 \LaTeX 中的交叉引用命令，
% 在文档中生成超文本链接。}
%
% \NEWfeature{v3.1}
% 有两个选项可用于控制超链接锚点的位置：~\footnote{这两个选项得名于
% 被它们取代的 \pkg{hypcap} 宏包。}
% \begin{description}
% \item{\texttt{hypcap=true} 或 \texttt{hypcap=false}}\\[\smallskipamount]
% 设为 |true| 时，所有超链接锚点——插图目录条目、|\ref| 与
% |\auto|\-|ref| 将链接到的位置——都放在（浮动）环境的开头。\par
% 设为 |false| 时，超链接锚点（通常）放在题注处。\par
% （默认取值为 |hypcap=|\x|true|。）
% \item{\texttt{hypcapspace=}\meta{长度}}\\[\smallskipamount]
% 若超链接恰好指向插图或表格的顶部，外观欠佳；可指定超链接锚点与
% （浮动）环境本身之间的垂直距离，例如 |hyp|\-|cap|\-|space=|\x|0pt|
% 可取消该距离。\par
% （默认取值为 |hyp|\-|cap|\-|space=|\x|0.5\base|\-|line|\-|skip|。）
% \end{description}
%
% 这两个设置在 |lst|\-|listing|（\pkg{listings} 宏包提供）、
% |long|\-|table|（\pkg{longtable} 宏包提供）、|super|\-|tabular|
% （\pkg{supertabular} 宏包提供）与 |x|\-|tabular|（\pkg{xtab}
% 宏包提供）中均无效；在这些环境中，超链接锚点始终按设置了
% |hyp|\-|cap=|\x|true| 与 |hyp|\-|cap|\-|space=|\x|0pt| 来放置。
%
% \pagebreak[3]请注意：\nopagebreak
% \begin{description}
% \item{\cs{captionof}\marg{类型}\csmarg{\purerm\ldots}
%   与~\cs{captionsetup}\csmarg{type=\textrm{\meta{类型}}}$+$\cs{caption}\csmarg{\purerm\ldots}}%
% \\[\smallskipamount]
% 未加载 \pkg{hyperref} 时，两者结果相同。但加载 \pkg{hyperref}、并要求
% |hyp|\-|cap=|\x|true| 时，超链接锚点的放置位置不同。例如：
% \begin{quote}
% |\begin{minipage}{\linewidth}|\\
% |  |\ldots\\
% |  \captionof{figure}{一个插图}|\\
% |\end{minipage}|
% \end{quote}
% 会把超链接锚点放在题注处。
% （若设置了 |hyp|\-|cap=|\x|true|，会就此给出警告。）
% 而
% \begin{quote}
% |\begin{minipage}{\linewidth}|\\
% |  \captionsetup{type=figure}|\\
% |  |\ldots\\
% |  \caption{一个插图}|\\
% |\end{minipage}|
% \end{quote}
% 会把超链接锚点放在 |mini|\-|page| 开头，因为
% |\caption|\-|setup{type=|\x|figure}| 不仅把题注类型设为 “figure”，
% 还会放置一个超链接锚点。
%
% \item{\cs{caption}\csoarg{}\csmarg{\purerm\ldots}
%   与~\cs{captionsetup}\csmarg{list=false}$+$\cs{caption}\csmarg{\purerm\ldots}}\\[\smallskipamount]
% 同样，未加载 \pkg{hyperref} 时两者结果相同；加载后差别在于细微之处。
% 例如，|\caption| 的可选参数还会写入 |aux| 文件，并被 |\name|\-|ref|
% 命令使用。因此若选用带空可选参数的 |\caption|，|\name|\-|ref| 也会得到
% 空结果。若不希望在插图目录或表格目录中生成条目，最好使用
% |\caption|\-|setup{list=|\x|false}|。
% \end{description}
%
% \pagebreak[3]
% \subsection{hypcap 宏包}
% \label{hypcap}
% \packagedescription{%
% \pkg{hyp\-cap} 宏包\cite{hypcap} 解决这样一个问题：使用 \pkg{hyper\-ref}
% 链接到浮动体时，链接锚点可能落在题注而非浮动体开头。自 \thispackage\
% \version{3.1} 起已自行解决该问题，通常不再需要 \pkg{hypcap} 宏包。}
%
% 若在 \pkg{hyper\-ref} 宏包之外再加载 \pkg{hyp\-cap} 宏包，它将从
% \thispackage 手中接管超链接锚点位置的控制，覆盖选项 |hyp|\-|cap=|
% 与 |hyp|\-|cap|\-|space=|。
%
% 因此，要手动放置超链接锚点，|\caption|\-|setup{type=|\x\meta{类型}|}|
% 已不再足够；此时需改用 |\cap|\-|start|（由 \pkg{hyp\-cap} 宏包提供）。
%
% 在自动放置方面，\pkg{hypcap} 宏包仅为浮动环境 |figure| 与 |table|
% 提供良好的锚点放置；相比之下，\thispackage 的 |hypcap=|\x|true|
% 选项还能为下列环境提供良好的锚点放置：|floating|\-|figure|
% （\pkg{float\-flt} 宏包）、|FP|\-|figure| 与 |FPtable|
% （\pkg{flt\-page} 宏包）、|fig|\-|window|（\pkg{pic\-in\-par} 宏包）、
% |par|\-|pic|（\pkg{pic\-ins} 宏包）、|SC|\-|figure|
% （\pkg{side\-cap} 宏包）、|three|\-|part|\-|table|
% （\pkg{three\-part\-table} 宏包）以及 |wrapfigure|
% （\pkg{wrap\-fig} 宏包）。
%
% \pagebreak[3]
% \subsection{listings 宏包}
% \label{listings}
% \packagedescription{%
% \pkg{listings} 宏包\cite{listings} 用于排版程序代码。}
%
% 若要为 |lst|\-|listing| 环境设置专门选项，可使用
% \begin{quote}
%   |\captionsetup[lstlisting]|\marg{选项}\quad。
% \end{quote}
%
% 请注意，\pkg{listings} 宏包有自己专门的选项控制题注位置与间距：
% |caption|\-|pos=|、|above|\-|caption|\-|skip=| 与
% |below|\-|caption|\-|skip=|。\See{详见 \pkg{listings} 文档。}
% 这些 \pkg{listings} 选项会覆盖 \pkg{caption} 的相应选项，但又可被
% |\caption|\-|setup[lst|\-|listing]{|\ldots|}| 再次覆盖，例如
% \begin{quote}|\caption|\-|setup[lst|\-|listing]{skip=|\x|10pt}|\quad。\end{quote}
%
% \pagebreak[3]
% \subsection{longtable 宏包}
% \label{longtable}
% \packagedescription{%
% \pkg{longtable} 宏包\cite{longtable} 提供一个行为类似 \texttt{tabular}
% 环境、但表格本身可跨多页的环境。}
%
% 若要为 |long|\-|table| 环境设置专门选项，可使用
% \begin{quote}
%   |\captionsetup[longtable]|\marg{选项}\quad。
% \end{quote}
% 这些选项会在 |table| 的常规选项之外一并执行。
%
% |margin| 与 |width| 设置通常会覆盖 |\LT|\-|cap|\-|width|，使
% |table| 与 |long|\-|table| 中的题注外观一致。但若把
% |\LT|\-|cap|\-|width| 设为不同于默认值 $=$ |4in| 的数值，
% \thispackage 会跟随该取值。（不过 |\LT|\-|cap|\-|width| 仍会被
% |\caption|\-|setup[long|\-|table]{width=|\x\meta{取值}|}| 覆盖，
% 即使其取值不同于 |4in|。）
%
% \begin{Note}
% |\caption|\-|of| 与 |\Continued|\-|Float| 对 |long|\-|table| \emph{不起作用}。
% 此外，此处 |\caption| 与 |\caption|\-|list|\-|entry| 都不会使 |table|
% 计数器加一；该计数器由 |long|\-|table| 环境自行递增。
% \NEWfeature{v3.1}
% 若需要不使 |table| 计数器加一的 |long|\-|table|，请使用
% |long|\-|table*| 环境（由 \pkg{ltcaption} 宏包提供，它属于
% \thispackage 套件，会自动加载）。
% \end{Note}
%
% \pagebreak[3]
% \subsection{picinpar 宏包}
% \label{picinpar}
% \packagedescription{%
% 与 \pkg{floatflt} 宏包类似，\pkg{picinpar} 宏包\cite{picinpar} 提供
% 不必占满整页宽度、可被文字环绕浮动的插图与表格。有关这些宏包之间
% 差异的详细讨论，请参阅 《\LaTeX\ Companion》\cite{TLC2}。}
%
% \NEWfeature{v3.1}
% 若要为 |fig|\-|window| 与 |tab|\-|window| 环境设置专门选项，可使用
% \begin{quote}
%   |\captionsetup[figwindow]|\marg{选项}\quad\emph{以及}\\
%   |\captionsetup[tabwindow]|\marg{选项}\quad。
% \end{quote}
% 这些选项会在 |figure| 或 |table| 的常规选项之外一并执行。
%
% \begin{Note}
% 除非显式用 |\caption|\-|setup[fig|\-|window]{|\ldots|}| 或
% |\caption|\-|setup[fig|\-|table]{|\ldots|}| 设置，否则这些插图与表格
% 不使用 |margin| 或 |width| 设置。
% \end{Note}
%
% \pagebreak[3]
% \subsection{picins 宏包}
% \label{picins}
% \packagedescription{%
% 与 \pkg{floatflt}、\pkg{picinpar} 宏包类似，\pkg{picins} 宏包\cite{picins}
% 提供不必占满整页宽度、可被文字环绕浮动的插图与表格。有关这些宏包
% 之间差异的详细讨论，请参阅 《\LaTeX\ Companion》\cite{TLC2}。}
%
% \NEWfeature{v3.1}
% 若要为 |par|\-|pic| 环境设置专门选项，可使用
% \begin{quote}
%   |\captionsetup[parpic]|\marg{选项}\quad。
% \end{quote}
% 这些选项会在 |figure| 或 |table| 的常规选项之外一并执行。
%
% 此外，|\pic|\-|caption[]{|\ldots|}| 不在插图目录中生成条目，
% |\pic|\-|caption*{|\ldots|}| 给出不带标签、不编号的题注。
%
% \begin{Note}
% 这些插图与表格不使用 |margin| 或 |width| 设置。
% \iffalse
% 当使用 |pic|\-|caption|\-|side| 或 |pic|\-|caption|\-|top|\-|side| 时。
% 当使用 |pic|\-|caption|\-|out|\-|side| 或 |pic|\-|caption|\-|in|\-|side|
% 时，不使用 |width| 设置（但使用 |margin| 设置）。
% \fi
% 但可通过 |\caption|\-|setup[par|\-|pic]{|\ldots|}| 显式设置来覆盖。
% \end{Note}
%
% 若希望 |\pic|\-|caption| 使用不同于 |figure| 的类型，请\emph{不要}
% 像 \pkg{picins} 文档建议的那样重定义 |\@cap|\-|type|；请改用
% \thispackage 提供的 |\pic|\-|caption|\-|type|\marg{类型} 命令，例如：
% \begin{quote}
%   |\piccaptiontype{table}|\\
%   |\piccaption{一个示例表格}|\\
%   |\parpic(50mm,10mm)[s]{|\ldots|}|
% \end{quote}
%
% \pagebreak[3]
% \subsection{rotating 宏包}
% \label{rotating}
% \packagedescription{%
% \pkg{rotating} 宏包\cite{rotating} 提供浮动环境
% \texttt{side\-ways\-figure} 与 \texttt{side\-ways\-table}，它们与普通
% 插图、表格类似，只是旋转了 90 或 270 度。此外还提供专门用于旋转
% 题注的命令 \cs{rot\-caption}。}
%
% \iffalse % 不，我们不这样做
% 若要为 |side|\-|ways|\-|figure| 与 |side|\-|ways|\-|table| 环境设置
% 专门选项，可使用
% \begin{quote}
%   |\captionsetup[sidewaysfigure]|\marg{选项}\\
%   |\captionsetup[sidewaystable]|\marg{选项}\quad。
% \end{quote}
% 这些选项会在 |figure| 或 |table| 的常规选项之外一并执行。
% \fi
%
% 命令 |\rot|\-|caption| 会由 \thispackage 扩展，因此
% |\rot|\-|caption*| 与 |\rot|\-|caption|\-|of| 可分别类比
% |\caption*| 与 |\caption|\-|of| 使用。
%
% \pagebreak[3]
% \subsection{setspace 宏包}
% \label{setspace}
% \packagedescription{%
% \pkg{setspace} 宏包\cite{setspace} 提供更改行距的选项与命令，
% 例如 \cs{usepackage}\x\csoarg{one\-half\-spacing}\x\csmarg{set\-space}
% 会使文档按一点五倍行距排版。}
%
% 若 \pkg{setspace} 宏包与 \thispackage 搭配使用，题注默认按单倍行距
% 排版。可通过指定 |font=|\x|one|\-|half|\-|spacing|、
% |font=|\x|double|\-|spacing| 或 |font={stretch=|\x\meta{系数}|}|
% 来更改。
% \See{另见 \longref{fonts}}
%
% \pagebreak[3]
% \subsection{sidecap 宏包}
% \label{sidecap}
% \packagedescription{%
% \pkg{sidecap} 宏包\cite{sidecap} 提供浮动环境
% \texttt{SC\-figure} 与 \texttt{SC\-table}；它们与普通插图、表格类似，
% 只是题注放在内容\emph{旁侧}。}
%
% 若要为 |SCfigure| 与 |SCtable| 环境设置专门选项，可使用
% \begin{quote}
%   |\captionsetup[SCfigure]|\marg{选项}\quad\emph{以及}\\
%   |\captionsetup[SCtable]|\marg{选项}\quad。
% \end{quote}
% 这些选项会在 |figure| 或 |table| 的常规选项之外一并执行。
%
% \pagebreak[3]
% \begin{Note}
% \pkg{sidecap} 宏包提供自己的对齐选项。一旦设置，它们会覆盖题注选项
% |jus|\-|ti|\-|fi|\-|ca|\-|tion=| 中为旁侧题注指定的对齐方式。
% \end{Note}
%
% \begin{Note}
% 除非显式用 |\caption|\-|setup[SC|\-|figure]{|\ldots|}| 或
% |\caption|\-|setup[SC|\-|table]{|\ldots|}| 设置，否则这些插图与表格
% 不使用 |margin| 或 |width| 设置。
% \end{Note}
%
% \medskip
%
% \begin{Annotation*}
% 除使用 \pkg{sidecap} 宏包外，也可使用功能更强、更灵活的
% \pkg{floatrow} 宏包\cite{floatrow} 排版旁侧题注。
% \end{Annotation*}
%
% \pagebreak[3]
% \subsubsection*{未公开功能}
% \pkg{sidecap} 宏包 \version{1.6} 带有一些未公开的宏包选项与命令，
% 可用于对旁侧题注做进一步定制：
%
% \pagebreak[3]
% \DescribeMacro{margincaption}
% 宏包选项
% \begin{quote}
%   |margincaption|\qquad{\small（例如 |\usepackage[margincaption]{sidecap}|）}
% \end{quote}
% 使所有旁侧题注放入页边距。
%
% \pagebreak[3]
% \DescribeMacro{\sidecaptionvpos}
% 命令
% \begin{quote}
%   |\sidecaptionvpos|\marg{浮动体类型}\marg{位置}
% \end{quote}
% 设置旁侧题注的垂直位置。\meta{位置} 可取 `|t|'（顶部对齐）、
% `|b|'（底部对齐）或 `|c|'（居中）。|table| 的默认取值为 `|t|'，
% |figure| 以及用 |\Declare|\-|Floating|\-|Environment| 定义的其他
% 浮动体默认取值为 `|b|'。
%
% \pagebreak[3]
% \subsection{subfigure 宏包}
% \label{subfigure}
% \packagedescription{%
% \pkg{subfigure} 宏包\cite{subfigure} 用于在单个插图或表格环境中
% 支持小型（即“子”）插图与表格。该宏包已过时，新用户应改用
% \pkg{subfig}。}
%
% 由于 \pkg{subfigure} 宏包已过时，这里仅以这样的方式支持它：在开发期
% 间使用过 \thispackage\ \version{1.x} 的旧文档仍应能按预期编译。
%
% 请改用同时支持 \thispackage\ \version{3.x} 的 \pkg{subfig} 或
% \pkg{subcaption} 宏包。
%
% \See{详见 \pkg{subfig} 宏包\cite{subfig} 文档。}
%
% \pagebreak[3]
% \subsection{supertabular 与 xtab}
% \label{supertabular}
% \packagedescription{%
% \pkg{supertabular}\cite{supertabular} 与 \pkg{xtab}\cite{xtab}
% 宏包提供可跨多页、与 \pkg{longtable} 宏包\cite{longtable} 提供的
% \texttt{long\-table} 环境相当类似的环境。关于这些功能强大的宏包
% 之间差异的详细讨论，请参阅 《\LaTeX\ Companion》\cite{TLC2}。}
%
% 若要为 |super|\-|tabular| 或 |x|\-|tabular| 环境设置专门选项，可使用
% \begin{quote}
%   |\captionsetup[supertabular]|\marg{选项}\quad\emph{或}\\
%   |\captionsetup[xtabular]|\marg{选项}\quad。
% \end{quote}
% 这些选项会在 |table| 的常规选项之外一并执行。
%
% \begin{Note}
% |\Continued|\-|Float| 对 |super|\-|tabular| 与 |x|\-|tabular| 不起作用。
% \end{Note}
%
% \pagebreak[3]
% \subsection{threeparttable 宏包}
% \label{threeparttable}
% \packagedescription{%
% \pkg{threeparttable} 宏包\cite{threeparttable} 为在表格内容与题注之后
% 设有结构化注释区的表格提供一套方案。该方案通过完全不再需要脚注，
% 解决了在表格中放置脚注这一老问题。}
%
% \NEWfeature{v3.1}
% 若要为 |three|\-|part|\-|table| 与 |measured|\-|figure| 环境设置
% 专门选项，可使用
% \begin{quote}
%   |\captionsetup[threeparttable]|\marg{选项}\quad\emph{以及}\\
%   |\captionsetup[measuredfigure]|\marg{选项}\quad。
% \end{quote}
% 这些选项会在 |table| 或 |figure| 的常规选项之外一并执行。
%
% \begin{Note}
% 由于此处题注宽度与内容相同，这些插图与表格不使用 |margin| 或
% |width| 设置；除非显式用 |\caption|\-|setup[three|\-|part|\-|table]{|\ldots|}|
% 或 |\caption|\-|setup[measured|\-|figure]{|\ldots|}| 设置。
% \end{Note}
%
% \medskip
%
% \begin{Annotation*}
% 除 \pkg{threeparttable} 宏包外，也可使用 \pkg{floatrow}
% 宏包\cite{floatrow} 实现这一目的。
% \end{Annotation*}
%
% \pagebreak[3]
% \subsection{wrapfig 宏包}
% \label{wrapfig}
% \packagedescription{%
% 与 \pkg{floatflt}、\pkg{picinpar}、\pkg{picins} 宏包类似，
% \pkg{wrapfig} 宏包\cite{wrapfig} 提供不必占满整页宽度、可被文字
% 环绕浮动的插图与表格。有关这些宏包之间差异的详细讨论，请参阅
% 《\LaTeX\ Companion》\cite{TLC2}。}
%
% \NEWfeature{v3.1}
% 若要为 |wrap|\-|figure| 与 |wrap|\-|table| 环境设置专门选项，可使用
% \begin{quote}
%   |\captionsetup[wrapfigure]|\marg{选项}\quad\emph{以及}\\
%   |\captionsetup[wraptable]|\marg{选项}\quad。
% \end{quote}
% 这些选项会在 |figure| 或 |table| 的常规选项之外一并执行。
%
% \begin{Note}
% 除非显式用 |\caption|\-|setup[wrap|\-|figure]{|\ldots|}| 或
% |\caption|\-|setup[wrap|\-|table]{|\ldots|}| 设置，否则这些插图与表格
% 不使用 |margin| 或 |width| 设置。
% \end{Note}
%
% \iffalse
% --------------------------------------------------------------------------- %
% \fi
%
% \clearpage
% \section{延伸阅读}
%
% 推荐进一步阅读以下资料：
%
% \begin{itemize}
% \item
%   \TeX\ 常见问题解答（FAQ）——关于 \TeX\ 与 \LaTeX 的常见问题：
%   \begin{quote}\url{http://faq.tug.org/}\end{quote}
%
% \item
%   一份法语 FAQ 可参见：
%   \begin{quote}\url{http://www.grappa.univ-lille3.fr/FAQ-LaTeX/}\end{quote}
%
% \item
%   Christian Faulhammer 与 Ulrich Schwarz 撰写的
%   “什么是最小可复现示例？”：
%   \begin{quote}\url{http://www.minimalbeispiel.de/mini-en.html}\end{quote}
%
% \item
%   Keith Reckdahl 撰写的 \textsf{epslatex}，包含在 \LaTeXe\ 文档中
%   插入图形的诸多技巧。该文档位于目录
%   \begin{quote}\url{http://mirror.ctan.org/info/epslatex/}\end{quote}
% \end{itemize}
%
% \iffalse
% --------------------------------------------------------------------------- %
% \fi
%
% \pagebreak[3]
% \section{致谢}
%
% 谨向 Katja Melzner、Steven D.~Cochran、Frank Mittelbach、Olga Lapko、
% David Carlisle、Carsten Heinz、Keith Reckdahl、Markus Kohm、
% Heiko Oberdiek 与 Herbert Vo\ss\ 致以谢意。衷心感谢各位提供的帮助、
% 想法、耐心、热忱与支持！
%
% 同时也感谢
% Harald Harders、
% Peter L\"offler、
% Peng Yu、
% Alexander Zimmermann、
% Matthias Pospiech、
% J\"urgen Wieferink、
% Christoph Bartoschek、
% Uwe St\"ohr、
% Ralf Stubner、
% Geoff Vallis、
% Florian Keiler、
% J\"urgen G\"obel、
% Uwe Siart、
% Sang-Heon Shim、
% Henrik Lundell、
% David Byers、
% William Asquith、
% Dirk Hoffmann 教授、
% Frank Martini、
% Danie Els、
% Philipp Woock、
% Fadi Semmo、
% Matthias Stevens 与
% Leo Liu，
% 他们都为本宏包的改进作出了贡献。
%
% \iffalse
% --------------------------------------------------------------------------- %
% \fi
%
% \clearpage\appendix
% \section{按字母顺序的参考索引}
%
% \makeatletter
% \newenvironment{Warnings}{%
%   \newif\ifMessage
%   \def\Message{%
%     \ifMessage\end{minipage}\par\vspace\itemsep\pagebreak[3]\fi
%     \begin{minipage}{\linewidth}%
%     \setlength\parindent{0pt}%
%     \setlength\parskip{\smallskipamount}%
%     \setlength\leftskip{0pt}%
%     \color{blue}%
%     \Messagetrue}%
%   \def\Description{%
%     \par\color{black}\leftskip=1.5em}%
%   \def\Or{\@testopt\@Or{或}}%
%   \def\@Or[##1]{%
%     \\{\small\color{black}\hspace{1.5em}\textit{-- ##1 --}}\\}%
% }{%
%   \ifMessage\end{minipage}\par\fi
% }
% \makeatother
%
% \subsection{选项}
% \newcommand\preambleoption{~%
%   \begingroup\def\thefootnote{\fnsymbol{footnote}}%
%     \footnotemark[2]%
%   \endgroup}
% \newcommand\preambleoptiontext{%
%   \begingroup\def\thefootnote{\fnsymbol{footnote}}%
%     \footnotetext[2]{这些选项仅可在文档导言区使用。}%
%   \endgroup}
%
% \begin{longtable*}{lll}
% 选项              & 简要说明 & 小节 \\
% \hline
% \endhead
% |aboveskip|       & 设置题注上方间距 & \ref{skips} \\
% |belowskip|       & 设置题注下方间距 & \ref{skips} \\
% |compatibility|\preambleoption & 强制（非）兼容模式 & \ref{compatibility} \\
% |figureposition|\preambleoption & 提示插图题注的位置 & \ref{skips} \\
% |font|(|+|)       & 设置字体 & \ref{fonts} \\
% |format|          & 设置格式 & \ref{formats} \\
% \iffalse
% |FPlist|          & FPfigure 的目录条目应链接到何处？ & \ref{fltpage} \\
% |FPref|           & 指向 FPfigure 的 |\ref| 应链接到何处？ & \ref{fltpage} \\
% \fi
% |hangindent|      & 设置悬挂缩进 & \ref{margins} \\
% |hypcap|          & 选择“hypcap”功能 & \ref{hyperref} \\
% |hypcapspace|     & 设置超链接与内容之间的距离 & \ref{hyperref} \\
% |indention|       & 设置缩进 & \ref{margins} \\
% |justification|   & 设置对齐方式 & \ref{justification} \\
% |labelfont|(|+|)  & 设置题注标签字体 & \ref{fonts} \\
% |labelformat|     & 设置题注标签格式 & \ref{formats} \\
% |labelsep|        & 设置标签分隔符 & \ref{formats} \\
% |labelseparator|  & ——同 |labelsep|—— & \ref{formats} \\
% |list|            & 开启或关闭目录条目 & \ref{lists} \\
% |listformat|      & 设置“插图/表格目录”条目格式 & \ref{lists} \\
% |margin|          & 设置边距 & \ref{margins} \\
% |margin*|         & 设置边距，但仅在未设置宽度时生效 & \ref{margins} \\
% |maxmargin|       & 设置最大边距 & \ref{margins} \\
% |minmargin|       & 设置最小边距 & \ref{margins} \\
% |name|            & 设置当前环境的名称 & \ref{names} \\
% |oneside|         & 选择单面模式 & \ref{margins} \\
% |options|         & 执行给定选项列表 & \\
% |parindent|       & 设置段落缩进 & \ref{margins} \\
% |parskip|         & 设置段落间距 & \ref{margins} \\
% |position|        & 提示题注位置 & \ref{skips} \\
% |singlelinecheck| & 开启或关闭单行检查 & \ref{justification} \\
% |skip|            & 设置内容与题注之间的间距 & \ref{skips} \\
% |strut|           & 开启或关闭 |\strut| 的使用 & \ref{formats} \\
% |style|           & 设置题注样式 & \ref{style} \\
% |subtype|         & 设置子题注类型 & \phantom{t}--~\footnotemark \\
% |tableposition|\preambleoption & 提示表格题注的位置 & \ref{skips} \\
% |textfont|(|+|)   & 设置题注文字字体 & \ref{fonts} \\
% |textformat|      & 设置题注文字格式 & \ref{formats} \\
% |twoside|         & 选择双面模式 & \ref{margins} \\
% |type|            & 设置题注类型并放置超链接锚点 & \ref{types} \\
% |type*|           & 仅设置题注类型 & \ref{types} \\
% |width|           & 设置固定的题注宽度 & \ref{margins} \\
% \end{longtable*}
% \preambleoptiontext
% \footnotetext{选项 \texttt{subtype} 见 \pkg{subcaption} 宏包文档。}
%
% \nopagebreak\parbox[t]{\linewidth}{% 防止分页
% \begin{Note*}
%   此处未列出已过时的选项。这些选项的清单见
%   \longref{caption1} 与
%   \longref{caption2}。
% \end{Note*}}
%
% \pagebreak[3]
% \subsection{命令}
%
% \begin{longtable*}{lll}
% 命令                          & & 小节 \\
% \hline
% \endhead
% |\abovecaptionskip|              & & \ref{skips} \\
% |\belowcaptionskip|              & & \ref{skips} \\
% |\caption|                       & & \ref{caption} \\
% |\caption*|                      & & \ref{caption} \\
% |\captionlistentry|              & & \ref{captionlistentry} \\
% |\captionof|                     & & \ref{caption} \\
% |\captionof*|                    & & \ref{caption} \\
% |\captionsetup|                  & & \ref{captionsetup} \\
% |\captionsetup*|                 & & \ref{captionsetup} \\
% |\centerfirst|                   & & \ref{justification} \\
% |\centerlast|                    & & \ref{justification} \\
% |\clearcaptionsetup|             & & \ref{captionsetup} \\
% |\clearcaptionsetup*|            & & \ref{captionsetup} \\
% |\ContinuedFloat|                & & \ref{ContinuedFloat} \\
% |\DeclareCaptionFont|            & & \ref{declare} \\
% |\DeclareCaptionFormat|          & & \ref{declare} \\
% |\DeclareCaptionFormat*|         & & \ref{declare} \\
% |\DeclareCaptionJustification|   & & \ref{declare} \\
% |\DeclareCaptionLabelFormat|     & & \ref{declare} \\
% |\DeclareCaptionLabelSeparator|  & & \ref{declare} \\
% |\DeclareCaptionLabelSeparator*| & & \ref{declare} \\
% |\DeclareCaptionListFormat|      & & \ref{declare} \\
% |\DeclareCaptionOption|          & & \ref{declare} \\
% |\DeclareCaptionStyle|           & & \ref{declare} \\
% |\DeclareCaptionSubType|         & & --~\footnotemark \\
% |\DeclareCaptionTextFormat|      & & \ref{declare} \\
% |\showcaptionsetup|              & & \ref{captionsetup} \\
% \end{longtable*}
% \footnotetext{\cs{DeclareCaptionSubType} 见
%   \pkg{subcaption} 宏包文档。}
%
% \clearpage
% \subsection{警告}
%
% \begin{Warnings}
%
% \Message
%   |\caption outside box or environment.|
% \Or
%   |\captionsetup{type=...} outside box or environment.|
% \Or
%   |\captionsetup{type*=...} or \captionof outside box|\\
%   | or environment.|
% \Description
%   你把 |\caption|、|\caption|\-|of| 或
%   |\caption|\-|setup{type=|\x\meta{类型}|}| 命令放在了某个盒子、组或环境
%   之外。不应这样做，因为可能引发不良副作用。
%   \par\See{\longref{types} 与 \longref{caption}}
%
% \Message
%   |\caption will not be redefined since it's already|\\
%   |redefined by a document class or package which is|\\
%   |unknown to the caption package.|
% \Description
%   若 \thispackage 检测到某个（未知的）文档类或宏包已对 |\caption| 做了
%   增强，它将不再重定义 |\caption|，因为那样会直接破坏该增强。结果是，
%   部分功能（如 |\caption*|、|\Continued|\-|Float|、|\caption|\-|setup|
%   可选参数的使用，以及选项 |list=| 与 |hypcap=|）将不可用。\par
%   若不在意原有增强，而希望改用 \thispackage 的全套功能，可尝试使用
%   不受支持的（!）选项 |compatibility=|\x|false|，并祝好运。
%   （但你将转而收到下一条警告。）
%   \par\See{\longref{classes} 与 \longref{compatibility}}
%
% \Message
%   |Forced redefinition of \caption since the|\\
%   |unsupported(!) package option `compatibility=false'|\\
%   |was given.|
% \Description
%   既然你执意指定了选项 |compatibility=|\x|false|，\thispackage 会尽力
%   满足你的要求。但根据所使用的文档类或其他宏包，这可能导致功能无法
%   使用甚至报错。请做好心理准备！
%   \par\See{\longref{compatibility}}
%
% \Message
%   |Hyperref support is turned off because hyperref has|\\
%   |stopped early.|
% \Description
%   若 \pkg{hyperref} 宏包在加载过程中提前中断（无论何种原因），
%   \thispackage 的 \pkg{hyperref} 支持将不可用。结果可能出现
%   \pkg{hyperref} 警告，以及指向插图或表格的超链接失效。
%   \par\See{\longref{hyperref}\,}
%
% \Message
%   |Ignoring optional argument [|\meta{位置}|] of \setcapwidth.|
% \Description
%   \Thispackage 尽力模拟 \KOMAScript 中与题注相关的命令。但使用本宏包时，
%   \KOMAScript 命令 |\setcapwidth| 的可选参数（尚）不能生效；若仍尝试
%   使用，将收到本警告。
%   \par\See{\longref{KOMA}}
%
% \Message
%   |Internal Warning: |\meta{警告消息}|.|
% \Description
%   你本不应看到这条警告：要么你使用了某个重定义 |figure| 或 |table|、
%   且 \thispackage 不认识的宏包，要么这是 \thispackage 的缺陷。
%   请发邮件向作者报告此问题。
%
% \Message
%   |\label without proper \caption|
% \Description
%   就 |\label| 而言，浮动环境的行为与其非浮动对应物不同：内部引用不会
%   在环境开头生成，而是在 |\caption| 处生成。因此 |\label| 命令必须放在
%   题注文字（|\caption| 的必填参数）的正\emph{后方}或\emph{内部}。
%
% \Message
%   |Option `|\meta{选项}|' was not in list `|\meta{选项列表}|'.|
% \Description
%   若尝试从某选项列表中移除特定选项（例如用
%   |\clear|\-|caption|\-|setup[po|\-|si|\-|tion]{table}|），而该选项在
%   列表中找不到，将收到本警告。若并非笔误且希望关闭此警告，请改用
%   |\clear|\-|caption|\-|setup*| 代替 |\clear|\-|caption|\-|setup|。
%   \par\See{\longref{captionsetup}}
%
% \Message
%   |Option list `|\meta{选项列表}|' undefined.|
% \Description
%   若尝试从某选项列表中移除特定选项（例如用
%   |\clear|\-|caption|\-|setup[for|\-|mat]{figure}|），而该选项列表
%   （尚未）定义，将收到本警告。若并非笔误且希望关闭此警告，请改用
%   |\clear|\-|caption|\-|setup*| 代替 |\clear|\-|caption|\-|setup|。
%   \par\See{\longref{captionsetup}}
%
% \Message
%   |Obsolete option `ignoreLTcapwidth' ignored.|
% \Description
%   本版本 \thispackage 不再模拟 \pkg{caption2} 宏包选项
%   |ignoreLTcapwidth|；通常直接删除该选项即可。
%   \par\See{\longref{caption2} 与 \longref{longtable}}
%
% \Message
%   |`ragged2e' support has been changed. Rerun to get|\\
%   |captions right.|
% \Description
%   \pkg{ragged2e} 宏包仅在确实需要时才由 \thispackage 加载。这至少需要
%   两次 \LaTeX 编译，因此首次编译时可能收到本警告。再次编译文档即可，
%   该警告应会消失。
%   \par\See{\longref{justification}}
%
% \Message
%   |Reference on page |\meta{页码}| undefined.|
% \Description
%   若采用 |twoside| 页面布局，\thispackage 需要跟踪页码以正确算出边距。
%   这至少需要两次 \LaTeX 编译，首次编译时可能收到本警告。再次编译
%   文档即可，该警告应会消失。
%   \par\See{\longref{margins}}
%
% \Message
%   |The caption type was already set to `|\meta{类型}|'.|
% \Description
%   本警告提示题注选项混用。例如在 |figure| 环境内部使用
%   |\caption|\-|setup{type=|\x|table}| 或
%   |\caption|\-|of{table}{|\ldots|}|，将导致该题注同时使用两套选项：
%   |figure| 的一套（由 |\caption|\-|setup[figure]{|\ldots|}| 指定）与
%   |table| 的一套（由 |\caption|\-|setup[table]{|\ldots|}| 指定）。\par
%   {\small（可用带星号形式 |\caption|\-|setup*{type=|\x\ldots|}| 关闭本警告。）\par}
%   \par\See{\longref{captionsetup}}
%
% \Message
%   |The option `hypcap=true' will be ignored for this|\\
%   |particular \caption.|
% \Description
%   \Thispackage 未为该特定题注找到合适的超链接锚点，因此决定忽略
%   |hypcap=|\x|true|（默认开启）设置。结果是指向该题注的链接（例如
%   插图目录中的条目，或由 |\ref|、|\auto|\-|ref| 生成）将链接到插图或
%   表格的题注，而非插图或表格本身。\par
%   这可能发生在非浮动环境内使用 |\caption|\-|of| 时，也可能发生在使用
%   某个重定义 |figure| 或 |table|、且 \thispackage 不认识的宏包时。\par
%   若可以接受但希望关闭本警告，只需在引发警告的 |\caption| 或
%   |\caption|\-|of| 命令之前放置
%   |\caption|\-|setup{hyp|\-|cap=|\x|false}|。
%   若不能接受，可自行用
%   |\caption|\-|setup{type=|\x\meta{浮动体类型}|}| 设置超链接锚点。
%   \par\See{\longref{hyperref}\,}
%
% \Message
%   |Unsupported document class (or package) detected,|\\
%   |usage of the caption package is not recommended.|
% \Description
%   要么你所用的文档类是 \thispackage 不认识的，要么文档中包含了同样
%   重定义 |\@make|\-|caption|（内部负责排版题注）的宏包。这意味着
%   \thispackage 要么会以不希望的方式改变题注设计，甚至可能完全无法
%   正常工作。
%   \par\See{\longref{classes} 与 \longref{compatibility}}
%
% \Message
%   |Unused \captionsetup[|\meta{类型}|].|
% \Description
%   你通过 |\caption|\-|setup[|\meta{类型}|]| 指定了选项，但之后并未使用。
%   原因可能是 \meta{类型} 拼写有误；使用了某个重定义 |figure| 或
%   |table|、且 \thispackage 不认识的宏包；或仅仅是此行之后根本没有
%   使用 \meta{类型} 环境。（若希望关闭本警告，可改用 |\caption|\-|setup*|
%   代替 |\caption|\-|setup|。）
%   \par\See{\longref{captionsetup}}
%
% \Message
%   |Usage of the |\meta{宏包}| package together with the|\\
%   |caption package is strongly not recommended.|\\
%   |Instead of loading the |\meta{宏包}| package you should|\\
%   |use the caption package option `tableposition=top'.|
% \Description
%   \meta{宏包} 同样处理题注间距。请确定究竟想用哪一个：\thispackage
%   还是另一个；两者并用可能导致题注上下间距错误。
%   \par\See{\longref{skips}}
%
% \end{Warnings}
%
% \pagebreak[3]
% \subsection{错误}
%
% \begin{Warnings}
%
% \Message
%   |Argument of \@caption has an extra }.|
% \Or
%   |Paragraph ended before \@caption was complete.|
% \Description
%   若要把特殊内容（例如一个 tabular）作为题注排版，即使不使用插图目录
%   或表格目录，也必须为 |\caption| 或 |\caption|\-|of| 提供目录条目用的
%   可选参数。
%   \par\See{\longref{caption}、\longref{lists} 与 \longref{hyperref}\,}
%
% \Message
%   |\caption outside float.|
% \Description
%   你把 |\caption| 命令放在了浮动环境，或 |long|\-|table|、|wrap|\-|figure|
%   之外。若确实如此，请改用 |\caption|\-|setup{type=|\x\meta{类型}|}|
%   $+$ |\caption|，或使用 |\caption|\-|of|。
%   \par\See{\longref{caption}}
%
% \Message
%   |\ContinuedFloat outside float.|
% \Description
%   你把 |\Continued|\-|Float| 命令放在了浮动环境之外。% 或 |long|\-|table|。
%   若确实如此，请使用组合
%   |\caption|\-|setup{type=|\x\meta{类型}|}| $+$ |\Continued|\-|Float|。\par
%   请注意，|\Continued|\-|Float| 在 |long|\-|table| 内部不起作用；
%   但若足以满足需要，可使用 |long|\-|table*| 环境排版不使 |table|
%   计数器加一的 |long|\-|table|。
%   \par\See{\longref{ContinuedFloat} 与 \longref{longtable}}
%
% \Message
%   |Continued `|\meta{类型}|' after `|\meta{类型}|'.|
% \Description
%   续接的插图或表格不允许被另一类型的浮动环境（或 |long|\-|table|）
%   打断，例如插在一幅插图与一幅续接插图之间的表格。
%   \par\See{\longref{ContinuedFloat}}
%
% \Message
%   |For a successful cooperation we need at least version|\\
%   |`|\meta{日期}|' of package |\meta{宏包}|, but only version|\\
%   |`|\meta{旧日期}|' is available.|
% \Description
%   \Thispackage 无法与如此陈旧的宏包配合工作；请将其更新到较新版本，
%   至少更新到所要求的版本。
%
% \Message
%   |Internal Error: |\meta{错误消息}|.|
% \Description
%   你本不应看到这条错误。请发邮件向作者报告此问题。
%
% \Message
%   |No float type '|\meta{类型}|' defined.|
% \Description
%   你在 |\caption|\-|setup{type=|\x\meta{类型}|}|、|\caption|\-|of|\marg{类型}
%   或 |\Declare|\-|Caption|\-|Sub|\-|Type| 中指定的 \meta{类型} 尚未定义。
%   \meta{类型} 应为 `\texttt{figure}' 或 `\texttt{table}'，或用
%   \pkg{newfloat} 宏包提供的 |\Declare|\-|Floating|\-|Environment|、
%   \pkg{float} 宏包\cite{float} 提供的 |\new|\-|float|，或
%   \pkg{floatrow} 宏包\cite{floatrow} 提供的
%   |\Declare|\-|New|\-|Float|\-|Type| 定义的其他浮动环境。
%
% \Message
%   |Not allowed in longtable* environment.|
% \Description
%   |long|\-|table*| 环境内部不允许使用 |\caption|。请改用不带标签的
%   |\caption*|，或改用普通的 |long|\-|table| 环境。
%
% \Message
%   |Not available in compatibility mode.|
% \Description
%   所需功能在兼容模式下不受支持。“兼容模式”是指 \thispackage 检测到
%   不兼容的文档类，或同样扩展 |\caption| 命令的不兼容宏包；在保留原有
%   扩展的前提下，\thispackage 的部分功能不受支持。
%   \par\See{\longref{compatibility}}
%
% \Message
%   |Only one \caption can be placed in this environment.|
% \Description
%   \pkg{fltpage} 与 \pkg{sidecap} 宏包提供的环境内部只能放置
%   \emph{一则}题注。（这是实现设计使然。）
%
% \Message
%   |Paragraph ended before \caption@makecurrent was complete.|
% \Or
%   |Paragraph ended before \caption@prepareanchor was complete.|
% \Description
%   若要排版多段落题注，即使不使用插图目录或表格目录，也必须为
%   |\caption| 或 |\caption|\-|of| 提供目录条目用的可选参数。
%   \par\See{\longref{caption}、\longref{lists} 与 \longref{hyperref}\,}
%
% \Message
%   |Something's wrong--perhaps a missing \caption|\\
%   |in the last figure or table.|
% \Description
%   似乎你使用了 |\sub|\-|caption| 命令（或另一则排版了子题注的命令），
%   却没有相应的 |\caption| 命令。这不受支持。
%
% \Message
%   |The option `labelsep=|\meta{名称}|' does not work|\\
%   |with `format=hang'.|
% \Or
%   |The option `labelsep=|\meta{名称}|' does not work|\\
%   |with \setcaphanging (which is set by default).|
% \Description
%   含 |\\| 命令的题注标签分隔符（如 |labelsep=|\x|newline|）不能与
%   悬挂题注格式（如 |format=|\x|hang|）搭配使用。请另选题注标签分隔符
%   （例如 |labelsep=|\x|colon|）或另选题注格式（例如
%   |format=|\x|plain|）。
%   \par\See{\longref{formats} 或 \longref{KOMA}}
%
% \Message
%   |The package option `caption=false' is obsolete.|\\
%   |Please pass this option to the subfig package instead|\\
%   |and do *not* load the caption package anymore.|
% \Description
%   你指定了选项 |caption=|\x|false|。这曾是一种变通办法：不加载整套
%   \pkg{caption} 宏包（保留文档类或其他宏包提供的题注功能），同时让
%   \pkg{subfig} 宏包继续工作。该机制已过时、不再提供；请改为把该选项
%   传给 \pkg{subfig} 宏包，并且不再加载 \thispackage。
%   \par\See{\pkg{subfig} 宏包\cite{subfig} 文档}
%
% \Message
%   |Undefined boolean value `|\meta{取值}|'.|
% \Description
%   你尝试为布尔选项（如 |singlelinecheck=| 或 |hypcap=|）设置了不合法的取值。
%   此处仅允许 |false|、|no|、|off|、|0|，或 |true|、|yes|、|on|、|1|。
%
% \Message
%   |Undefined format `|\meta{名称}|'.|
% \Description
%   你尝试设置一个并不存在的题注格式。也许是笔误！？
%   \par\See{\longref{formats}}
%
% \Message
%   |Undefined label format `|\meta{名称}|'.|
% \Description
%   你尝试设置一个并不存在的题注标签格式。也许是笔误！？
%   \par\See{\longref{formats}}
%
% \Message
%   |Undefined label separator `|\meta{名称}|'.|
% \Description
%   你尝试设置一个并不存在的题注标签分隔符。也许是笔误！？
%   \par\See{\longref{formats}}
%
% \Message
%   |Undefined list format `|\meta{名称}|'.|
% \Description
%   你尝试设置一个并不存在的目录条目格式。也许是笔误！？
%   \par\See{\longref{lists}}
%
% \Message
%   |Undefined position `|\meta{名称}|'.|
% \Description
%   你尝试为题注位置设置了不合法的取值。也许是笔误！？
%   \par\See{\longref{skips}}
%
% \Message
%   |Undefined style `|\meta{名称}|'.|
% \Description
%   你尝试设置一个并不存在的题注样式。也许是笔误！？
%   \par\See{\longref{style}}
%
% \Message
%   |Usage of the `position' option is incompatible|\\
%   |to the `|\meta{宏包}|' package.|
% \Description
%   所列宏包同样处理题注间距。请确定究竟想用哪一个：\thispackage 的
%   |position=| 选项，还是另一套机制；两者并用会导致题注上下间距错误，
%   因而不受支持。
%   \par\See{\longref{skips}}
%
% \Message
%   |Undefined text format `|\meta{名称}|'.|
% \Description
%   你尝试设置一个并不存在的题注文字格式。也许是笔误！？
%   \par\See{\longref{formats}}
%
% \Message
%   |You can't use both, the (obsolete) caption2 *and*|\\
%   |the (current) caption package.|
% \Description
%   本错误消息已说明一切：不能两者并用。请只使用 \thispackage。
%   \par\See{\longref{caption2}}
%
% \end{Warnings}
%
% \iffalse
% --------------------------------------------------------------------------- %
% \fi
%
% \clearpage
% \section{版本历史}
% \label{history}
%
% 版本 $1.0$ 完成于 1994 年，提供了若干用于定制题注外观的选项。
% 此外该版本已支持 \pkg{rotating} 与 \pkg{subfigure} 宏包。
% 版本 $1.1$ 引入了 |center|\-|last| 选项；版本 $1.2$ 增加了对
% \pkg{float} 宏包的支持；版本 $1.3$ 改进了对 \pkg{subfigure} 宏包的
% 支持；版本 $1.4$ 则带来了新选项 |nooneline|。
%
% \pkg{caption2} 宏包 $2.0$ 是常规 \pkg{caption} 宏包的实验性分支版本。
% 当时对新功能以及对 \pkg{longtable} 等额外宏包适配的需求十分强烈，
% 因此它在 1995 年作为 beta 测试版公开发布，并附以尚不完整的文档。
% 此外它还改进了对 \pkg{subfigure} 宏包的支持。（改进版 $2.1$ 于
% 2002 年提供。）
%
% 2003 年作者终于抽出时间，与 Frank Mittelbach 和 Steven Cochran 合作
% 完成了 \pkg{caption} 宏包新的正式发行版 $3.0$。该版本于 2003 年 12 月
% 发布，取代了已被搁置的 \pkg{caption2} 宏包。主体部分经过重写，用户
% 界面得到彻底重构；此外还额外支持了 \pkg{hyperref}、\pkg{hypcap}、
% \pkg{listings}、\pkg{sidecap} 与 \pkg{supertabular} 宏包。
%
% 此前各版本均面向标准 \LaTeX\ 文档类 \cls{article}、\cls{report} 与
% \cls{book} 设计；2007 年发布的当前版本 $3.1$ 还支持 \AmS、
% \KOMAScript、\NTG、\SmF\ 文档类以及 \cls{beamer} 类。此外它还适配了
% \pkg{french} Babel 选项、\pkg{frenchle} 与 \pkg{frenchpro} 宏包，
% 以及 \pkg{floatflt}、\pkg{fltpage}、\pkg{picinpar}、\pkg{picins}、
% \pkg{setspace}、\pkg{threeparttable} 与 \pkg{wrapfig} 宏包。
% 同时引入了新的选项与命令，包括 |font+|、|figure|\-|within| 与
% |table|\-|within|、|list| 与 |list|\-|format|、|max|\-|margin| 与
% |min|\-|margin|、|\caption|\-|list|\-|entry|、
% |\Declare|\-|Caption|\-|List|\-|Format| 等。其他益处包括新的兼容性
% 检查 \see*{\longref{compatibility}}、新的“hypcap”功能
% \see*{\longref{hyperref}\,}，以及子题注功能
% \see*{\pkg{subcaption} 宏包文档}。
%
% \iffalse
% --------------------------------------------------------------------------- %
% \fi
%
% \clearpage
% \section{与旧版本的兼容性}
%
% \subsection{caption v\texorpdfstring{$1.x$}{1.x}}
% \label{caption1}
%
% 本版本 \thispackage 仍支持本宏包 $1.0$ 至 $1.4$ 版本所提供的旧选项与
% 命令，因此编译旧文档时应不会出现问题；但请勿将旧选项、旧命令与新的
% 混用。这种用法不受支持，可能引发不良副作用。
%
% 下面简要列出已过时的选项，以及它们在本版本 \thispackage 中的替代方式：
%
% \begin{center}\small
% \begin{tabular}{ll}
% \pkg{caption} \version{1.x} & \pkg{caption} \version{3.x}\\
% \hline
% |normal|        & |format=plain|\\
% |hang|          & |format=hang|\\
% |isu|           & |format=hang|\\
% |center|        & |justification=centering|\\
% |centerlast|    & |justification=centerlast|\\
% |nooneline|     & |singlelinecheck=off|\\
% |scriptsize|    & |font=scriptsize|\\
% |footnotesize|  & |font=footnotesize|\\
% |small|         & |font=small|\\
% |normalsize|    & |font=normalsize|\\
% |large|         & |font=large|\\
% |Large|         & |font=Large|\\
% |up|            & |labelfont=up|\\
% |it|            & |labelfont=it|\\
% |sl|            & |labelfont=sl|\\
% |sc|            & |labelfont=sc|\\
% |md|            & |labelfont=md|\\
% |bf|            & |labelfont=bf|\\
% |rm|            & |labelfont=rm|\\
% |sf|            & |labelfont=sf|\\
% |tt|            & |labelfont=tt|\\
% \end{tabular}
% \end{center}
%
% 除用于设置字体的选项外，还可以用 |\re|\-|new|\-|command| 重定义
% |\caption|\-|size|、|\caption|\-|font| 与 |\caption|\-|label|\-|font|
% 命令，以此作为另一种更灵活的方式改变题注字体。该机制已由下列命令
% 取代：
% \begin{quote}
%   |\DeclareCaptionFont{|\ldots|}{|\ldots|}|\qquad 以及\\
%   |\captionsetup{font=|\ldots|,labelfont=|\ldots|}|\qquad。
% \end{quote}
% \SeeUserDefined
%
% 在 \version{1.x} 中，题注边距通过
% \begin{quote}
%   |\setlength{\captionmargin}{|\ldots|}|\quad  设置。
% \end{quote}
% 现已由下列写法取代：
% \begin{quote}
%   |\captionsetup{margin=|\ldots|}|\qquad。
% \end{quote}
% \See{\longref{margins}}
%
% 例如旧式代码
% \begin{quote}
%   |\usepackage[hang,bf]{caption}|\\
%   |\renewcommand\captionfont{\small\sffamily}|\\
%   |\setlength\captionmargin{10pt}|
% \end{quote}
% 仍可正常工作，但如今应写为
% \begin{quote}
%   |\usepackage[format=hang,labelfont=bf,font={small,sf},|\\
%   |            margin=10pt]{caption}|
% \end{quote}
% 或
% \begin{quote}
%   |\usepackage{caption}|\\
%   |\captionsetup{format=hang,labelfont=bf,font={small,sf},|\\
%   |              margin=10pt}|\qquad。
% \end{quote}
%
% 较为少见的选项 |ruled| 允许对用 \pkg{float} 宏包定义的 |ruled|
% 浮动体部分套用题注设置；本版本 caption 宏包同样会模拟该选项。但不再
% 推荐使用它，因为本版本 \thispackage 提供了更灵活的方式来更改这些
% 浮动环境的题注：
% \begin{quote}
%   |\DeclareCaptionStyle{ruled}{|\ldots|}|
% \end{quote}
% 或
% \begin{quote}
%   |\captionsetup[ruled]{|\ldots|}|\qquad。
% \end{quote}
% \SeeUserDefined[, \longref{captionsetup}, 与 \longref{float}]
%
% \subsection{caption2 v\texorpdfstring{$2.x$}{2.x}}
% \label{caption2}
%
% 尽管 \pkg{caption} 宏包与其实验性、现已过时的变体 \pkg{caption2}
% 所做的事情非常相近，但二者的实现设计差异很大，因此无法提供完全兼容。
% 为此，本宏包发行版中仍保留名为 |caption2.sty| 的文件，以便使用
% \pkg{caption2} 的旧文档仍能顺利编译。
%
% 新建文档应改用 \thispackage 的当前版本。多数情况下，只需把命令
% \begin{quote}
%   |\usepackage[...]{caption2}|
% \end{quote}
% 替换为
% \begin{quote}
%   |\usepackage[...]{caption}|\qquad。
% \end{quote}
% 但部分选项与命令不会被模拟，之后可能出现错误消息。希望本节能帮助你
% 消除这些错误。若从 \pkg{caption2} 迁移到 \pkg{caption} 时遇到问题，
% 请随时发邮件求助。
%
% 除上一节列出的过时选项外，下列选项同样会被模拟：
%
% \begin{center}\small
% \begin{tabular}{ll}
% \pkg{caption2} \version{2.x} & \pkg{caption} \version{3.x}\\
% \hline
% |flushleft|   & |justification=raggedright|\\
% |flushright|  & |justification=raggedleft|\\
% |oneline|     & |singlelinecheck=on|\\
% \end{tabular}
% \end{center}
%
% 在 \version{2.x} 中，题注边距通过
% \begin{quote}\leavevmode\hbox{%
%   |\setcaptionmargin{|\ldots|}| 或
%   |\setcaptionwidth{|\ldots|}|\quad。
% }\end{quote}
% 设置。现已由下列写法取代：
% \begin{quote}\leavevmode\hbox{%
%   |\captionsetup{margin=|\ldots|}| 或
%   |\captionsetup{width=|\ldots|}|\quad。
% }\end{quote}
% \See{\longref{margins}}
%
% 在 \version{2.x} 中，缩进通过
% \begin{quote}
%   |\captionstyle{indent}|\\
%   |\setlength\captionindent{|\ldots|}|\quad。
% \end{quote}
% \pagebreak[3]
% 设置。现已改为
% \nopagebreak[3]
% \begin{quote}
%   |\captionsetup{format=plain,indention=|\ldots|}|\quad。
% \end{quote}
%
% 所谓单行检查，由命令
% |\oneline|\-|captions|\-|false|（关闭检查）与
% |\oneline|\-|captions|\-|true|（开启检查）控制。现已由
% |\caption|\-|setup{single|\-|line|\-|check=|\x|off}|
% 或
% |\caption|\-|setup{single|\-|line|\-|check=|\x|on}|\quad 取代。
% \See{\longref{justification}}
%
% 下列命令
% \begin{quote}
%   |\captionlabeldelim|、|\captionlabelsep|、|\captionstyle|、\\
%   |\defcaptionstyle|、|\newcaptionstyle| 与 |\renewcaptionstyle|
% \end{quote}
% 没有简单的替代方式，因此本版本 \thispackage 不会模拟它们。（使用这些
% 命令会导致错误消息。）改写这类代码并不总是轻而易举；但认真阅读本
% 手册，你应能找到合适的选项与命令作为替代。
%
% \iffalse
% ...（若干示例）
% \fi
%
% \version{2.x} 的选项 |ignore|\-|LT|\-|cap|\-|width| 同样没有替代方式。
% 但多数情况下可直接弃用该选项，因为在本版本 \thispackage 中，
% |\LT|\-|cap|\-|width| 的值本来就会被忽略（除非把它设置成不同于
% 默认值 |4in| 的值）。
% \See{\longref{longtable}}
%
% \subsection{caption v\texorpdfstring{$3.0$}{3.0}}
% \label{caption3}
%
% \thispackage\ \version{3.0} 除标准 \LaTeX\ 文档类
% \cls{article}、\cls{report} 与 \cls{book} 外，不支持其他文档类。
% 因此当时默认设置是固定的；如今 \thispackage 支持更多文档类，默认设置
% 会随所用文档类而定。
%
% 例如：在 \version{3.0} 中，题注默认 |jus|\-|ti|\-|fi|\-|ca|\-|tion|
% 始终为 |jus|\-|ti|\-|fied|；如今使用标准文档类时仍是
% |jus|\-|ti|\-|fied|，但搭配 \cls{beamer} 文档类时则以
% |ragged|\-|right| 作为默认格式。
%
% 选择“旧”默认值的简便方式是：加载 \thispackage 时使用选项
% |style=base|（或之后通过 |\caption|\-|setup| 设置）。
%
% \medskip
%
% 另一个新功能是自动兼容性检查：发现不兼容时会发出警告消息；若发现严重
% 不兼容，\thispackage 的部分功能将被停用。请注意，新增到 \thispackage
% 的只是这一检查；因此收到此类警告时，旧版本 \thispackage 同样存在
% 不兼容，只是没有发出警告，“仅仅”表现为副作用或引发问题。所以这些
% 警告只表示可能出现某些不良副作用或问题，并非表示一定会发生。无论
% 如何，在此类情形下使用 \thispackage 都应谨慎。
%
% \bigskip
%
% \begin{Note*}
% \thispackage\ \version{3.0} 曾提供选项 |caption=|\x|false|，作为一种
% 变通办法：不加载整套 \pkg{caption} 宏包（保留文档类或其他宏包提供的
% 题注功能），同时让 \pkg{subfig} 宏包\cite{subfig} 继续工作。该机制
% 已过时、不再提供；请改为把该选项传给 \pkg{subfig} 宏包，并且不再
% 加载 \thispackage。
% \end{Note*}
%
% \iffalse
% --------------------------------------------------------------------------- %
% \fi
%
% \clearpage
% \section{供文档类作者使用的命令}
%
% 文档类可以直接加载 \pkg{caption} 宏包，而不必自行定义
% |\@make|\-|caption| 的变体。但若作者出于任何原因决定不这样做，
% 还有第二种方案：
%
% \NEWfeature{v3.5}
% 若文档类定义了 |\caption@document|\-|class|，则在加载 \pkg{caption}
% 内核时，会以 |caption-|\meta{文档类}|.sto| 的形式加载对 \pkg{caption}
% 宏包的适配文件。该适配文件的目标是更改 \pkg{caption} 宏包的默认设置，
% 使直接加载它（不带任何宏包选项）不会损害题注的外观。这样终端用户便可
% 通过添加宏包选项，对题注外观的某些方面进行微调，而不会对其他方面造成
% 不希望的副作用。
%
% 两种情况下都可使用命令 |\Set|\-|Caption|\-|Default| 调整
% \pkg{caption} 宏包的默认值：
%
% \DescribeMacro{\SetCaptionDefault}
% 命令
% \begin{quote}
%   |\SetCaptionDefault*|\marg{选项}\marg{默认值}
% \end{quote}
% 用于更改选项 \meta{选项} 的默认值，例如：
% \begin{quote}
%   |\SetCaptionDefault{format}{hang}|\\
%   |\SetCaptionDefault{justification}{raggedright}|\\
%   |\DeclareCaptionLabelSeparator{silly}{ +++ }|\\
%   |\SetCaptionDefault{labelseparator}{silly}|
% \end{quote}
% 本命令可用于下列选项：
%   |box|、|font|、|format|、|justification|、|label|\-|font|、
%   |label|\-|format|、|label|\-|separator|、|list|\-|format|、
%   |margin|、|parbox|、|position|、|single|\-|line|\-|check|、
%   |text|\-|font|、|text|\-|format|，以及 |width|。
%
% 带星号形式仅在设置 |margin| 或 |width| 的默认长度时才有区别：此时
% 内部使用 |\def| 而非 |\edef| 保存数值，因此实际边距在使用时（而非
% 设置时）才计算。
%
% \DescribeMacro{\SetCaptionFallback}
% 命令
% \begin{quote}
%   |\SetCaptionFallback|\marg{选项}\marg{回退值}
% \end{quote}
% 用于更改选项 \meta{选项} 的回退值，例如：
% \begin{quote}
%   |\SetCaptionFallback{labelformat}{simple}|
% \end{quote}
% 目前只有预定义标签格式 `original' 会在当前（浮动）环境未定义
% |\fnum@|\meta{环境} 时使用回退值。
%
% \DescribeMacro{\AtCaptionPackage}
% 适配文件会在加载 \pkg{caption} 内核 |caption3.sty| 时加载；因此即使
% 未加载 \pkg{caption} 宏包，所给适配仍会对 \pkg{floatrow} 与
% \pkg{subfig} 宏包生效。
%
% 但有些适配必须针对 \pkg{caption} 宏包进行；为此，适配文件作者可使用
% 钩子
% \begin{quote}|\AtCaptionPackage|\marg{代码}\end{quote}
% 钩子中的 \meta{代码} 会在所列宏包选项求值之前立即执行。
%
% 例如，\KOMAScript\ 文档类的适配文件使用
% |\At|\-|Begin|\-|Caption| 重定义 \pkg{caption} 宏包选项
% |figure|\-|position| 与 |table|\-|position|，以发出警告：在当前
% 环境下使用这些选项收效甚微。
%
% \DescribeMacro{\AfterCaptionPackage}
% 另一个附加钩子
% \begin{quote}|\AfterCaptionPackage|\marg{代码}\end{quote}
% 同样提供给适配文件作者使用。钩子中的 \meta{代码} 会在
% \pkg{caption} 宏包的最末尾执行。
%
% \iffalse
% --------------------------------------------------------------------------- %
% \fi
%
% \clearpage
% \section{供 babel 语言宏包作者使用的命令}
%
% \NEWfeature{v3.5}
% \DescribeMacro{\caption@switchdefault}
% 若已加载 \pkg{caption3} 内核宏包，则可使用命令
% \begin{quote}|\caption@switchdefault|\marg{选项}\marg{含 \#1 的代码}\marg{新默认值}\end{quote}
% 它通过 \meta{代码} 保存所给选项的旧默认值，随后设置新默认值。
%
% \DescribeMacro{\caption@restoredefault}
% 命令
% \begin{quote}|\caption@restoredefault|\marg{选项}\marg{旧值}\end{quote}
% 用于恢复此前存放在 \meta{旧值} 中的旧默认值。
%
% 示例：
% \begin{quote}
%   |\DeclareCaptionLabelSeparator{french}{|\ldots|}|\\
%   |\addto\extrasfrench{%|\\
%   |  \caption@switchdefault{labelseparator}{\babel@save #1}{french}}|\\
%   |% 注意：此处不需要 \caption@restoredefault|
% \end{quote}
%
% 另一示例：
% \begin{quote}
%   |\caption@switchdefault{font}{\let\original@captionfont#1}{it}%|\\
%   \ldots\\
%   |\caption@restoredefault{font}{\original@captionfont}%|\\
% \end{quote}
%
% 由于这两个命令在 \pkg{caption3} 内核内均以 |\def| 定义，可在 babel
% 语言宏包内预先定义：
% \begin{quote}
%   |\providecommand\caption@switchdefault[3]{}|\\
%   |\providecommand\caption@restoredefault[2]{}|
% \end{quote}
% 这样使用前便无需再检测它们是否存在。
%
% 此外，若已加载 \pkg{caption} 宏包（即已定义
% |\caption@makecaption|），babel 语言宏包不应重定义
% |\@make|\-|caption|。（遗憾的是，此处不能使用 \LaTeX{} 内核命令
% |\@if|\-|package|\-|loaded{caption}|，因为它只能在文档导言区使用。）
%
% 请注意，如下代码
% \begin{quote}
%   |\addto\extraswhatever{%|\\
%   |  \@ifundefined{caption@makecaption}{%|\\
%   |    \babel@save\@makecaption|\\
%   |    \let\@makecaption\my@makecaption|\\
%   |  }{%|\\
%   |    \caption@switchdefault{|\ldots|}{\babel@save #1}{|\ldots|}%|\\
%   |  }}|
% \end{quote}
% 是错误的，因为它混淆了 \pkg{caption3} 内核与 \pkg{caption} 宏包是否
% 存在。加载 \pkg{caption3} 内核并不等于也加载了 \pkg{caption} 宏包：
% 后者重定义 |\caption|、|\@caption| 与 |\@make|\-|caption|，而
% \pkg{caption3} 内核并不重定义其中任何一个，只为 \pkg{caption}、
% \pkg{floatrow} 与 \pkg{subfig} 宏包提供（默认）值。
%
% 目前内建支持以下 babel 语言选项：|arabic|、|farsi|、|french|、
% |hungarian|、|latvian| 与 |magyar|。可通过定义
% \begin{quote}
%   |\def\caption@ifbabel@|\meta{宏包}|{false}|
% \end{quote}
% 停用对应支持，其中 \meta{宏包} 可为 |arabi|、|french|、|frenchle|、
% |hungarian| 或 |latvian|。
%
% \iffalse
% --------------------------------------------------------------------------- %
% \fi
%
% \pagebreak[4]
% \section*{内部命令}
%
% 下列内部命令同样可供其他宏包安全使用，因为其接口在未来版本中不会
% 改变：
% \begin{quote}\begin{tabular}{ll}
%   内部命令                  & 当前使用者  \\
% \hline
%   |\caption@kernel@addcontentsline| & \pkg{floatrow} \\
%   |\caption@settype|                & \pkg{newfloat} \\
%   |\caption@xlabel|                 & \pkg{cleveref} \\
% \end{tabular}\end{quote}
%
% 若你的宏包希望使用 \pkg{caption} 宏包的其他内部命令，请发邮件联系，
% 以便共同定义合适的接口。
%
% \pkg{lwarp} 宏包同样使用了大量内部命令，相关工作正在 Git 分支
% \texttt{feature/lwarp} 上推进：
% \url{https://gitlab.com/axelsommerfeldt/caption/-/tree/feature/lwarp}
%
% \iffalse
% --------------------------------------------------------------------------- %
% \fi
%
% \StopEventually{%^^A
% \begin{thebibliography}{99}
%
%   \bibitem{TLC2}
%   Frank Mittelbach 与 Michel Goossens:\\
%   \newblock {\em 《{\LaTeX} 伴侣》（第 2 版）},\\
%   \newblock Addison-Wesley, 2004.
%
%   \bibitem{beamer}
%   Till Tantau:\\
%   \href{http://www.ctan.org/pkg/beamer}%
%        {\emph{《Beamer 类用户指南，3.07 版》}},\\
%   2007 年 3 月 11 日
%
%   \bibitem{KOMAScript}
%   Markus Kohm \& Jens-Uwe-Morawski:\\
%   \href{http://www.ctan.org/pkg/koma-script}%
%        {\emph{KOMA-Script —— 一套功能全面的 \LaTeXe\ 套装}},\\
%   2007-01-09
%
%   \bibitem{NTG}
%   Victor Eijkhout:\\
%   \href{http://www.ctan.org/pkg/ntgclass}%
%        {\emph{《荷兰语 \LaTeX\ 文档类入门》}},\\
%   1989 年 9 月 3 日
%
%   \bibitem{algorithms}
%   Rog\'erio Brito:\\
%   \href{http://www.ctan.org/pkg/algorithms}%
%        {\emph{《Algorithms》}},\\
%   2006 年 6 月 2 日
%
%   \bibitem{float}
%   Anselm Lingnau:\\
%   \href{http://www.ctan.org/pkg/float}%
%        {\emph{《改进的浮动体环境》}},\\
%   2001/11/08
%
%   \bibitem{floatflt}
%   Mats Dahlgren:\\
%   \href{http://www.ctan.org/pkg/floatflt}%
%        {\emph{《floatflt 宏包使用导引》}},\\
%   1998/06/05
%
%   \bibitem{floatrow}
%   Olga Lapko:\\
%   \href{http://www.ctan.org/pkg/floatrow}%
%        {\emph{《floatrow 宏包文档》}},\\
%   2009/08/02
%
%   \bibitem{fltpage}
%   Sebastian Gross:\\
%   \href{http://www.ctan.org/pkg/fltpage}%
%        {\emph{《欢迎试用 fltpage 宏包 beta 版！》}},\\
%   1998/11/13
%
%   \bibitem{hyperref}
%   Sebastian Rahtz \& Heiko Oberdiek:\\
%   \href{http://www.ctan.org/pkg/hyperref}%
%        {\emph{《\LaTeX 中的超文本标记：hyperref 使用手册》}},\\
%   2012 年 11 月
%
%   \bibitem{hypcap}
%   Heiko Oberdiek:\\
%   \href{http://www.ctan.org/pkg/oberdiek}%
%        {\emph{《hypcap 宏包 —— 调整题注锚点》}},\\
%   2011/02/16
%
%   \bibitem{listings}
%   Carsten Heinz \& Brooks Moses:\\
%   \href{http://www.ctan.org/pkg/listings}%
%        {\emph{《listings 宏包》}},\\
%   2007/02/22
%
%   \bibitem{longtable}
%   David Carlisle:\\
%   \href{http://www.ctan.org/pkg/longtable}%
%        {\emph{《longtable 宏包》}},\\
%   2004/02/01
%
%   \bibitem{picinpar}
%   Friedhelm Sowa:\\
%   \href{http://www.ctan.org/pkg/picinpar}%
%        {\emph{《段落中的插图》}},\\
%   1993 年 7 月 13 日
%
%   \bibitem{picins}
%   Joachim Bleser 与 Edmund Lang:\\
%   \href{http://www.ctan.org/pkg/picins}%
%        {\emph{《PicIns 用户手册 3.0 版》}},\\
%   1992 年 9 月
%
%   \bibitem{rotating}
%   Sebastian Rahtz 与 Leonor Barroca:\\
%   \href{http://www.ctan.org/pkg/rotating}%
%        {\emph{《\LaTeX 中旋转对象的样式选项》}},\\
%   1997/09/26
%
%   \bibitem{setspace}
%   Erica M. S. Harris \& Geoffrey Tobin:\\
%   \href{http://www.ctan.org/pkg/setspace}%
%        {\emph{《\LaTeX 文档宏包“setspace”》}},\\
%   2000 年 12 月 1 日
%
%   \bibitem{sidecap}
%   Rolf Niepraschk \& Hubert G\"a\ss lein:\\
%   \href{http://www.ctan.org/pkg/sidecap}%
%        {\emph{《sidecap 宏包》}},\\
%   2003/06/06
%
%   \bibitem{subfigure}
%   Steven D. Cochran:\\
%   \href{http://www.ctan.org/pkg/subfigure}%
%        {\emph{《subfigure 宏包》}},\\
%   2005/03/15
%
%   \bibitem{subfig}
%   Steven D. Cochran:\\
%   \href{http://www.ctan.org/pkg/subfig}%
%        {\emph{《subfig 宏包》}},\\
%   2005/07/05
%
%   \bibitem{supertabular}
%   Johannes Braams 与 Theo Jurriens:\\
%   \href{http://www.ctan.org/pkg/supertabular}%
%        {\emph{《supertabular 环境》}},\\
%   2004/02/20
%
%   \bibitem{threeparttable}
%   Donald Arseneau:\\
%   \href{http://www.ctan.org/pkg/threeparttable}%
%        {\emph{《三段式表格：标题、tabular 环境、注释》}},\\
%   2010 年 3 月 5 日
%
%   \bibitem{wrapfig}
%   Donald Arseneau:\\
%   \href{http://www.ctan.org/pkg/wrapfig}%
%        {\emph{《WRAPFIG.STY 3.6 版》}},\\
%   2003 年 1 月 31 日
%
%   \bibitem{xtab}
%   Peter Wilson:\\
%   \href{http://www.ctan.org/pkg/xtab}%
%        {\emph{《xtab 宏包》}},\\
%   2011/07/31
%
% \end{thebibliography}
% }
%
% \iffalse
% --------------------------------------------------------------------------- %
% \fi
%
% \DoNotIndex{\\,\_,\ ,\@@par}
% \DoNotIndex{\@bsphack}
% \DoNotIndex{\@car,\@cdr,\@classoptionslist,\@cons,\@currext,\@currname}
% \DoNotIndex{\@ehc,\@ehd,\@empty,\@esphack,\@expandtwoargs}
% \DoNotIndex{\@for,\@firstofone,\@firstoftwo}
% \DoNotIndex{\@gobble,\@gobblefour,\@gobbletwo,\@hangfrom}
% \DoNotIndex{\if@minipage,\@ifnextchar,\@ifpackagelater,\@ifpackageloaded}
% \DoNotIndex{\@ifstar,\@ifundefined,\@latex@error,\@minipagefalse,\@minipagetrue}
% \DoNotIndex{\@namedef,\@nameuse}
% \DoNotIndex{\@onlypreamble,\@parboxrestore,\@plus,\@ptionlist}
% \DoNotIndex{\@removeelement,\@restorepar,\@secondoftwo,\@setminipage,\@setpar}
% \DoNotIndex{\@tempa,\@tempboxa,\@tempdima,\@tempdimb,\@tempdimc,\@tempb,\@tempc}
% \DoNotIndex{\@testopt}
% \DoNotIndex{\@undefined,\@unprocessedoptions,\@unusedoptionlist}
% \DoNotIndex{\p@,\z@}
% \DoNotIndex{\active,\addtocounter,\addtolength,\advance,\aftergroup}
% \DoNotIndex{\baselineskip,\begin,\begingroup,\bfseries,\box}
% \DoNotIndex{\catcode,\centering,\changes,\csname,\def,\divide,\do,\downarrow}
% \DoNotIndex{\edef,\else,\empty,\end,\endcsname,\endgraf,\endgroup,\expandafter}
% \DoNotIndex{\fi,\footnotesize,\global}
% \DoNotIndex{\hangindent,\hbox,\hfil,\hsize,\hskip,\hspace,\hss}
% \DoNotIndex{\ifcase,\ifdim,\ifnum,\ifodd,\ifvoid,\ifvmode}
% \DoNotIndex{\ifx,\ignorespaces,\itshape}
% \DoNotIndex{\kernel@ifnextchar}
% \DoNotIndex{\Large,\large,\leavevmode,\leftmargini,\leftskip,\let,\linewidth}
% \DoNotIndex{\llap,\long,\m@ne,\margin,\mdseries,\message}
% \DoNotIndex{\newcommand,\newdimen,\newlength,\newline,\newif,\newsavebox}
% \DoNotIndex{\next,\nobreak,\nobreakspace,\noexpand,\noindent,\numberline}
% \DoNotIndex{\normalcolor,\normalfont,\normalsize,\or,\par,\parbox,\parfillskip}
% \DoNotIndex{\parindent,\parskip,\prevdepth,\protect,\protected@edef,\protected@write}
% \DoNotIndex{\providecommand,\quad}
% \DoNotIndex{\raggedleft,\raggedright,\relax,\renewcommand,\RequirePackage}
% \DoNotIndex{\rightskip,\rmfamily}
% \DoNotIndex{\sbox,\scriptsize,\scshape,\setbox,\setlength,\sffamily,\slshape}
% \DoNotIndex{\small,\string,\space,\strut}
% \DoNotIndex{\textheight,\the,\toks@,\typeout,\ttfamily}
% \DoNotIndex{\unvbox,\uparrow,\upshape,\usebox,\usepackage}
% \DoNotIndex{\value,\vbox,\vsize,\vskip,\wd,\width,\z@skip}
% \DoNotIndex{\AtBeginDocument,\AtEndOfPackage,\CurrentOption,\DeclareOption}
% \DoNotIndex{\ExecuteOptions,\GenericWarning,\IfFileExists,\InputIfFileExists}
% \DoNotIndex{\NeedsTeXFormat,\MessageBreak}
% \DoNotIndex{\PackageError,\PackageInfo,\PackageWarning,\PackageWarningNoLine}
% \DoNotIndex{\PassOptionsToPackage,\ProcessOptions,\ProvidesPackage}
%
% \iffalse
% --------------------------------------------------------------------------- %
% \fi
%
% \setlength{\parskip}{0pt plus 1pt}
% \renewcommand*\Note[2][Note]{\par{\small\emph{#1:} #2}\par}
% \let\endNote\undefined
%
% \changes{v1.0}{1994/10/27}{首次发布}
% \changes{v1.1}{1994/11/03}{新选项 \texttt{centerlast}}
% \changes{v1.2}{1994/11/28}{支持 \pkg{float} 宏包}
% \changes{v1.3}{1995/01/09}{子题注中支持 \cs{captionlabelfont}}
% \changes{v1.4}{1995/01/30}{新选项 \texttt{nooneline}}
% \changes{v1.4b}{1995/04/05}{缺陷修复版}
% \changes{v3.0}{2003/12/20}{重写；新增大量命令与功能}
% \changes{v3.0d}{2004/11/28}{拆分为两个宏包：\pkg{caption} 与 \pkg{caption3}}
%
% \iffalse
% --------------------------------------------------------------------------- %
% \fi
%
% \clearpage
% \section{实现}
%
% \iffalse
%<*package>
% \fi
%
% \subsection{引言}
%
% \Thispackage\ 由两部分组成——内核（|caption3.sty|）与主宏包
% （|caption.sty|）。
%
% \Thispackage\ 重定义 \LaTeX\ 命令 |\caption|、|\@caption| 与
% |\@makecaption|，并将后者映射到 |\caption@@make|，使用户能够控制
% |figure|、|table| 等浮动环境题注的外观。此外，对其他宏包（如
% \pkg{longtable} 或 \pkg{supertabular}）提供的题注机制也做了类似处理：
% 把相应内部命令（如 |\LT@makecaption| 或 |\ST@caption|）映射到
% \pkg{caption3} 内核所提供的命令。因此可以把 \thispackage\ 理解为一个
% 分层宏包：它在 \LaTeXe\ 或各宏包提供的题注机制与 \pkg{caption3} 内核
% 提供的题注机制之间提供适配层。
%
% \subsection{标识}
%
% 至少需要 \LaTeX2e\ 1994/12/01 版本。
%    \begin{macrocode}
\NeedsTeXFormat{LaTeX2e}[1994/12/01]
%    \end{macrocode}
%
% 若 \LaTeX\ 内核不认识版本声明，则跳过它们
% （依据 \url{https://www.latex-project.org/publications/2018-FMi-TUB-tb122mitt-version-rollback.pdf} 中的建议）。
%    \begin{macrocode}
\providecommand\DeclareRelease[3]{}
\providecommand\DeclareCurrentRelease[2]{}
%    \end{macrocode}
%
% 声明所有受支持的发行版本。
%    \begin{macrocode}
\DeclareRelease{v1}{1994-10-27}{caption_1995-04-05.sty}
\DeclareRelease{v3.0}{2003-12-20}{caption_2007-04-16.sty}
\DeclareRelease{v3.1}{2007-09-01}{caption_2010-01-09.sty}
\DeclareRelease{v3.2}{2011-08-06}{caption_2011-11-10.sty}
\DeclareRelease{v3.3}{2013-01-10}{caption_2019-09-01.sty}
\DeclareRelease{v3.4}{2019-11-24}{caption_2020-07-29.sty}
\DeclareRelease{v3.5}{2020-08-30}{caption_2020-10-26.sty}
\DeclareCurrentRelease{v3.6}{2022-02-20}
%    \end{macrocode}
%
% 标识当前宏包版本。
%    \begin{macrocode}
\ProvidesPackage{caption}[2023/08/05 v3.6o Customizing captions (AR)]
%    \end{macrocode}
%
% \subsection{加载内核}
%
% 加载版本匹配的 \pkg{caption} 内核。
%    \begin{macrocode}
\RequirePackage{caption3}[2023/07/31] % needs v2.4d or newer
%    \end{macrocode}
%
% \subsection{检查未知文档类}
%
%    \begin{macrocode}
\caption@ifdocumentclass{unknown}%
  {\caption@WarningNoLine{%
     Unknown document class (or package),\MessageBreak
     standard defaults will be used}}%
  {}
%    \end{macrocode}
%
% \subsection{检查不兼容的宏包}
%
% \changes{v3.0p}{2007/04/06}{新增对过时 \pkg{caption2} 宏包的检查}
%    \begin{macrocode}
\@ifpackageloaded{caption2}{%
  \caption@Error{%
    You can't use both, the (obsolete) caption2 *and*\MessageBreak
    the (current) caption package}%
  \endinput
}{}
%    \end{macrocode}
%
% 检查与 `caption' 宏包 \version{1} 时期曾能配合工作、现已不兼容的宏包。
% \changes{v3.1}{2007/06/15}{新增对不兼容 \pkg{ftcap} 宏包的检查}
% \changes{v3.1}{2007/06/15}{新增对不兼容 \pkg{nonfloat} 宏包的检查}
% \changes{v3.1}{2007/06/18}{新增对不兼容 \pkg{topcapt} 宏包的检查}
% \changes{v3.6}{2021/01/02}{修订对 \pkg{ftcap}、\pkg{nonfloat}、\pkg{topcapt} 等不兼容宏包的检查}
%    \begin{macrocode}
\caption@AtBeginDocument{%
  \caption@setincompatible{\@ifpackageloaded{captcont}}{captcont package}%
  \caption@setincompatible{\@ifpackageloaded{ftcap}}{ftcap package}%
  \caption@setincompatible{\@ifpackageloaded{nonfloat}}{nonfloat package}%
  \caption@setincompatible{\@ifpackageloaded{topcapt}}{topcapt package}}
%    \end{macrocode}
%
% \begin{macro}{\caption@setincompatible}
% \changes{v3.4}{2019/09/09}{新增本宏}
% \changes{v3.6}{2020/12/22}{适配回退方案}
% |\caption@setincompatible|\marg{条件}\marg{消息文本}
%    \begin{macrocode}
\newcommand*\caption@setincompatible[2]{%
  #1{%
    \caption@Error{%
      The #2 is incompatible to the\MessageBreak
      current version of the `caption' package,\MessageBreak
      so try \string\usepackage{caption}[=v1] instead}%
  }{}}
\@onlypreamble\caption@setincompatible
%    \end{macrocode}
% \end{macro}
%
% \subsection{位置}
%
%    \begin{macrocode}
\DeclareCaptionPosition{auto}{#2}{%
  \caption@ifsubcaption@above{#2}{%
    \ifvmode
      \caption@autoposition{#1}{#2}%
    \else
      #2%
    \fi}}
%    \end{macrocode}
%
% \subsection{选项声明}
%
% \subsubsection{杂项选项}
%
%    \begin{macrocode}
\DeclareCaptionOption*{config}[caption]{%
   \InputIfFileExists{#1.cfg}%
     {\typeout{*** Local configuration file #1.cfg used ***}}%
     {\caption@Warning{Configuration file #1.cfg not found}}}
%    \end{macrocode}
%
% \begin{macro}{\caption@selectlistentry}
% \changes{v3.2}{2010/11/07}{新增本宏}
%  |\caption@selectlistentry|\marg{标题或条目}
%    \begin{macrocode}
\newcommand*\caption@selectlistentry[1]{%
  \caption@ifinlist{#1}{heading}{%
    \let\caption@iflistheading\@firstoftwo
  }{\caption@ifinlist{#1}{default,list-entry,entry}{%
    \let\caption@iflistheading\@secondoftwo
  }{%
    \caption@Error{Undefined list-entry selection `#1'}%
  }}}
%    \end{macrocode}
% \end{macro}
%
% \changes{v3.3}{2012-03-25}{新增选项 \opt{list-entry}}
%    \begin{macrocode}
\DeclareCaptionOption{list-entry}{\caption@selectlistentry{#1}}
\captionsetup{list-entry=default}
%    \end{macrocode}
%
% \changes{v3.2}{2010/11/07}{新增选项 \opt{parboxrestore}}
% \changes{v3.6}{2021/01/07}{移除（未文档化的）选项 \opt{parboxrestore}}
%
% \changes{v3.1f}{2007/11/28}{新增选项 \opt{@minipage=}}
%    \begin{macrocode}
\DeclareCaptionOption{@minipage}{%
  \caption@ifinlist{#1}{auto,default}%
    {\let\caption@if@minipage\@gobbletwo}%
    {\caption@set@bool\caption@if@minipage{#1}}}
\captionsetup{@minipage=default}
%    \end{macrocode}
%
% \changes{v3.1}{2007/07/21}{新增选项 `compatibility'}
% \changes{v3.1}{2007/09/01}{修订 caption v1.x 兼容选项}
% \changes{v3.1d}{2007/10/24}{修复 07-10-24 版本中 caption v1.x$ 兼容选项的缺陷}
% \changes{v3.5a}{2020/09/02}{修复 caption v1.x 兼容选项}
% \changes{v3.6}{2020/12/22}{选项 `compatibility' 适配回退概念}
% \changes{v3.6h}{2022/07/03}{重新加入命令 \cs{caption@ifcompatibility}}
%
% |compatibility=| 选项仅支持布尔值，因为过去只有这些值被写入文档。
%
%    \begin{macrocode}
\DeclareCaptionOption{compatibility}[1]{%
  \caption@ifinlist{#1}{1,true,yes,on}{%
    \caption@Error{%
      Obsolete compatibility value `#1',\MessageBreak
      please use either \string\usepackage{caption}[=v1]\MessageBreak
      or \string\usepackage{caption-light} instead}%
  }{\caption@ifinlist{#1}{0,false,no,off}{%
      \let\caption@ifcompatibility\@secondoftwo % 不再使用；仅为向后兼容保留
  }{%
    \caption@Error{Undefined compatibility value `#1'}%
  }}}
\@onlypreamble@key{caption}{compatibility}
\caption@setbool{compatibility}{0} % 默认：尽量不与 v1.x 兼容
%    \end{macrocode}
%
% \subsubsection{caption v1.x 兼容选项}
%
%    \begin{macrocode}
\DeclareCaptionOptionNoValue*{normal}{%
  \caption@setup{format=plain,justification=justified}}
\DeclareCaptionOptionNoValue*{isu}{%
  \caption@setup{format=hang,justification=justified}}
\DeclareCaptionOptionNoValue*{hang}{%
  \caption@setup{format=hang,justification=justified}}
\DeclareCaptionOptionNoValue*{center}{%
  \caption@setup{format=plain,justification=centering}}
\DeclareCaptionOptionNoValue*{centerlast}{%
  \caption@setup{format=plain,justification=centerlast}}
%    \end{macrocode}
%
%    \begin{macrocode}
\DeclareCaptionOptionNoValue*{scriptsize}{\def\captionfont{\scriptsize}}
\DeclareCaptionOptionNoValue*{footnotesize}{\def\captionfont{\footnotesize}}
\DeclareCaptionOptionNoValue*{small}{\def\captionfont{\small}}
\DeclareCaptionOptionNoValue*{normalsize}{\def\captionfont{\normalsize}}
\DeclareCaptionOptionNoValue*{large}{\def\captionfont{\large}}
\DeclareCaptionOptionNoValue*{Large}{\def\captionfont{\Large}}
%    \end{macrocode}
%
%    \begin{macrocode}
\DeclareCaptionOptionNoValue*{up}{\caption@setup{labelfont+=up}}
\DeclareCaptionOptionNoValue*{it}{\caption@setup{labelfont+=it}}
\DeclareCaptionOptionNoValue*{sl}{\caption@setup{labelfont+=sl}}
\DeclareCaptionOptionNoValue*{sc}{\caption@setup{labelfont+=sc}}
\DeclareCaptionOptionNoValue*{md}{\caption@setup{labelfont+=md}}
\DeclareCaptionOptionNoValue*{bf}{\caption@setup{labelfont+=bf}}
\DeclareCaptionOptionNoValue*{rm}{\caption@setup{labelfont+=rm}}
\DeclareCaptionOptionNoValue*{sf}{\caption@setup{labelfont+=sf}}
\DeclareCaptionOptionNoValue*{tt}{\caption@setup{labelfont+=tt}}
%    \end{macrocode}
%
%    \begin{macrocode}
\DeclareCaptionOptionNoValue*{nooneline}{\caption@setup{slc=0}}
%    \end{macrocode}
%
%    \begin{macrocode}
\caption@setbool{ruled}{0}
\DeclareCaptionOptionNoValue*{ruled}{\caption@setbool{ruled}{1}}
%    \end{macrocode}
%
% \subsubsection{caption2 v2.x 兼容选项}
% \changes{v3.0h}{2005/07/08}{新增 caption2 \version{2.x} 兼容选项}
% \changes{v3.1}{2007/09/01}{修订 caption2 v2.x 兼容选项}
% \changes{v3.1d}{2007/10/24}{修复 07-10-24 版本中 caption v2.x$ 兼容选项的缺陷}
% \changes{v3.5a}{2020/09/02}{修复 caption v2.x 兼容选项}
%
%    \begin{macrocode}
\DeclareCaptionOptionNoValue*{flushleft}{%
  \caption@setup{format=plain,justification=raggedright}}
\DeclareCaptionOptionNoValue*{flushright}{%
  \caption@setup{format=plain,justification=raggedleft}}
%    \end{macrocode}
%    \begin{macrocode}
\DeclareCaptionOptionNoValue*{oneline}{\caption@setup{slc=1}}
%    \end{macrocode}
%    \begin{macrocode}
\DeclareCaptionOptionNoValue*{ignoreLTcapwidth}{%
  \caption@WarningNoLine{Obsolete option `ignoreLTcapwidth' ignored}}
%    \end{macrocode}
%
% \subsubsection{已过时的 caption v3.0 选项}
% \changes{v3.1}{2007/05/07}{宏包选项 \opt{caption} 标记为过时，
%        现改为报错}
%
%    \begin{macrocode}
\DeclareCaptionOption*{caption}{%
  \caption@setbool{temp}{#1}%
  \caption@ifbool{temp}{}{%
    \caption@Error{%
      The package option `caption=#1' is obsolete.\MessageBreak
      Please pass this option to the subfig package instead\MessageBreak
      and do *not* load the caption package anymore}}}
%    \end{macrocode}
%
% \subsubsection{fltpage 宏包支持选项}
%
% \changes{v3.1}{2007/09/01}{新增选项 `FPlist' 与 `FPref'}
% 这些选项用于控制目录条目以及 |\ref|、|\pageref| 或 |\autoref| 的链接
% 指向何处。默认值为 |FPlist=caption| 与 |FPref=figure|，二者并不一致，
% 但与 \pkg{fltpage} 宏包的常规行为保持兼容。
%    \begin{macrocode}
\DeclareCaptionOption{FPlist}[1]{\caption@setFPoption{list}{#1}}
\DeclareCaptionOption{FPref}[1]{\caption@setFPoption{ref}{#1}}
\@onlypreamble@key{caption}{FPlist}
\@onlypreamble@key{caption}{FPref}
%    \end{macrocode}
%    \begin{macrocode}
\newcommand*\caption@setFPoption[2]{%
  \edef\caption@tempa{\@car#2\@nil}%
  \caption@setbool{FP#1cap}{\if c\caption@tempa 1\else 0\fi}}
\@onlypreamble\caption@setFPoption
%    \end{macrocode}
%    \begin{macrocode}
\captionsetup{FPlist=caption,FPref=figure}
%    \end{macrocode}
%
% \subsubsection{hyperref 宏包支持选项}
%
% \changes{v3.1}{2007/06/21}{新增选项 `hypcap' 与 `hypcapspace'}
% 使用 |hypcap=off| 可关闭 \pkg{hypcap} 支持（默认为 |on|）。
%    \begin{macrocode}
\DeclareCaptionOption{hypcap}[1]{\caption@setbool{hypcap}{#1}}
\DeclareCaptionOption{hypcapspace}{\def\caption@hypcapspace{#1}}
%    \end{macrocode}
%    \begin{macrocode}
\captionsetup{hypcap=1,hypcapspace=.5\baselineskip}
%    \end{macrocode}
%
% \subsection{文档类支持}
%
% 执行通过 |\At|\-|Caption|\-|Package| 声明的代码。我们在处理选项之前
% 立即执行它，以便通过 |\At|\-|Caption|\-|Package| 完成的设置可被用户
% 选项覆盖。
%
%    \begin{macrocode}
\caption@documentclass@code
\let\caption@documentclass@code\@undefined
\let\AtCaptionPackage\@firstofone
%    \end{macrocode}
%
% \subsection{选项处理}
%
% \changes{v3.0m}{2007/03/04}{新增 verbose 选项}
% \changes{v3.1}{2007/07/02}{从宏包中移除 verbose 选项}
% \changes{v3.1}{2007/03/17}{不再处理全局选项}
%
%    \begin{macrocode}
\caption@SetupOptions{caption}{\caption@setkeys{#1}{#2}}
\caption@ProcessOptions*{caption}
%    \end{macrocode}
%
% \subsection{\cs{caption}、\cs{@caption} 与 \cs{@makecaption}}
%
% \begin{macro}{\caption@flags}
% \changes{v3.3}{2018/08/08}{新增本计数器及相关宏}
% 由于支持续接浮动体与子图，管理插图（及表格）计数器并非易事；尤其
% caption 宏包的早期版本已经证明：在局部修正计数器并非良策，而让其
% 依赖 |position=| 设置（如 \pkg{subfig} 宏包所做）同样不可取。
%
% 过去使用三个布尔开关存放标志，但由于部分宏包（如 \pkg{floatrow}
% 与 \pkg{tabu}）在某些情形下会竭力把 \LaTeX{} 计数器变化限制在局部，
% 因此把这些标志改为 \LaTeX{} 计数器值似乎更妥当。
%
% 另见：
% \url{https://gitlab.com/axelsommerfeldt/caption/-/issues/27} 与
% \url{https://tex.stackexchange.com/questions/444942/subcaptionbox-in-tabu-wrong-numbering-of-subfigures}
%
%    \begin{macrocode}
\newcounter{caption@flags}
%    \end{macrocode}
% \end{macro}
%
% \begin{macro}{\caption@clrflags}
% \changes{v3.3}{2018/08/08}{新增本宏}
% |\caption@clrflags| 清除所有标志。注意此处必须使用
% |\set|\-|counter|（而非 \TeX{} 原语），因为部分宏包会重定义
% |\set|\-|counter|，以把计数器变化限制在局部。
%    \begin{macrocode}
\newcommand*\caption@clrflags{%
  \setcounter{caption@flags}{0}}
%    \end{macrocode}
% \end{macro}
%
% \begin{macro}{\caption@ifflags}
% \changes{v3.6}{2020/12/28}{新增本宏}
% |\caption@ifflags|\marg{是代码}\marg{否代码}\\
% 检测是否设置了任一标志。
%    \begin{macrocode}
\newcommand*\caption@ifflags{%
  \ifcase\c@caption@flags
    \expandafter\@secondoftwo
  \else
    \expandafter\@firstoftwo
  \fi}
%    \end{macrocode}
% \end{macro}
%
% \begin{macro}{\caption@setflag}
% \changes{v3.3}{2013/04/14}{新增本宏}
% \changes{v3.3}{2018/08/08}{重写为使用计数器而非开关}
% |\caption@setflag|\marg{编号}\\
% 设置给定编号对应的标志（\meta{编号} 取 1、2、4、8、16、\ldots 之一）。
% 注意此处必须使用 |\add|\-|to|\-|counter|（而非 |\advance|），因为部分
% 宏包会重定义 |\add|\-|to|\-|counter|，以把计数器变化限制在局部。
%    \begin{macrocode}
\newcommand*\caption@setflag[1]{%
  \caption@ifflag{#1}{}{\addtocounter{caption@flags}{#1}}}
%    \end{macrocode}
% \end{macro}
%
% \begin{macro}{\caption@ifflag}
% \changes{v3.3}{2011/12/28}{新增本开关}
% \changes{v3.3}{2018/08/08}{重写（并改名）为使用计数器而非开关}
% |\caption@ifflag|\marg{编号}\marg{是代码}\marg{否代码}\\
% 检测给定编号对应的标志是否已设置（\meta{编号} 取 1、2、4、8、16、\ldots 之一）。
%    \begin{macrocode}
\newcommand*\caption@ifflag[1]{%
  \@tempcnta\c@caption@flags
  \divide\@tempcnta #1\relax
  \ifodd\@tempcnta
    \expandafter\@firstoftwo
  \else
    \expandafter\@secondoftwo
  \fi}
%    \end{macrocode}
% \end{macro}
%
% \begin{macro}{\caption@setcontinued}
% \changes{v3.6}{2020/09/20}{新增本宏}
% |\caption@setcontinued| 设置“续接浮动体”标志，即使用了
% |\continued|\-|float|。
%    \begin{macrocode}
\newcommand*\caption@setcontinued{\caption@setflag1}
%    \end{macrocode}
% \end{macro}
%
% \begin{macro}{\caption@ifcontinued}
% \changes{v3.3}{2011/12/28}{新增本开关}
% \changes{v3.3}{2018/08/08}{\cs{ifcaption@ContinuedFloat} 由 \cs{caption@ifflag1} 取代}
% \changes{v3.6}{2020/09/20}{以 \cs{caption@ifcontinued} 形式恢复}
% |\caption@if|\-|continued|\marg{是代码}\marg{否代码}\\
% 检测“续接浮动体”标志是否已设置。
%    \begin{macrocode}
\newcommand*\caption@ifcontinued{\caption@ifflag1}
%    \end{macrocode}
% \end{macro}
%
% \begin{macro}{\caption@setcaption}
% \changes{v3.6}{2020/09/20}{新增本宏}
% |\caption@setcaption| 设置“题注”标志，即使用了 |\caption|、
% |\phantom|\-|caption| 或 |\caption|\-|list|\-|entry|。
%    \begin{macrocode}
\newcommand*\caption@setcaption{\caption@setflag2}
%    \end{macrocode}
% \end{macro}
%
% \begin{macro}{\caption@ifcaption}
% \changes{v3.3}{2011/12/28}{新增本开关}
% \changes{v3.3}{2018/08/08}{\cs{ifcaption@caption} 由 \cs{caption@ifflag2} 取代}
% \changes{v3.6}{2020/09/20}{以 \cs{caption@ifcaption} 形式恢复}
% |\caption@if|\-|caption|\marg{是代码}\marg{否代码}\\
% 检测“题注”标志是否已设置。
%    \begin{macrocode}
\newcommand*\caption@ifcaption{\caption@ifflag2}
%    \end{macrocode}
% \end{macro}
%
% \begin{macro}{\caption@setsubcaption}
% \changes{v3.6}{2020/09/20}{新增本宏}
% |\caption@setsubcaption| 设置“子题注”标志，即使用了 |\subcaption|、
% |\subcaptionbox| 等命令。
%    \begin{macrocode}
\newcommand*\caption@setsubcaption{%
  \caption@ifcaption{\caption@setflag8}{\caption@setflag4}}
%    \end{macrocode}
% \end{macro}
%
% \begin{macro}{\caption@ifsubcaption@above}
% \changes{v3.3}{2011/12/28}{新增本开关}
% \changes{v3.3}{2018/08/08}{\cs{ifcaption@subcaption} 由 \cs{caption@ifflag4} 取代}
% \changes{v3.6}{2020/09/20}{以 \cs{caption@ifsubcaption@above} 与 \cs{caption@ifsubcaption@below} 形式恢复}
% |\caption@if|\-|sub|\-|caption@above|\marg{是代码}\marg{否代码}\\
% 检测“子题注”标志是否在“题注”标志之前设置，即子题注是否排版在相应
% 题注的\emph{上方}。
%    \begin{macrocode}
\newcommand*\caption@ifsubcaption@above{\caption@ifflag4}
%    \end{macrocode}
% \end{macro}
%
% \begin{macro}{\caption@ifsubcaption@below}
% \changes{v3.3}{2011/12/28}{新增本开关}
% \changes{v3.3}{2018/08/08}{\cs{ifcaption@subcaption} 由 \cs{caption@ifflag4} 取代}
% \changes{v3.6}{2020/09/20}{以 \cs{caption@ifsubcaption@above} 与 \cs{caption@ifsubcaption@below} 形式恢复}
% |\caption@if|\-|sub|\-|caption@below|\marg{是代码}\marg{否代码}\\
% 检测“子题注”标志是否在“题注”标志之后设置，即子题注是否排版在相应
% 题注的\emph{下方}。
%    \begin{macrocode}
\newcommand*\caption@ifsubcaption@below{\caption@ifflag8}
%    \end{macrocode}
% \end{macro}
%
% \begin{macro}{\caption@caption}
% \changes{v3.0f}{2005/05/22}{缺陷修复：加入 \cs{ContinuedFloat}}
% \changes{v3.0f}{2005/06/13}{移除对 \cs{caption*[]} 的支持}
% \changes{v3.0l}{2007/01/19}{加入对 \textsf{combine} 宏包的最小支持}
% \changes{v3.1}{2007/02/04}{重写，不再保存并使用原始定义}
% \changes{v3.6}{2022/01/05}{\cs{caption@Error} 的使用改为 \cs{caption@OutsideFloat}}
%  下面是 |\caption| 与 |\caption*| 的定义。除支持带星号形式外，这段代码
%  还按所支持的各类宏包做了适配。此处使用 |\caption@dblarg| 而非
%  |\@dblarg|，使 |\caption{}|（参数为空）仍会生成目录条目，而
%  |\caption[]{}| 不会。
%    \begin{macrocode}
\def\caption@caption{%
  \caption@iftype
    {\caption@checkgrouplevel\@empty\caption
     \caption@star
       {\caption@refstepcounter\@captype}%
       {\caption@dblarg{\@caption\@captype}}}%
    {\caption@OutsideFloat\caption
     \caption@withoptargs\@gobbletwo}}
%    \end{macrocode}
% \end{macro}
%
% \begin{macro}{\caption@OutsideFloat}
% \changes{v3.6}{2022/01/05}{新增本宏}
% |\caption@OutsideFloat}| 发出“\meta{命令} 位于浮动体之外”的错误。
%    \begin{macrocode}
\newcommand*\caption@OutsideFloat[1]{%
  \caption@Error{\string#1 outside float}}
%    \end{macrocode}
% \end{macro}
%
% \begin{macro}{\caption@star}
% \changes{v3.1}{2007/06/24}{由 \cs{caption@caption} 改名为 \cs{caption@star}}
% \changes{v3.2c}{2011/08/21}{移除 \cs{global}}
%  用于处理 |\caption| 后可选 |*| 的辅助宏。
%    \begin{macrocode}
\newcommand*\caption@star[2]{%
  \@ifstar{\caption@startrue#2[]}{#1#2}}
%    \end{macrocode}
% \end{macro}
%
% \begin{macro}{\caption@@caption}
% \changes{v3.1}{2007/02/04}{重写，不再保存并使用原始定义}
% \changes{v3.1f}{2007/11/28}{缺陷修复：加入对 \texttt{@minipage} 标志的处理}
% \changes{v3.1f}{2007/11/28}{加入对选项 \opt{@minipage} 的支持}
% \changes{v3.1g}{2008/03/16}{内部实现：改用 \cs{caption@prepareanchor}}
% \changes{v3.2d}{2011/10/21}{适配 \textsf{memoir} 文档类}
% \changes{v3.2f}{2012/01/14}{\cs{M@gettitle} 的使用移至 \cs{caption@prepareanchor}}
% \changes{v3.3}{2013/01/06}{适配 \textsf{nag} 宏包}
% \changes{v3.6}{2022/01/06}{加入 \cs{caption@nobreak} 的使用，以改进对 \textsf{beamer} 文档类的支持}
%  与上文相同，本版本已按所支持的宏包做了适配。此外，代码由
%  |\caption@begin|\-|ex| 与 |\caption@end| 嵌套，而非
%  |\begin|\-|group| 与 |\end|\-|group|；并使用
%  |\caption@par|\-|box|\-|restore| 代替 |\@parbox|\-|restore|。
%    \begin{macrocode}
\long\def\caption@@caption#1[#2]#3{%
%    \end{macrocode}
%    \begin{macrocode}
  \ifcaption@star \else
    \caption@prepareanchor{#1}{#2}%
    \memcaptioninfo{#1}{\csname the#1\endcsname}{#2}{#3}%
    \@nameuse{nag@hascaptiontrue}%
  \fi
%    \end{macrocode}
%    \begin{macrocode}
  \par
  \caption@nobreak
  \caption@beginex{#1}{#2}{#3}%
    \caption@setfloatcapt{%
      \caption@parboxrestore
      \if@minipage
        \@setminipage
      \fi
      \caption@normalsize
      \ifcaption@star
        \let\caption@makeanchor\@firstofone
      \fi
      \@makecaption{\csname fnum@#1\endcsname}%
                   {\ignorespaces\caption@makeanchor{#3}}\par
      \caption@nobreak
      \caption@if@minipage\@minipagetrue\@minipagefalse}%
  \caption@end}
%    \end{macrocode}
% \cls{memoir} 文档类相关代码：
%    \begin{macrocode}
\providecommand\memcaptioninfo[4]{}
%    \end{macrocode}
% \end{macro}
%
% \begin{macro}{\caption@prepareanchor}
% \changes{v3.1g}{2008/03/16}{新增本宏}
% \changes{v3.2f}{2012/01/14}{\cs{M@gettitle} 的使用从 \cs{caption@@caption} 移出}
%    \begin{macrocode}
\newcommand*\caption@prepareanchor[2]{%
  \caption@makecurrent{#1}{#2}%
  \caption@ifhypcap\caption@@start\relax
  \M@gettitle{#2}}
%    \end{macrocode}
% \cls{memoir} 文档类相关代码：
%    \begin{macrocode}
\providecommand\M@gettitle[1]{}
%    \end{macrocode}
% \end{macro}
%
% \begin{macro}{\caption@nobreak}
% \changes{v3.6}{2022/01/06}{新增本宏}
% |\caption@nobreak| 用于需要在 |\@caption| 或 |\beamer@makecaption|
% 执行 |\nobreak| 的位置使用。除非使用 \cls{beamer} 文档类，否则它定义为
% 空操作。
%    \begin{macrocode}
\newcommand*\caption@nobreak{}
%    \end{macrocode}
% \end{macro}
%
% \begin{macro}{\caption@parboxrestore}
% \changes{v3.1}{2007/03/03}{新增本宏定义}
% \changes{v3.1b}{2007/09/17}{缺陷修复：加入 CR 的重定义}
% \changes{v3.2}{2010/11/07}{加入对选项 \opt{parboxrestore} 的支持}
% \changes{v3.5g}{2020/10/26}{新实现；由 \cs{caption@boxrestore} 改名为 \cs{caption@parboxrestore}}
% \changes{v3.6}{2021/01/02}{现在始终使用 \cs{@parboxrestore}}
% |\caption@parboxrestore| 直接展开为 |\@par|\-|box|\-|restore|。
%    \begin{macrocode}
\newcommand*\caption@parboxrestore{\@parboxrestore}
%    \end{macrocode}
% \end{macro}
%
% \begin{macro}{\caption@makecaption}
% \changes{v3.1i}{2008/05/06}{加入 \cs{caption@rule} 的使用}
% \changes{v3.6}{2022/01/06}{加入 \cs{caption@nobreak} 的使用，以改进对 \textsf{beamer} 文档类的支持}
%  |\@makecaption|\marg{标签}\marg{文字}\par
%  其基本工作与（标准 \LaTeX\ 文档类的）原始代码相同，但会顾及
%  |position=| 设置，并使用 \pkg{caption} 内核中的 |\caption@@make|
%  最终排版题注。
%    \begin{macrocode}
\long\def\caption@makecaption#1#2{%
  \caption@nobreak
  \caption@make@above
  \caption@nobreak
  \caption@@make{#1}{#2}%
  \caption@nobreak
  \caption@make@below
  \caption@nobreak}
%    \end{macrocode}
% \end{macro}
%
% \begin{macro}{\caption@make@above}
% \begin{macro}{\caption@make@below}
% \changes{v3.3a}{2019/04/02}{新增本宏}
% |\caption@make@above|\\
% |\caption@make@below|\par
% 根据 |\caption@iftop| 排版题注上方或下方的内容。
%    \begin{macrocode}
\newcommand*\caption@make@above{%
  \caption@iftop
    \caption@belowskip
    {\caption@rule\caption@aboveskip}}
%    \end{macrocode}
%    \begin{macrocode}
\newcommand*\caption@make@below{%
  \caption@iftop
    {\caption@aboveskip\caption@rule}%
    \caption@belowskip}
%    \end{macrocode}
% \end{macro}
% \end{macro}
%
% \begin{macro}{\caption@redefine}
% \changes{v3.1}{2007/05/09}{加入对 \cs{@xfloat} 的重定义}
% \changes{v3.1b}{2007/09/22}{兼容模式下现在也重定义 \cs{@xfloat}}
% \changes{v3.1k}{2009/03/30}{修补代码现纳入宏 \cs{caption@redefine}}
% \changes{v3.5}{2020/08/30}{加入对 \cs{@xdblfloat} 的重定义}
%    \begin{macrocode}
\newcommand*\caption@redefine{%
%    \end{macrocode}
%   在 \version{3.4} 之前，只有当前定义为人所熟知时才重定义
%   |\caption| 与 |\@caption|，以便旧时（\pkg{caption} 宏包
%   \version{1.x} 时期，当时 |\caption| 与 |\@caption| 不由本宏包重定义）
%   编写的文档仍能顺利编译。2003 年引入 \version{3.0} 时这样做尚有
%   道理，但如今未免过于谨慎，也并非用户所需。因此从 \version{3.4} 起，
%   仅当 `compatibility' 选项设为 `true' 时才不重定义 |\caption| 与
%   |\@caption|；从 \version{3.6} 起则始终重定义。
%    \begin{macrocode}
  \let\caption\caption@caption
  \let\@caption\caption@@caption
%    \end{macrocode}
%   重定义 |\@x|\-|float|，以便在浮动环境内部使用按类型区分的选项、
%   设置 \pkg{hyperref} 锚点等。
%    \begin{macrocode}
  \let\caption@ORI@xfloat\@xfloat
  \let\@xfloat\caption@xfloat
%    \end{macrocode}
%   同样重定义 |\@xdbl|\-|float|，以便在此应用 |figure*| 或 |table*|
%   的特殊选项。
%    \begin{macrocode}
  \let\caption@ORI@xdblfloat\@xdblfloat
  \let\@xdblfloat\caption@xdblfloat
%    \end{macrocode}
%    \begin{macrocode}
}
%    \end{macrocode}
% 有些宏包（例如 \pkg{hyperref}）也会重定义 |\caption| 与 |\@caption|。
% 因此此处必须使用 |\At|\-|Begin|\-|Document|，以确保最终生效的是本
% 宏包的定义。
%    \begin{macrocode}
\caption@AtBeginDocument{%
  \caption@redefine
  \let\caption@redefine\relax}
%    \end{macrocode}
% |\@make|\-|caption| 的重定义立即执行。
%    \begin{macrocode}
\let\@makecaption\caption@makecaption
%    \end{macrocode}
% \end{macro}
%
% \begin{macro}{\caption@xfloat}
% \changes{v3.1}{2007/05/09}{加入对 \cs{@xfloat} 的重定义}
% \changes{v3.5}{2020/08/30}{新增并使用 \cs{caption@xfloat@hook}}
% \changes{v3.6}{2020/12/27}{对 \cs{@xfloat} 的重定义并入 \cs{caption@xfloat}}
%    \begin{macrocode}
\def\caption@xfloat#1[#2]{%
  \caption@ORI@xfloat{#1}[#2]%
  \caption@settype{#1}%
  \caption@setanchor
  \caption@xfloat@hook}
%    \end{macrocode}
% 钩子，可用 |\g@addto@macro\caption@x|\-|float@hook{|\ldots|}| 扩充。
% （题注类型可在 |\@cap|\-|type| 中找到。）
%    \begin{macrocode}
\newcommand*\caption@xfloat@hook{}
%    \end{macrocode}
% \end{macro}
%
% \begin{macro}{\caption@xdblfloat}
% \changes{v3.5}{2020/08/30}{加入对 \cs{@xdblfloat} 的重定义}
% \changes{v3.5}{2020/08/30}{新增并使用 \cs{caption@xdblfloat@hook}}
% \changes{v3.6}{2020/12/27}{对 \cs{@xdblfloat} 的重定义并入 \cs{caption@xdblfloat}}
%    \begin{macrocode}
\def\caption@xdblfloat#1[#2]{%
  \caption@ORI@xdblfloat{#1}[#2]% expands to \@xfloat{#1}[#2] + extra stuff
  \caption@setoptions{#1*}%
  \caption@xdblfloat@hook}
%    \end{macrocode}
% 钩子，可用 |\g@addto@macro\caption@xdbl|\-|float@hook{|\ldots|}| 扩充。
% （题注类型可在 |\@cap|\-|type| 中找到。）
%    \begin{macrocode}
\newcommand*\caption@xdblfloat@hook{}
%    \end{macrocode}
% \end{macro}
%
% \subsection{\cs{setcaptiontype} 与 \cs{setcaptionsubtype}}
%
% \changes{v3.0d}{2004/08/03}{新增选项 \opt{type=}}
% \changes{v3.1}{2007/06/25}{新增选项 \opt{type*=}}
% \changes{v3.1f}{2007/11/16}{新增选项 \opt{subtype} 与 \opt{subtype*}}
% \changes{v3.2}{2010/10/24}{缺陷修复 10-10-17：选项 \opt{type=} 现使用 \cs{setcaptiontype} 而非 \cs{caption@settype}}
%
%    \begin{macrocode}
\caption@AtBeginDocument{%
%    \end{macrocode}
%    \begin{macrocode}
  \DeclareCaptionOption{type}{\setcaptiontype{#1}}%
  \DeclareCaptionOption{type*}{\setcaptiontype*{#1}}%
%    \end{macrocode}
%    \begin{macrocode}
  \DeclareCaptionOptionNoValue{subtype}{\setcaptionsubtype\relax}%
  \DeclareCaptionOptionNoValue{subtype*}{\setcaptionsubtype*}%
%    \end{macrocode}
%    \begin{macrocode}
}
%    \end{macrocode}
%
% \emph{重要提示：}
% 与 |\captionof| 一样，选项 |type=| 只应在组、盒子或环境内部使用，
% 且不会检查参数是否为合法的浮动环境。
%
% \begin{macro}{\setcaptiontype}
% \changes{v3.2}{2010/10/24}{新增本宏}
% \changes{v3.5}{2020/08/24}{定义方式适配 \pkg{caption-light} 宏包}
% \changes{v3.6}{2022/02/20}{加入可选参数}
% \changes{v3.6e}{2022/04/17}{重写，修正任务顺序}
% \changes{v3.6f}{2022/04/18}{\cs{@ifnextchar} 由 \cs{kernel@ifnextchar} 取代}
%  |\setcaptiontype*|\oarg{附加选项}\marg{类型}\\
%  与 |\captionsetup{type=|\meta{类型}|}| 类似，但即使 |\caption|\-|setup|
%  已被重定义也能工作。不带星号的版本还会设置 \pkg{hyperref} 锚点
%  （当 |hypcap=|\-|true| 且未加载 \pkg{hypcap} 宏包时）。
%    \begin{macrocode}
\def\setcaptiontype{%
  \caption@teststar\@setcaptiontype\@gobble\@iden}
%    \end{macrocode}
%    \begin{macrocode}
\newcommand*\@setcaptiontype[1]{%
  \kernel@ifnextchar[%]
    {\@setcaptiontype@#1}%
    {\@@setcaptiontype#1\relax}}
%    \end{macrocode}
%    \begin{macrocode}
\def\@setcaptiontype@#1[#2]{%
  \@@setcaptiontype{#1}{\caption@setoptions{#2}}}
%    \end{macrocode}
%    \begin{macrocode}
\newcommand*\@@setcaptiontype[3]{%
  \caption@settype{#3}#2%
  \caption@parboxrestore@light
  #1\caption@setanchor}
%    \end{macrocode}
% \end{macro}
%
% \begin{macro}{\caption@settype}
% \changes{v3.1}{2007/05/09}{新增本宏}
% \changes{v3.1b}{2007/09/22}{加入 \cs{caption@xlabel} 的使用}
% \changes{v3.1d}{2007/10/23}{修订 \cs{caption@checkgrouplevel} 的错误消息}
% \changes{v3.1f}{2007/12/03}{替换粗劣的 \cs{caption@iftypewarning}}
% \changes{v3.1g}{2008/01/12}{加入 \cs{ifcaptionsetup@star} 的使用}
% \changes{v3.2}{2010/10/26}{错误消息从此处移至 \cs{caption@refstepcounter}}
% \changes{v3.3}{2011/12/28}{加入对 \cs{ifcaption@caption} 与 \cs{ifcaption@subcaption} 的重置}
% \changes{v3.4e}{2020/01/02}{加入检查，防止子类型被误用}
% \changes{v3.4e}{2020/01/02}{加入对子类型是否已定义的检查}
% \changes{v3.5d}{2020/09/28}{修订本宏}
% \changes{v3.6}{2020/12/27}{Some code out-sourced to \cs{caption@reset@currentlabel}}
% \changes{v3.6}{2021/01/02}{Usage of \cs{caption@initposition} added}
% \changes{v3.6}{2021/01/09}{Usage of \cs{flushsubcaptionlistentries} added}
% \changes{v3.6e}{2022/04/17}{移除带星号形式与可选参数}
% \changes{v3.6g}{2022/04/18}{移除 \cs{ifcaptionsetup@star} 的使用}
%  |\caption@settype|\marg{类型}\\
%  设置 |\@captype| 并执行与其关联的选项（使用 |\caption@set|\-|options|）。
%  此外还检查 |\current|\-|group|\-|level|（若可用），并使用宏
%  |\caption@|(|sub|)|type|\-|hook|（供本宏包对 \pkg{float} 的支持使用）。\par
%    \begin{macrocode}
\newcommand*\caption@settype{%
  \caption@initposition
  \caption@clrflags
  \aftergroup\flushsubcaptionlistentries
  \caption@@settype{}}
%    \end{macrocode}
%    \begin{macrocode}
\newcommand*\caption@@settype[2]{%
% #1 = "" 或 "sub"
% #2 = <类型>，例如："figure" 或 "subfigure"
  \caption@Debug{#1type:=#2}%
  \@nameuse{caption@check#1type}{#2}%
    {\caption@checkgrouplevel{#1}{\setcaption#1type}%
%    \end{macrocode}
%    \begin{macrocode}
     \edef\caption@tempa{#2}%
     \expandafter\ifx\csname @#1captype\endcsname\caption@tempa \else
       \@nameuse{caption@#1type@warning}%
     \fi
     \expandafter\let\csname @#1captype\endcsname\caption@tempa
%    \end{macrocode}
%    \begin{macrocode}
     \@nameuse{caption@#1typehook}%
%    \end{macrocode}
%    \begin{macrocode}
     \caption@setoptions{#2}%
     \ifx\caption@opt\relax
       \@nameundef{caption@#1type@warning}%
     \else
       \@namedef{caption@#1type@warning}{\caption@Warning{%
         The #1type was already set to
         `\csname @#1captype\endcsname'\MessageBreak}}%
     \fi}}
%    \end{macrocode}
% \end{macro}
%
% \begin{macro}{\setcaptionsubtype}
% \changes{v3.2a}{2011/08/15}{新增本宏}
% \changes{v3.5d}{2020/09/28}{修订本宏}
% \changes{v3.6}{2022/01/05}{\cs{caption@Error} 的使用改为 \cs{caption@OutsideFloat}}
% \changes{v3.6}{2022/01/07}{通过检查 \cs{caption@ifsubtype} 抑制重复使用}
% \changes{v3.6e}{2022/04/17}{重写}
%  |\setcaptionsubtype*|\\
%  与 |\captionsetup{subtype}| 类似，但即使 |\caption|\-|setup| 已被重定义
%  也能工作。设置 |\@sub|\-|cap|\-|type| 并执行与其关联的选项。不带星号
%  的版本还会设置 \pkg{hyperref} 锚点（当 |hypcap=|\-|true| 且未加载
%  \pkg{hypcap} 宏包时）。
%    \begin{macrocode}
\newcommand*\setcaptionsubtype{%
  \caption@teststar\@setcaptionsubtype\@gobble\@iden}
%    \end{macrocode}
%    \begin{macrocode}
\newcommand*\@setcaptionsubtype[1]{%
  \caption@iftype
    {\caption@ifsubtype
       {\caption@Debug{type=\@captype, subtype=\@subcaptype}}%
       {\caption@@settype{sub}{sub\@captype}%
        #1\caption@setanchor}}%
    {\caption@OutsideFloat\setcaptionsubtype}}
%    \end{macrocode}
% \end{macro}
%
% \begin{macro}{\caption@parboxrestore@light}
% \changes{v3.2}{2010/10/24}{新增本宏}
% \changes{v3.5g}{2020/10/26}{由 \cs{caption@boxrestore@mini} 改名为 \cs{caption@parboxrestore@mini}}
% \changes{v3.6}{2021/01/02}{由 \cs{caption@parboxrestore@mini} 改名为 \cs{caption@parboxrestore@light}}
% 重置 |\par|，使 |\@caption| 中第一个 |\par| 的行为与浮动环境中基本一致。
% 供 |\set|\-|caption|\-|type| 使用。
%    \begin{macrocode}
\newcommand*\caption@parboxrestore@light{%
  \let\par\@@par
  \parindent\z@ \parskip\z@skip
  \sloppy}
%    \end{macrocode}
% \end{macro}
%
% \begin{macro}{\caption@setanchor}
% \changes{v3.6e}{2022/04/17}{代码自 \cs{caption@settype} 中抽出}
%  |\caption@setanchor|\\
%  设置 \pkg{hyperref} 锚点。
%    \begin{macrocode}
\newcommand*\caption@setanchor{%
  \caption@reset@currentlabel
  \caption@start}
%    \end{macrocode}
% \end{macro}
%
% \begin{macro}{\caption@reset@currentlabel}
% \changes{v3.6}{2020/12/27}{本宏自 \cs{caption@settype} 中分离}
%  重定义 |\@currentlabel|，使 |\caption| 之前的 |\label| 给出提示而非
%  错误引用。
%    \begin{macrocode}
\newcommand*\caption@reset@currentlabel{%
  \let\@currentlabel\caption@undefinedlabel
% \let\@currentHlabel\@undefined
  \ifx\caption@x@label\@undefined
    \let\caption@x@label\label
    \let\label\caption@xlabel
  \fi}
%    \end{macrocode}
% \end{macro}
%
% \begin{macro}{\caption@iftype}
% \changes{v3.1}{2007/04/11}{新增本宏定义}
%  由于经常需要检查 |\@cap|\-|type| 是否已定义（即是否位于浮动环境
%  内部），引入本辅助宏。
%    \begin{macrocode}
\newcommand*\caption@iftype{\caption@ifdefined\@captype}
%    \end{macrocode}
% \end{macro}
%
% \begin{macro}{\caption@ifsubtype}
% \changes{v3.6}{2022/01/07}{新增本宏定义}
%  由于经常需要检查 |\@sub|\-|cap|\-|type| 是否已定义（即是否位于浮动
%  环境的子题注组内部），引入本辅助宏。
%    \begin{macrocode}
\newcommand*\caption@ifsubtype{\caption@ifdefined\@subcaptype}
%    \end{macrocode}
% \end{macro}
%
% \begin{macro}{\caption@checktype}
% \changes{v3.4g}{2020/01/02}{新增本宏定义}
% \changes{v3.4m}{2020/07/29}{加入对 \pkg{subfig} 宏包中 \cs{sf@counterlist} 的检查}
% |\caption@checktype|\marg{类型}\marg{代码}\par
% 供 |\caption@settype| 使用：给定类型并非正规类型（而是子类型，即由
%  |\DeclareCaptionSubType| 或 |\newsubfloat|\cite{subfig} 定义）时，
%  展开所给代码或发出错误。
%    \begin{macrocode}
\newcommand*\caption@checktype[1]{%
  \caption@ifdefined@subtype{#1}\@firstoftwo{\caption@ifin@list\sf@counterlist{#1}}%
    {\caption@Error{Should not be used with subtype `#1'}%
     \@gobble}%
    \@firstofone}
%    \end{macrocode}
% \end{macro}
%
% \begin{macro}{\caption@checksubtype}
% \changes{v3.4g}{2020/01/02}{新增本宏定义}
% \changes{v3.4m}{2020/07/29}{加入对 \pkg{subfig} 宏包中 \cs{sf@counterlist} 的检查}
% \changes{v3.6c}{2022/04/06}{移除对 \pkg{subfig} 宏包中 \cs{sf@counterlist} 的检查}
% \changes{v3.6c}{2022/04/06}{加入 \cs{caption@subtypesource} 的使用}
% |\caption@checksubtype|\marg{类型}\marg{代码}\par
% 供 |\caption@settype| 使用：给定类型并非正规子类型（即未由
%  |\Declare|\-|Caption|\-|Sub|\-|Type| 定义）时，展开所给代码或发出错误。
%    \begin{macrocode}
\newcommand*\caption@checksubtype[1]{%
  \caption@ifdefined@subtype{#1}%
    \@firstofone
    {\@ifundefined{c@#1}%
       {\caption@Error{Undefined counter `#1'}}%
       {\caption@subtypesource\caption@subtype@source{#1}%
        \caption@Error{The counter `#1' was defined by\MessageBreak\caption@subtype@source}}%
     \@gobble}}
%    \end{macrocode}
% \end{macro}
%
% \begin{macro}{\caption@checkgrouplevel}
% \changes{v3.1}{2007/07/15}{新增本宏定义}
% \changes{v3.1f}{2007/11/28}{错误改为警告（兼容）}
% \changes{v3.3}{2013/04/21}{现使用 \cs{caption@ifeTeX}}
% \changes{v3.5}{2019/09/11}{现依赖 \eTeX{}}
% 检查 |\captionsetup{type=|\ldots|}| 或 |\caption| 是否在组内部执行；
% 若不在组内部，则发出警告消息。
%    \begin{macrocode}
\newcommand*\caption@checkgrouplevel[2]{%
  \@ifundefined{#1caption@grouplevel}{%
     \caption@ifundefined\caption@grouplevel{\let\caption@grouplevel\z@}{}%
     \ifnum\currentgrouplevel>\caption@grouplevel\relax
       \expandafter\edef\csname #1caption@grouplevel\endcsname{%
         \the\currentgrouplevel}%
     \else
       \caption@Warning{\string#2\MessageBreak outside box or environment}%
     \fi
  }{}}%
%    \end{macrocode}
% \end{macro}
%
% \begin{macro}{\caption@undefinedlabel}
% \changes{v3.1}{2007/07/03}{新增本宏定义}
% \changes{v3.1b}{2007/09/22}{内容改为 \texttt{??}}
% \changes{v3.1i}{2008/04/13}{内容改为 \cs{caption@xref}(\ldots)}
% \changes{v3.1k}{2009/03/30}{缺陷修复 08-12-08：\cs{caption@xref} 现为健壮命令}
% \changes{v3.2c}{2011/09/13}{文字由 `before' 改为 `without proper'}
% \changes{v3.3}{2016/02/21}{文字由 `without proper \cs{caption}' 改为 `without proper reference'}
% 在（浮动）环境内部，|\\currentlabel| 默认使用本标签。（见上文）
%  \iffalse\Note{此处所用命令都必须健壮。}\fi
%    \begin{macrocode}
\newcommand*\caption@undefinedlabel{%
  \protect\caption@xref{\caption@labelname}{\on@line}}
%    \end{macrocode}
%    \begin{macrocode}
\DeclareRobustCommand*\caption@xref[2]{%
  \caption@WarningNoLine{\noexpand\label without proper reference#2}%
  \@setref\relax\@undefined{#1}}
%    \end{macrocode}
%    \begin{macrocode}
\newcommand*\caption@labelname{??}
%    \end{macrocode}
% \end{macro}
%
% \begin{macro}{\caption@xlabel}
% \changes{v3.1b}{2007/09/22}{新增本宏定义}
% \changes{v3.1f}{2007/11/11}{修订}
% \changes{v3.1i}{2008/04/13}{缺陷修复：\cs{pageref} 现在应始终可用}
% \changes{v3.2c}{2011/08/19}{现在也支持可选参数}
% 浮动环境内 |\label| 的新实现。|\label| 通过 |\caption@withoptargs|
% 重定义，因此 |#1| 为可选参数（若有），|#2| 为必填参数。
%    \begin{macrocode}
\newcommand*\caption@xlabel{%
  \caption@withoptargs\caption@@xlabel}
%    \end{macrocode}
%    \begin{macrocode}
\newcommand*\caption@@xlabel[2]{%
  \caption@@@xlabel
  \def\caption@labelname{#2}%
  \caption@x@label#1{#2}}
%    \end{macrocode}
%    \begin{macrocode}
\newcommand*\caption@@@xlabel{%
  \global\let\caption@@@xlabel\@empty
  \@bsphack
    \protected@write\@auxout{}%
      {\string\providecommand*\string\caption@xref[2]{%
        \string\@setref\string\relax\string\@undefined{\string##1}}}%
  \@esphack}
%    \end{macrocode}
% \end{macro}
%
% \begin{macro}{\captionof}
% \changes{v3.1}{2007/05/09}{改用 \cs{caption@settype}，而非 \cs{def}\cs{@captype}}
% \changes{v3.1c}{2007/09/22}{缺陷修复：移除对 \cs{caption@setfloatcapt} 的重定义}
% \changes{v3.1e}{2007/10/30}{更好的兼容性：现通过 \cs{AtBeginDocument} 定义}
% \changes{v3.2}{2010/10/24}{缺陷修复 10-10-17：现使用 \cs{setcaptiontype} 而非 \cs{caption@settype}}
% \changes{v3.5}{2020/08/24}{定义方式适配 \pkg{caption-light} 宏包}
%  |\captionof|\marg{类型}\oarg{目录条目}\marg{标题}\\
%  |\captionof*|\oarg{目录条目}\marg{标题}\par
% \Note{本命令将通过 \cs{AtBeginDocument} 定义，因此
%  \cs{usepackage}\csmarg{caption,capt-of} 仍可使用。
%  （与 \version{1.x} 兼容）}
%    \begin{macrocode}
\caption@AtBeginDocument{%
  \def\captionof{\caption@teststar\caption@of{\caption*}\caption}}
%    \end{macrocode}
%    \begin{macrocode}
\providecommand*\caption@of[2]{\setcaptiontype*{#2}#1}
%    \end{macrocode}
% \end{macro}
%
% \subsection{\cs{phantomcaption}}
%
% \begin{macro}{\phantomcaption}
% \changes{v3.2}{2010/10/31}{缺陷修复 10-09-01b：新增本宏}
% \changes{v3.6}{2022/01/05}{\cs{caption@Error} 的使用改为 \cs{caption@OutsideFloat}}
%  |\phantomcaption|\par
%  使计数器加一并生成可供 |\label| 使用的引用，但不排版任何内容。
%    \begin{macrocode}
\newcommand*\phantomcaption{%
  \caption@iftype
    {\caption@refstepcounter\@captype}%
    {\caption@OutsideFloat\phantomcaption}}
%    \end{macrocode}
% \end{macro}
%
% \subsection{\cs{captionlistentry}}
%
% \begin{macro}{\captionlistentry}
% \changes{v3.1}{2007/07/28}{新增本宏}
% \changes{v3.1f}{2007/11/16}{缺陷修复 07-11-09：加入“空格技巧”}
% \changes{v3.1g}{2008/01/20}{新增带星号形式}
% \changes{v3.1l}{2010/01/09}{适配当前版本的 nameref 宏包}
% \changes{v3.6}{2020/12/27}{带可选参数时使用 \cs{caption@refstepcounter@} 而非 \cs{caption@refstepcounter}}
% \changes{v3.6}{2022/01/05}{\cs{caption@Error} 的使用改为 \cs{caption@OutsideFloat}}
% \changes{v3.6f}{2022/04/18}{\cs{@ifnextchar} 由 \cs{kernel@ifnextchar} 取代}
%  |\captionlistentry|\oarg{浮动体类型}\marg{目录条目}\\
%  |\captionlistentry*|\oarg{浮动体类型}\marg{目录条目}
%    \begin{macrocode}
\newcommand*\captionlistentry{%
  \caption@teststar\caption@listentry\@firstoftwo\@secondoftwo}
%    \end{macrocode}
%    \begin{macrocode}
\newcommand*\caption@listentry[1]{%
  \kernel@ifnextchar[%]
    {\caption@listentry@{#1}}%
    {\caption@@listentry{#1}}}
%    \end{macrocode}
%    \begin{macrocode}
\def\caption@listentry@#1[#2]{%
  \caption@@@listentry{#1}{#2}{\caption@refstepcounter@}}
%    \end{macrocode}
%    \begin{macrocode}
\newcommand*\caption@@listentry[1]{%
  \caption@iftype
    {\caption@@@listentry{#1}{\@captype}{\caption@refstepcounter}}%
    {\caption@OutsideFloat\captionlistentry
     \@gobble}}
%    \end{macrocode}
%    \begin{macrocode}
\newcommand*\caption@@@listentry[4]{%
  \@bsphack
    #1{\caption@gettitle{#4}}%
      {#3{#2}\caption@makecurrent{#2}{#4}}%
    \caption@addcontentsline{#2}{#4}%
  \@esphack}
%    \end{macrocode}
% \end{macro}
%
% \subsection{\cs{captiontext}}
%
% \begin{macro}{\captiontext}
% \changes{v3.6}{2022/01/05}{新增本宏}
% \changes{v3.6}{2022/01/05}{\cs{caption@Error} 的使用改为 \cs{caption@OutsideFloat}}
%  |\captiontext*|\oarg{编号}\marg{题注文字}\par
%  排版题注，但不使计数器加一、不写入目录、不添加间距。带星号形式还会
%  跳过单行检查，且不使用盒子、parbox 或边距。
%    \begin{macrocode}
\newcommand*\captiontext{%
  \caption@iftype
    {\caption@text\@captype}%
    {\caption@OutsideFloat\captiontext
     \caption@withoptargs\@gobbletwo}}
\newcommand*\caption@text[1]{%
  \caption@teststar{\caption@@text{#1}}\caption@@@make\caption@@make}
\newcommand*\caption@@text[2]{%
  \kernel@ifnextchar[%]
    {\caption@@@text{#1}#2\@iden}%
    {\caption@@@text{#1}#2\@gobble[0]}}
\long\def\caption@@@text#1#2#3[#4]#5{%
  \begingroup
    #3{\csname c@#1\endcsname #4\relax}%
    #2{\caption@fnum{#1}}{#5}%
  \endgroup}
%    \end{macrocode}
% \end{macro}
%
% \subsection{\env{captiongroup} 环境}
%
% \begin{environment}{captiongroup}
% \changes{v3.6}{2022/02/20}{新增本环境定义}
% \env{captiongroup} 只是一个内含 |\set|\-|caption|\-|type| 的环境，仅此
% 而已。本环境也提供带星号形式，使用 |\set|\-|caption|\-|type*|。
%    \begin{macrocode}
\newenvironment{captiongroup}{\setcaptiontype}{}
\newenvironment{captiongroup*}{\setcaptiontype*}{}
%    \end{macrocode}
% \end{environment}
%
% \subsubsection{\env{captionblock} 环境}
%
% \begin{macro}{\caption@minipage}
% \changes{v3.6}{2022/02/20}{新增本宏定义}
% 这是一个以 |\set|\-|caption|\-|type| 为首行内容的 \env{minipage}。
%    \begin{macrocode}
\newcommand*\caption@minipage{%
  \@testopt\caption@iminipage b}
%    \end{macrocode}
% 其余所有可选参数都通过 \pkg{caption} 内核提供的通用辅助宏
%  |\caption@with|\-|opt|\-|args| 传递。
%    \begin{macrocode}
\def\caption@iminipage[#1]{%
  \caption@withoptargs{\caption@iiminipage{#1}}}
%    \end{macrocode}
% `B' 与 `T' 会加入 |\vspace{0pt}|；其他取值（及附加可选参数）原样传递
% 给 |minipage| 环境。
%    \begin{macrocode}
\newcommand*\caption@iiminipage[3]{%
  \let\caption@endminipage@hook\@empty
  \if#1B%
    \minipage[b]#2{#3}%
    \def\caption@endminipage@hook{\vspace{0pt}}%
  \else\if#1T%
    \minipage[t]#2{#3}%
    \vspace{0pt}%
  \else
    \minipage[#1]#2{#3}%
  \fi\fi
  \@floatboxreset
  \setcaptiontype}
%    \end{macrocode}
%    \begin{macrocode}
\newcommand*\caption@endminipage{%
  \caption@endminipage@hook
  \endminipage}
%    \end{macrocode}
% \end{macro}
%
% \begin{environment}{captionblock}
% \changes{v3.6}{2022/02/20}{新增本环境定义}
% \env{captionblock} 是一个以 |\set|\-|caption|\-|type| 为首行内容的
%  \env{minipage}。
%    \begin{macrocode}
\newenvironment{captionblock}{\caption@minipage}{\caption@endminipage}
%    \end{macrocode}
% \end{environment}
%
% \subsection{\cs{captionbox}}
%
% \begin{macro}{\captionbox}
% \changes{v3.2}{2010/12/17}{新增本宏}
% \changes{v3.2a}{2011/08/16}{\cs{captionsetup} 由 \cs{captionbox@settype} 与 \cs{caption@setposition} 取代}
% \changes{v3.3}{2012/04/09}{适配修订版 \cs{caption@withoptargs}}
% \changes{v3.3}{2017/03/19}{固定宽度使用时，内容不再放入 hbox}
% \changes{v3.3c}{2019/09/01}{适配为可由 \pkg{subcaption} 提供的 \cs{subfloat} 使用}
% \changes{v3.5}{2020/07/28}{适配 \pkg{caption3} v2.0}
% \changes{v3.6}{2020/12/24}{加入对参数 \meta{内部位置} 的检查}
% 一个含内容与题注的 |\parbox|，二者以一条不可见的 |\hrule| 分隔。
%    \begin{macrocode}
\newcommand*\captionbox{%
  \caption@withoptargs{\caption@ibox\@gobble}}
%    \end{macrocode}
%    \begin{macrocode}
\newcommand\caption@ibox[3]{%
  \kernel@ifnextchar[%]
    {\caption@iibox{#1}{#2}{#3}}%
    {\caption@iibox@{#1}{#2}{#3}}}
%    \end{macrocode}
%    \begin{macrocode}
\long\def\caption@iibox#1#2#3[#4]{%
  \@testopt{\caption@iiibox{#1}{#2}{#3}[{#4}]}\captionbox@innerpos@default}
%    \end{macrocode}
%    \begin{macrocode}
\long\def\caption@iibox@#1#2#3#4{%
  \setbox\@tempboxa\hbox{#4}%
  \caption@iiibox{#1}{#2}{#3}[\wd\@tempboxa][\captionbox@innerpos@default]{\unhbox\@tempboxa}}
%    \end{macrocode}
%    \begin{macrocode}
\long\def\caption@iiibox#1{%
  \caption@iiiibox{#1}\vbox\vtop}
%    \end{macrocode}
%    \begin{macrocode}
\long\def\caption@iiiibox#1#2#3#4#5[#6][#7]#8{%
  \caption@checkdecl{justification}{#7}{%
    \begingroup
    #1*% set \caption@position
    \caption@iftop{%
      \endgroup
      \parbox[t]{#6}{%
        #1\relax
        \caption@setposition t%
        #2{\caption#4{#5}}%
        \captionbox@hrule
        \csname caption@justification@#7\endcsname
        #8}%
    }{%
      \endgroup
      \parbox[b]{#6}{%
        #1\relax
        \caption@setposition b%
        \csname caption@justification@#7\endcsname
        #8%
        \captionbox@hrule
        #3{\caption#4{#5}}}%
    }}}
%    \end{macrocode}
%    \begin{macrocode}
\newcommand*\captionbox@innerpos@default{c}
\newcommand*\captionbox@hrule{\hrule\@height\z@\relax}
%    \end{macrocode}
%    \begin{macrocode}
\providecommand*\caption@justification@c{\centering}
\providecommand*\caption@justification@l{\raggedright}
\providecommand*\caption@justification@r{\raggedleft}
\providecommand*\caption@justification@s{}
%    \end{macrocode}
% \end{macro}
%
% \subsection{\cs{continuedfloat}}
%
% \begin{macro}{\c@continuedfloat}
% \changes{v3.4}{2015/09/16}{计数器由 \cs{ContinuedFloat} 改名为 \cs{continuedfloat}}
% |\thecontinuedfloat| 预设为 |\@empty|，因此续接计数器通常不会纳入题注
% 标签或引用。
%    \begin{macrocode}
\newcounter{continuedfloat}
\let\c@ContinuedFloat\c@continuedfloat     % 为兼容而定义
%    \end{macrocode}
%    \begin{macrocode}
\def\thecontinuedfloat{\theContinuedFloat}
\let\theContinuedFloat\@empty              % 为兼容而定义
%    \end{macrocode}
% \end{macro}
%
% \begin{macro}{\continuedfloat@captype}
% \changes{v3.6}{由 \cs{caption@CFtype} 改名为 \cs{continuedfloat@captype}}
% 由于只有一个 `continuedfloat' 计数器（而非每种题注类型各一个），需要把
% 题注类型绑定到该计数器，确保计数器持有非零值时不会被（误）用于其他
% 浮动体类型。
%    \begin{macrocode}
\newcommand*\continuedfloat@captype{??}
%    \end{macrocode}
% \end{macro}
%
% \begin{macro}{\caption@reset@continuedfloat}
%  |\caption@reset@continuedfloat|\marg{题注类型}\\
%  将续接计数器重置为零，并把所给题注类型绑定到该计数器。
%    \begin{macrocode}
\newcommand*\caption@reset@continuedfloat[1]{%
  \xdef\continuedfloat@captype{#1}%
  \@stpelt{continuedfloat}}
%    \end{macrocode}
% \end{macro}
%
% \begin{macro}{\caption@set@continuedfloat}
% \changes{v3.6}{2020/12/27}{\cs{continuedfloat} 的部分代码抽出到本宏}
%  |\caption@set@continuedfloat|\marg{是代码}\marg{否代码}\\
%  在递增续接计数器之前，检查当前题注类型是否与绑定到
%  `continuedfloat' 计数器的类型一致。
%    \begin{macrocode}
\newcommand*\caption@set@continuedfloat{%
  \ifx\@captype\continuedfloat@captype
    \stepcounter{continuedfloat}%
    \expandafter\@firstoftwo
  \else
    \caption@Error{Continued `\@captype' after `\continuedfloat@captype'}%
    \expandafter\@secondoftwo
  \fi}
%    \end{macrocode}
% \end{macro}
%
% \begin{macro}{\continuedfloat}
% \changes{v3.0i}{2005/11/12}{加入可选参数}
% \changes{v3.0g}{2005/06/28}{缺陷修复：\cs{newcounter} 由 \cs{newcount} 取代}
% \changes{v3.1}{2007/07/02}{更改 \pkg{hyperref} 扩展格式}
% \changes{v3.1}{2007/07/07}{加入 \cs{caption@setoptions}}
% \changes{v3.1g}{2008/01/29}{新增带星号形式}
% \changes{v3.2f}{2011/12/17}{移除未文档化的可选参数}
% \changes{v3.2f}{2011/12/17}{加入对 \cs{caption@@@continuedfloat} 的重定义}
% \changes{v3.3}{2011/12/28}{现使用 \cs{ifcaption@ContinuedFloat}}
% \changes{v3.3}{2013/05/01}{加入 \cs{caption@restorecounters} 的使用}
% \changes{v3.3}{2013/12/16}{缺陷修复：不再使用 \cs{caption@ContinuedFloat}，因其会破坏 \pkg{subfig} 提供的 \cs{ContinuedFloat}}
% \changes{v3.4}{2015/09/16}{由 \cs{ContinuedFloat} 改名为 \cs{continuedfloat}}
% \changes{v3.5f}{2020/10/06}{移除 \cs{caption@restorecounters} 的使用}
% \changes{v3.6}{2020/12/27}{计数器题注类型的检查抽出到 \cs{caption@check@continuedfloat}}
% \changes{v3.6}{2022/01/05}{\cs{caption@Error} 的使用改为 \cs{caption@OutsideFloat}}
%  |\continuedfloat|\\
%  |\continuedfloat*|\par
%  主要用于设置“续接浮动体”标志、递增续接计数器，并执行所给选项。
%    \begin{macrocode}
\def\continuedfloat{%
  \caption@teststar\caption@continuedfloat\@firstoftwo\@secondoftwo}
%    \end{macrocode}
%    \begin{macrocode}
\newcommand*\caption@continuedfloat[1]{%
  \caption@iftype
    {#1{\caption@@refstepcounter\@captype}{}%
     \caption@@continuedfloat\@captype}%
    {\caption@OutsideFloat\continuedfloat}}
%    \end{macrocode}
%    \begin{macrocode}
\newcommand*\caption@@continuedfloat{%
  \caption@clrflags
  \caption@set@continuedfloat
    {\caption@setcontinued
     \caption@@@continuedfloat}%
    {\@gobble}}
%    \end{macrocode}
%    \begin{macrocode}
\newcommand*\caption@@@continuedfloat[1]{%
  \caption@setoptions{ContinuedFloat}% 出于兼容性考虑
  \caption@setoptions{continuedfloat}%
  \caption@setoptions{continued#1}%
  \expandafter\l@addto@macro\csname the#1\endcsname\thecontinuedfloat
%    \end{macrocode}
%  使用 \pkg{hyperref} 宏包时，存在如下问题：使用 |\continued|\-|float|
%  会产生重复的超链接——主浮动体与其续接浮动体的 |\@current|\-|Href|
%  相同。因此必须确保每次都生成唯一的标签与引用。做法是扩展
%  |\the|\-|H|\-|figure| 与 |\the|\-|H|\-|table|，使续接浮动体采用
%  如下形式
%  \begin{quote}\meta{类型}|.|\meta{类型编号}|\alph{|\meta{续接编号}|}|\end{quote}
%  而非
%  \begin{quote}\meta{类型}|.|\meta{类型编号}\quad\quad。\end{quote}
%  \par{\small（本实现采用了 Steven Douglas Cochran 的思路。）}
%  \Note{若设置了 \pkg{hyperref} 宏包选项 \texttt{naturalnames=true}，本措施无效。}
%    \begin{macrocode}
  \@ifundefined{theH#1}{}{%
    \expandafter\l@addto@macro\csname theH#1\endcsname{%
      \@alph\c@continuedfloat}}%
%    \end{macrocode}
% 确保选项与 \pkg{hyperref} 修补只应用一次（每组一次）。
%    \begin{macrocode}
  \let\caption@@@continuedfloat\@gobble}
%    \end{macrocode}
% \end{macro}
%
% \begin{macro}{\ContinuedFloat}
% 出于兼容考虑，仍支持已过时的 |\Continued|\-|Float|。
%    \begin{macrocode}
\def\ContinuedFloat{\continuedfloat} % 为兼容而定义
%    \end{macrocode}
% \end{macro}
%
% \begin{macro}{\caption@ContinuedFloat}
% \changes{v3.0f}{2005/06/11}{新增本宏}
% \changes{v3.3}{2013/12/16}{缺陷修复：修订并改用 \cs{def} 而非 \cs{newcommand} 定义}
% |\caption@ContinuedFloat|\marg{题注类型}\par
% 本宏供 \pkg{subfig} 宏包提供的 \cs{ContinuedFloat} 使用。由于
%  \pkg{subfig} 自行处理计数器，我们尽量不通过更改计数器相关标志加以
%  干预。
%    \begin{macrocode}
\def\caption@ContinuedFloat{%
  \caption@set@continuedfloat\caption@@@continuedfloat\@gobble}
%    \end{macrocode}
% \end{macro}
%
% \begin{macro}{\nextfloat}
% \changes{v3.6}{2020/09/20}{新增本命令}
% \changes{v3.6}{2021/01/09}{加入 \cs{flushsubcaptionlistentries} 的使用}
% \changes{v3.6}{2022/01/05}{\cs{caption@Error} 的使用改为 \cs{caption@OutsideFloat}}
% |\nextfloat| 重置标志（表示浮动环境的下一内容从此处开始），并设置新的
%  \pkg{hyperref} 锚点。
%    \begin{macrocode}
\providecommand*\nextfloat{%
    \caption@teststar\caption@nextfloat\@firstoftwo\@secondoftwo}
%    \end{macrocode}
%    \begin{macrocode}
\newcommand*\caption@nextfloat[1]{%
  \caption@iftype
    {\caption@@nextfloat\@captype
     #1{}{\caption@reset@currentlabel\caption@start}}%
    {\caption@OutsideFloat\nextfloat}}
%    \end{macrocode}
%    \begin{macrocode}
\newcommand*\caption@@nextfloat[1]{%
  \flushsubcaptionlistentries
  \caption@clrflags
  \ifx\caption@@@continuedfloat\@gobble
    % 下一个续接浮动体
    \caption@set@continuedfloat\caption@setcontinued\relax
  \fi}
%    \end{macrocode}
% \end{macro}
%
% \subsection{内部辅助宏}
%
% \begin{macro}{\caption@refstepcounter}
% \changes{v3.1}{2007/07/28}{新增本宏}
% \changes{v3.2}{2010/10/26}{错误消息从 \cs{caption@@@settype} 移至此处}
% \changes{v3.2c}{2011/09/04}{加入对 \textsf{memoir} 文档类 \cs{donemaincaptiontrue} 的支持}
% \changes{v3.3}{2011/12/28}{加入 \cs{ifcaption@caption}、\cs{ifcaption@subcaption} 与 \cs{ifcaption@ContinuedFloat} 的使用}
% \changes{v3.5e}{2020/09/28}{修复 \cs{stepcounter} 重定义中的缺陷；现使用 \cs{edef} 而非 \cs{def}}
% \changes{v3.6}{2020/09/20}{标志处理抽出到 \cs{caption@ifrefstepcounter}}
% \changes{v3.6}{2020/12/27}{\cs{caption@@@refstepcounter} 改名为 \cs{caption@refstepcounter@}}
% \changes{v3.6}{2021/01/09}{适配 \cs{caption@prepare@stepcounter} 的当前接口}
% |\caption@refstepcounter|\marg{计数器}\\
% 递增浮动体计数器（使用 |\ref|\-|step|\-|counter|），并重置续接计数器。
%    \begin{macrocode}
\newcommand*\caption@refstepcounter{%
  \@nameuse{donemaincaptiontrue}% 支持 memoir 文档类
  \caption@ifrefstepcounter\caption@@refstepcounter\caption@@refcounter}
%    \end{macrocode}
%    \begin{macrocode}
\newcommand*\caption@@refstepcounter[1]{%
  \caption@prepare@stepcounter{#1}%
  \caption@refstepcounter@{#1}}
%    \end{macrocode}
%    \begin{macrocode}
\newcommand*\caption@@refcounter[1]{%
  \let\caption@stepcounter@ORI\stepcounter
  \def\stepcounter##1{%
    \edef\caption@tempa{#1}%
    \edef\caption@tempb{##1}%
    \ifx\caption@tempa\caption@tempb
      % Suppress incrementation of counter
    \else
      \caption@stepcounter@ORI{##1}%
    \fi}%
  \caption@refstepcounter@{#1}%
  \let\stepcounter\caption@stepcounter@ORI}
%    \end{macrocode}
%    \begin{macrocode}
\newcommand*\caption@refstepcounter@{\refstepcounter}
%    \end{macrocode}
% \end{macro}
%
% \begin{macro}{\caption@stepcounter}
% \changes{v3.6}{2020/09/20}{新增本宏}
% \changes{v3.6}{2020/12/27}{标志处理抽出到 \cs{caption@ifstepcounter}}
% \changes{v3.6}{2020/12/27}{\cs{caption@@@stepcounter} 改名为 \cs{caption@stepcounter@}}
% \changes{v3.6}{2021/01/09}{适配 \cs{caption@prepare@stepcounter} 的当前接口}
% |\caption@stepcounter|\marg{计数器}\\
% 递增浮动体计数器（使用 |\step|\-|counter|），并重置续接计数器。
%    \begin{macrocode}
\newcommand*\caption@stepcounter{%
  \caption@ifstepcounter\caption@@stepcounter\caption@@counter}
%    \end{macrocode}
%    \begin{macrocode}
\newcommand*\caption@@stepcounter[1]{%
  \caption@prepare@stepcounter{#1}%
  \caption@stepcounter@{#1}}
%    \end{macrocode}
%    \begin{macrocode}
\newcommand*\caption@@counter[1]{} % nothing to do
%    \end{macrocode}
%    \begin{macrocode}
\newcommand*\caption@stepcounter@{\stepcounter}
%    \end{macrocode}
% \end{macro}
%
% \begin{macro}{\caption@prepare@stepcounter}
% \changes{v3.3}{2011/12/28}{新增本宏}
% \changes{v3.6}{2021/01/09}{去掉参数 \meta{来源提示}；适配 \cs{caption@addsubcontentslines} 的当前接口}
% |\caption@prepare@stepcounter|\marg{计数器}
% 输出缓存的目录条目并重置续接计数器。
%    \begin{macrocode}
\newcommand*\caption@prepare@stepcounter{%
  \caption@addsubcontentslines\relax
  \caption@reset@continuedfloat}
%    \end{macrocode}
% \end{macro}
%
% \begin{macro}{\caption@ifrefstepcounter}
% \changes{v3.6}{2020/09/20}{新增本宏}
% |\caption@ifrefstepcounter|\marg{是代码}\marg{否代码}\\
% 检测是否应递增计数器，并相应更新标志。
%    \begin{macrocode}
\newcommand*\caption@ifrefstepcounter{%
  \caption@ifcounter
%    \end{macrocode}
% 若本环境中已排版过一则题注，则需要（再次）递增计数器。
%    \begin{macrocode}
    {\caption@ifcaption}%
%    \end{macrocode}
% 若题注及其下方的子题注均已排版，则额外发出警告。（仅当设置
%  |position=b| 时这样做，因为此情形下子题注编号最可能出错。）
%    \begin{macrocode}
    {\caption@ifsubcaption@below{\caption@iftop\@secondoftwo\@firstoftwo}{\@secondoftwo}}%
%    \end{macrocode}
% 无论哪种情形，都设置“题注”标志以表示已排版题注。
%    \begin{macrocode}
    {\caption@setcaption}}
%    \end{macrocode}
% \end{macro}
%
% \begin{macro}{\caption@ifstepcounter}
% \changes{v3.6}{2020/12/27}{新增本宏}
% |\caption@ifstepcounter|\marg{是代码}\marg{否代码}\\
% 检测是否应递增计数器，并相应更新标志。
%    \begin{macrocode}
\newcommand*\caption@ifstepcounter{%
  \caption@ifcounter
%    \end{macrocode}
% 若本环境中已排版过一则题注及其上方的子题注，则需要（再次）递增计数器。
%    \begin{macrocode}
    {\caption@ifcaption\caption@ifsubcaption@above\@secondoftwo}%
%    \end{macrocode}
% 在此情形下额外发出警告。（仅当设置 |position=t| 时这样做，因为此情形下
% 子题注编号最可能出错。）
%    \begin{macrocode}
    {\caption@iftop}%
%    \end{macrocode}
% 无论哪种情形，都设置“子题注”标志以表示已排版子题注。
%    \begin{macrocode}
    {\caption@setsubcaption}}
%    \end{macrocode}
% \end{macro}
%
% \begin{macro}{\caption@ifcounter}
% \changes{v3.6}{2020/12/28}{新增本宏}
% |\caption@ifcounter|\marg{是否递增计数器}\marg{是否发出警告}\marg{设置标志}\marg{是代码}\marg{否代码}\\
% 内含 |\caption@if|\-|ref|\-|step|\-|counter| 与 |\caption@if|\-|step|\-|counter|
%  的公共代码。
%    \begin{macrocode}
\newcommand*\caption@ifcounter[3]{%
  \caption@ifflags
%    \end{macrocode}
% 检测是否应（再次）递增计数器；若是，再检测是否应额外发出警告。
% （计数器递增后，所有标志将被清除，以表示新一轮开始。）
%    \begin{macrocode}
    {#1%
       {#2\caption@Warning\@gobble{Ambiguous sub-caption(s), use \string\nextfloat}%
        \caption@clrflags
        \let\caption@next\@firstoftwo}%
       {\let\caption@next\@secondoftwo}}%
%    \end{macrocode}
% 若此前未设置任何标志，则始终递增计数器。
%    \begin{macrocode}
    {\let\caption@next\@firstoftwo}%
%    \end{macrocode}
% 设置“题注”或“子题注”标志之一。
%    \begin{macrocode}
  #3%
%    \end{macrocode}
%    \begin{macrocode}
  \caption@next}
%    \end{macrocode}
% \end{macro}
%
% \begin{macro}{\caption@dblarg}
% \changes{v3.1}{2007/02/05}{新增本宏}
% \changes{v3.3}{2012/03/25}{加入对选项 \opt{list-entry} 的支持}
% \changes{v3.6n}{2023/07/31}{仅空目录条目展开为 \cs{relax}，其余保持不变}
%  与 |\@dblarg| 相比加入了 |\relax|，使 |\caption{}| 展开为
%  |\caption[\relax]{}|（而非 |\caption[]{}|）。此外还加入了对选项
%  \opt{list-entry} 的支持。
%    \begin{macrocode}
\newcommand\caption@dblarg[1]{%
  \kernel@ifnextchar[%]
    {\caption@ydblarg{#1}}%
    {\caption@xdblarg{#1}}}
%    \end{macrocode}
%    \begin{macrocode}
\newcommand\caption@xdblarg[2]{%
  \caption@iftokens{#2}{#1[{#2}]{#2}}{#1[\relax]{}}}
%    \end{macrocode}
%    \begin{macrocode}
\long\def\caption@ydblarg#1[#2]#3{%
  \caption@iflistheading{#1[{#3}]{#3}}{#1[{#2}]{#3}}}
%    \end{macrocode}
%    \begin{macrocode}
%    \end{macrocode}
% \end{macro}
%
% \begin{macro}{\caption@begin}
% \changes{v3.0a}{2004/01/23}{\cs{caption@settype} 改为 \cs{caption@setfloattype}}
% \changes{v3.0e}{2005/04/12}{加入对默认标签格式的检查}
% \changes{v3.0f}{2005/06/11}{加入 \cs{caption@resetContinuedFloat} 的使用}
% \changes{v3.1}{2007/05/09}{移除 \cs{caption@setfloattype} 与 \cs{caption@resetContinuedFloat} 的使用}
% \changes{v3.3}{2011/12/29}{移除 \cs{caption@fixedposition} 的使用}
%  对 |\caption| 的处理始终由 |\caption@begin|（或 |\caption@beginex|）与
%  |\caption@end| 包裹。\par
%  |\caption@begin|\marg{类型} 执行以下任务：
%  \begin{enumerate}
%  \item[1.] 必要时覆盖 |position=| 设置（例如设为 |auto|，或在
%            |supertabular| 内部使用时）。
%  \item[2.] 开启一个新组。
%  \item[3.] 若题注标签格式非默认，则定义 |\fnum@|\meta{类型}。
%  \end{enumerate}
%    \begin{macrocode}
\newcommand*\caption@begin[1]{%
  \caption@fixposition
  \begingroup
    \caption@setfnum{#1}}
%    \end{macrocode}
% \end{macro}
%
% \begin{macro}{\caption@beginex}
% \changes{v3.0b}{2004/05/16}{加入对 \opt{listof=} 的支持}
% \changes{v3.0l}{2007/02/04}{参数由 2 个增至 3 个，加入对空标题的检查}
% \changes{v3.1}{2007/07/01}{现生成目录条目，而非重定义 \cs{addcontentsline}}
% \changes{v3.2}{2010/10/26}{缺陷修复 09-05-15：加入 \cs{lst@@caption} 的设置}
% \changes{v3.5h}{2022/03/01}{加入钩子，移除 \cs{lst@@caption} 的设置}
%  |\caption@beginex|\marg{类型}\marg{目录条目}\marg{标题}\par
%  除执行 |\caption@begin| 的相同任务外，还会：
%  \begin{enumerate}
%  \item[4.] 应用钩子 |\caption@beginex@hook|
%  \item[5.] 在相应目录中加入条目。
%  \item[6.] 根据参数 \meta{标题} 设置 |\caption@ifempty|。
% \end{enumerate}
%    \begin{macrocode}
\newcommand\caption@beginex[3]{%
  \caption@begin{#1}%
  \caption@beginex@hook
  \caption@addcontentsline{#1}{#2}%
  \caption@ifempty{#3}{}}
%    \end{macrocode}
%    \begin{macrocode}
\newcommand*\caption@beginex@hook{}
%    \end{macrocode}
% \end{macro}
%
% \begin{macro}{\caption@end}
% \changes{v3.3}{2011/12/29}{移除 \cs{caption@fixedposition} 的使用}
% |\caption@end| 关闭该组。
%    \begin{macrocode}
\newcommand*\caption@end{%
  \endgroup}
%    \end{macrocode}
% \end{macro}
%
% \changes{v3.0a}{2004/01/23}{加入内部钩子 \cs{caption@@begin} 与 \cs{caption@@end}}
% \changes{v3.0h}{2005/10/06}{移除内部钩子 \cs{caption@@begin} 与 \cs{caption@@end}}
%
% \begin{macro}{\caption@normalsize}
% \changes{v3.1}{2007/05/05}{新增本宏定义}
% \changes{v3.5}{2020/07/27}{适配 \pkg{caption3} v2.0}
%  |\@caption| 将用它代替 |\normalsize|。
%    \begin{macrocode}
\newcommand*\caption@normalsize{%
  \caption@font@normal\@unused}
%    \end{macrocode}
% \end{macro}
%
% \begin{macro}{\caption@setfloatcapt}
%  为支持 \pkg{float} 宏包而需要：该宏包中题注不直接排版，而是被收进
%  名为 |\@floatcapt| 的 |\vbox| 中。
%    \begin{macrocode}
\let\caption@setfloatcapt\@firstofone
%    \end{macrocode}
% \end{macro}
%
% \begin{macro}{\caption@makecurrent}
%  为支持 \pkg{nameref} 与 \pkg{hyperref} 宏包而需要。
% \changes{v3.2f}{2012/01/14}{缺陷修复：加入 \cs{caption@gettitle} 的使用}
%    \begin{macrocode}
\newcommand*\caption@makecurrent[1]{\caption@gettitle}
%    \end{macrocode}
% \end{macro}
%
% \begin{macro}{\caption@makeanchor}
% \begin{macro}{\caption@start}
% \begin{macro}{\caption@@start}
% \begin{macro}{\caption@freezeHref}
% \begin{macro}{\caption@defrostHref}
%  这些都为支持 \pkg{hyperref} 宏包而需要。
%    \begin{macrocode}
\let\caption@makeanchor\@firstofone
\let\caption@start\relax
\let\caption@@start\relax
\let\caption@freezeHref\relax
\let\caption@defrostHref\relax
%    \end{macrocode}
% \end{macro}
% \end{macro}
% \end{macro}
% \end{macro}
% \end{macro}
%
% \begin{macro}{\caption@gettitle}
% \changes{v3.1l}{2010/01/09}{新增本宏及其使用，使 \cs{@currentlabelname} 的定义适配较新版本的 \pkg{nameref}}
% \changes{v3.6o}{2023/08/05}{修订对旧版 \pkg{nameref} 的支持，使其与 \LaTeX\ 2023-06-01 兼容}
% 为支持 \pkg{nameref} 宏包而需要。
%    \begin{macrocode}
\newcommand\caption@gettitle[1]{%
  \caption@ifdefined\NR@gettitle
    {\NR@gettitle{#1}}%                                         % 支持较新版本的 nameref 宏包
    {\caption@ifdefined\NR@sect{\def\@currentlabelname{#1}}{}}} % 支持旧版本的 nameref 宏包
%    \end{macrocode}
% \end{macro}
%
% \subsection{对子题注的支持}
% \changes{v3.1f}{2007/11/16}{加入对子题注的支持}
%
% \begin{macro}{\subcaptionsetup}
% \changes{v3.6l}{2023/03/22}{本宏定义自 \pkg{subcaption} 宏包移至此处}
% |\subcaptionsetup*|\oarg{环境}\marg{选项}
% 是 |\caption|\-|setup*[sub|\meta{环境}|]|\marg{选项} 的简写；尤其
%  |\sub|\-|caption|\-|setup|\marg{选项} 所设选项对所有子题注均有效。
%    \begin{macrocode}
\newcommand*\subcaptionsetup{%
  \caption@teststar\@subcaptionsetup*{}}
%    \end{macrocode}
%    \begin{macrocode}
\newcommand*\@subcaptionsetup[1]{%
  \@testopt{\captionsetup@sub{#1}}{}}
%    \end{macrocode}
%    \begin{macrocode}
\def\captionsetup@sub#1[#2]#3{%
  \captionsetup#1[sub#2]{#3}}
%    \end{macrocode}
% \end{macro}
%
% \begin{macro}{\caption@subtypehook}
% \changes{v3.1g}{2008/03/20}{加入对 \cs{caption@setlist} 的重定义}
% \changes{v3.1k}{2009/10/09}{缺陷修复 09-08-12b：现在也在此恢复 \cs{@makecaption}}
% \changes{v3.2}{2011/06/26}{现使用 \cs{caption@@addcontentsline} 而非 \cs{addcontentsline}}
% \changes{v3.2a}{2011/08/17}{加入对 \cs{phantomcaption} 的重定义}
% \changes{v3.2a}{2011/08/17}{加入 \cs{caption@warmup} 的使用}
% \changes{v3.2c}{2011/09/30}{缺陷修复 11-09-26：加入 \cs{if@subfloatrow} 的使用}
% \changes{v3.3}{2011/12/28}{缺陷修复 11-12-16：\cs{caption@l@stepcounter} 由 \cs{stepcounter} 取代}
% \changes{v3.3}{2012/03/15}{加入 \cs{caption@setbox} 的使用}
% \changes{v3.3}{2013/05/01}{加入对 \cs{caption@settype} 的重定义}
% \changes{v3.3}{2013/06/17}{补上遗漏的对 \cs{captionlistentry} 的重定义}
% \changes{v3.3}{2018/09/06}{适配 \pkg{chkfloat} 宏包}
% \changes{v3.4f}{2020/01/03}{加入对 \cs{@caption} 的重定义，因其被 \env{threeparttable} 环境重定义}
% \changes{v3.5}{2020/07/27}{适配 \pkg{caption3} v2.0}
% \changes{v3.6}{2020/09/20}{标志处理抽出到 \cs{caption@stepcounter}}
% \changes{v3.6d}{2022/04/17}{移除 \cs{if@subfloatrow} 的使用，改为重定义 \env{subfloatrow*}}
% 钩子，供 \cs{caption@setsubtype} 内部使用。\par
% （注意：若位于 |sub|\-|float|\-|row| 环境内部，必须保持 \pkg{floatrow}
%  宏包的 |\@make|\-|caption| 代码完好。）
%    \begin{macrocode}
\newcommand*\caption@subtypehook{%
  \ifx\caption\caption@subcaption \else
    \caption@warmup
    \caption@stepcounter\@captype
    \caption@ifcaption
      {}%
      {\let\caption@add@contentsline\caption@addsubcontentsline
       \let\caption@addsubcontentslines\@gobble}%
%    \end{macrocode}
%    \begin{macrocode}
    \c@continuedfloat=0\relax
    \let\caption@setfloatcapt\@firstofone
    \let\caption@chkfloat\@gobbletwo
%    \end{macrocode}
%    \begin{macrocode}
    \caption@set{box}{none}%
    \caption@clearmargin
    \caption@iflist{}{\let\caption@setlist\@gobble}%
    \caption@setoptions{sub}%
    \caption@setoptions{subfloat}% 为与 subfig 宏包兼容
%    \end{macrocode}
%    \begin{macrocode}
% 重定义 \setcaptiontype
    \renewcommand*\caption@settype[1]{%
      \caption@Error{##1 inside \@subcaptype}}%
%    \end{macrocode}
%    \begin{macrocode}
% 重定义 \caption
    \let\caption\caption@subcaption
    \let\@caption\caption@@caption
    \let\phantomcaption\caption@subphantom
    \let\captionlistentry\caption@sublistentry
    \let\captiontext\caption@subtext
%   \let\@captype\@subcaptype
    \let\caption@refstepcounter\caption@refstepcounter@
%    \end{macrocode}
%    \begin{macrocode}
% 恢复 \@makecaption
    \let\@makecaption\caption@makecaption
  \fi}%
%    \end{macrocode}
% \end{macro}
%
% \begin{macro}{\caption@subcaption}
% 生成一则子题注。
%    \begin{macrocode}
\newcommand*\caption@subcaption{%
  \caption@checkgrouplevel{sub}\subcaption
  \caption@star
    {\caption@refstepcounter\@subcaptype}%
    {\caption@dblarg{\@caption\@subcaptype}}}
%    \end{macrocode}
% \end{macro}
%
% \begin{macro}{\caption@subphantom}
% \changes{v3.2a}{2011/08/17}{新增本宏}
% 与 |\phantomcaption| 相同，但用于子题注。
%    \begin{macrocode}
\newcommand*\caption@subphantom{%
  \caption@checkgrouplevel{sub}\phantomsubcaption
  \caption@refstepcounter\@subcaptype}
%    \end{macrocode}
% \end{macro}
%
% \begin{macro}{\caption@sublistentry}
% \changes{v3.3}{2013/06/17}{新增本宏}
% 与 |\captionlistentry| 相同，但用于子题注。
%    \begin{macrocode}
\newcommand*\caption@sublistentry{%
  \caption@teststar\@subcaptionlistentry\@firstoftwo\@secondoftwo}
%    \end{macrocode}
%    \begin{macrocode}
\newcommand*\@subcaptionlistentry[1]{%
  \@testopt{\caption@listentry{#1}}{\@subcaptype}}
%    \end{macrocode}
% \end{macro}
%
% \begin{macro}{\caption@subtext}
% \changes{v3.2a}{2011/08/17}{新增本宏}
% 与 |\captiontext| 相同，但用于子题注。
%    \begin{macrocode}
\newcommand*\caption@subtext{%
% \caption@checkgrouplevel{sub}\subcaptiontext
  \caption@text\@subcaptype}
%    \end{macrocode}
% \end{macro}
%
% \begin{macro}{caption@clearsubcontentslines}
% 清除待处理的子题注目录条目。
%    \begin{macrocode}
\newcommand*\caption@clearsubcontentslines{%
  \global\let\caption@subcontentslines\@empty}
%    \end{macrocode}
%    \begin{macrocode}
\caption@clearsubcontentslines
%    \end{macrocode}
% \end{macro}
%
% \begin{macro}{\caption@addsubcontentsline}
% \changes{v3.1i}{2008/08/02}{目录条目内现允许 \cs{label}、\cs{index} 与 \cs{glossary}}
% \changes{v3.2}{2011/06/26}{现使用 \cs{caption@@addcontentsline} 而非 \cs{addcontentsline}}
% \changes{v3.2c}{2011/08/19}{现使用 \cs{caption@gobble} 而非 \cs{@gobble}}
% \changes{v3.3}{2012/01/15}{加入调试消息}
% \changes{v3.3}{2013/03/10}{单个 \cs{protect} 由 \cs{noexpand} 取代}
% \changes{v3.6m}{2023/07/28}{对 \cs{label} 等的重定义由 \cs{caption@prepareslc} 取代}
% 加入一条待处理的子题注目录条目。
%    \begin{macrocode}
\newcommand*\caption@addsubcontentsline[4]{%
  \caption@Debug{\string\caption@addsubcontentsline{#1}{#2}}%
  \begingroup
  \caption@prepareslc
  \protected@edef\@tempa{\endgroup
    \noexpand\g@addto@macro\noexpand\caption@subcontentslines{%
      \noexpand\@namedef{the#2}{\csname the#2\endcsname}%
      \ifx\@currentHref\@undefined \else
        \noexpand\def\noexpand\@currentHref{\@currentHref}%
      \fi
      \noexpand\caption@@@addcontentsline{#1}{#2}{#3}{#4}}}%
  \@tempa}
%    \end{macrocode}
% \end{macro}
%
% \begin{macro}{\flushsubcaptionlistentries}
% \changes{v3.3}{2012/01/15}{由 \cs{caption@addsubcontentslines} 改名为 \cs{flushsubcaptionlistentries}}
% \changes{v3.3}{2013/03/10}{缺陷修复：补上遗漏的对 \cs{caption@addsubcontentslines} 的重定义}
% \changes{v3.6}{2021/01/09}{弃用对 \cs{chapter} 与 \cs{appendix} 的修补}
% 写出待处理的子题注目录条目。
% （注意：出于兼容考虑，保留 |\caption@add|\-|sub|\-|contents|\-|lines|
%  的参数，尽管它已不再使用。）
%    \begin{macrocode}
\newcommand*\flushsubcaptionlistentries{%
  \caption@addsubcontentslines\relax}
%    \end{macrocode}
%    \begin{macrocode}
\renewcommand*\caption@addsubcontentslines[1]{%
  \caption@Debug{\string\flushsubcaptionlistentries}%
  \begingroup
    \caption@subcontentslines
  \endgroup
  \caption@clearsubcontentslines}
%    \end{macrocode}
% \end{macro}
%
% \subsection{Babel 宏包支持}
%
% \subsubsection{arabic 与 farsi babel 选项}
% \changes{v3.5}{2020/08/02}{加入 \pkg{arabic} babel 支持}
%
%    \begin{macrocode}
\caption@IfBabelArabi{%
%    \end{macrocode}
%
% \pkg{arabicore} 宏包以与 \pkg{caption} 宏包不兼容的方式重定义
%  |\@make|\-|caption|、|\fnum@figure| 与 |\fnum@table|，因此必须在此
%  恢复本宏包或原始定义。
%
%    \begin{macrocode}
  \let\@makecaption\caption@makecaption
  \def\fnum@figure{\figurename\nobreakspace\thefigure}%
  \def\fnum@table{\tablename\nobreakspace\thetable}%
%    \end{macrocode}
%
%    \begin{macrocode}
}{}
%    \end{macrocode}
%
% \subsubsection{frenchle 与 frenchpro 宏包}
% \changes{v3.0h}{2005/10/03}{加入 \pkg{frenchle/pro} 宏包支持}
% \changes{v3.5}{2020/08/01}{更新 \pkg{frenchle/pro} 宏包支持}
%
%    \begin{macrocode}
\caption@IfFrenchLe{%
%    \end{macrocode}
%
% \begin{macro}{\f@ffrench}
% \begin{macro}{\f@tfrench}
%  |\f@ffrench| 与 |\f@tfrench| 在法语模式下对应 |\fnum@figure| 与
%  |\fnum@table|。二者含有附加代码，用于排版题注分隔符
%  |\caption|\-|separator|，而非通常的冒号。由于这会与本宏包的
%  |\@make|\-|caption| 代码冲突，必须在此移除该附加代码。
%    \begin{macrocode}
  \def\f@ffrench{\ifx\listoffigures\relax\else\figurename~\thefigure\fi}%
  \def\f@tfrench{\ifx\listoftables\relax\else\tablename~\thetable\fi}%
  \let\@eatDP\@undefined
%    \end{macrocode}
% \end{macro}
% \end{macro}
%
% 由于 \pkg{frenchle} 宏包（在 |\GO|\-|french| 内部）覆盖
%  |\@make|\-|caption|，需要事先保存、事后恢复。
%    \begin{macrocode}
  \caption@AtBeginDocument{%
    \let\caption@original@makecaption\@makecaption}%
  \g@addto@macro\GOfrench{%
    \let\@makecaption\caption@original@makecaption}%
%    \end{macrocode}
%
%    \begin{macrocode}
}{}
%    \end{macrocode}
%
% \subsubsection{hungarian 与 magyar babel 选项}
% \changes{v3.2}{2009/03/29}{加入 \pkg{magyar} babel 支持}
% \changes{v3.3}{2018/08/26}{修订 \pkg{magyar} babel 支持}
% \changes{v3.5}{2020/08/02}{更新 \pkg{magyar} babel 支持}
%
%    \begin{macrocode}
\caption@IfBabelHungarian{%
%    \end{macrocode}
%
% 由于 \pkg{magyar} 宏包倾向于覆盖 |\@make|\-|caption|，需事先保存、
% 事后恢复。
% \Note{目前它不会被覆盖，这是 \pkg{caption3} 内为 \pkg{floatrow} 宏包
%       所做变通的副作用。}
%
%    \begin{macrocode}
  \caption@AtBeginDocument{%
    \let\caption@original@makecaption\@makecaption}%
  \caption@ifdefined\extrashungarian{%
    \addto\extrashungarian{%
      \let\@makecaption\caption@original@makecaption}}{}%
  \caption@ifdefined\extrasmagyar{%
    \addto\extrasmagyar{%
      \let\@makecaption\caption@original@makecaption}}{}%
%    \end{macrocode}
%
%    \begin{macrocode}
}{}
%    \end{macrocode}
%
% \subsubsection{latvian babel 选项}
% \changes{v3.5}{2020/08/02}{更新 \pkg{latvian} babel 支持}
%
%    \begin{macrocode}
\caption@IfBabelLatvian{%
%    \end{macrocode}
%
% \pkg{latvian} 宏包以与 \pkg{caption} 宏包不兼容的方式重定义
%  |\@make|\-|caption|，因此必须在此恢复本宏包定义。
%
%    \begin{macrocode}
  \let\@makecaption\caption@makecaption
%    \end{macrocode}
%
%    \begin{macrocode}
}{}
%    \end{macrocode}
%
% \subsection{宏包支持}
% \changes{v3.1}{2007/04/11}{不再支持 `float'、`longtable' 等宏包选项，已移除}
%
% \begin{macro}{\caption@IfPackageLoaded}
% \changes{v3.0f}{2005/05/29}{加入第三个参数 \meta{代码}，无需额外检查}
% \changes{v3.1}{2007/07/04}{重写，并由 \cs{caption@ifpackage} 改名为 \cs{caption@IfPackageLoaded}}
% \changes{v3.3}{2018/12/26}{缺陷修复：补上遗漏的花括号}
%  |\caption@IfPackageLoaded|\marg{宏包}\oarg{版本}\marg{是代码}\marg{否代码}\par
%  某种程度上是 |\@ifpackageloaded| 与 |\@ifpackagelater| 的组合。若
%  \meta{宏包} 尚未加载，检查将在 |\AtBeginDocument|（再次）执行，因此该
%  宏包也可在稍后加载。
%    \begin{macrocode}
\newcommand\caption@IfPackageLoaded[1]{%
  \@testopt{\caption@@IfPackageLoaded{#1}}{}}
\@onlypreamble\caption@IfPackageLoaded
%    \end{macrocode}
%    \begin{macrocode}
\long\def\caption@@IfPackageLoaded#1[#2]#3#4{%
  \@ifpackageloaded{#1}\@firstofone{%
    \caption@Debug{#1 package is not loaded (yet)\@gobble}%
    \caption@AtBeginDocument}{%
      \caption@If@Package@Loaded{#1}[{#2}]{#3}{#4}}}
\@onlypreamble\caption@@IfPackageLoaded
%    \end{macrocode}
%    \begin{macrocode}
\long\def\caption@If@Package@Loaded#1[#2]{%
  \@ifpackageloaded{#1}{%
    \caption@InfoNoLine{#1 package is loaded}%
    \@ifpackagelater{#1}{#2}\@firstoftwo{%
      \caption@Error{%
        For a successful cooperation we need at least version\MessageBreak
          `#2' of package #1,\MessageBreak
        but only version\MessageBreak
          `\csname ver@#1.\@pkgextension\endcsname'\MessageBreak
        is available}%
      \@secondoftwo}%
  }{\@secondoftwo}}
\@onlypreamble\caption@If@Package@Loaded
%    \end{macrocode}
% \end{macro}
%
% \begin{macro}{\caption@clearmargin}
% \changes{v3.1}{2007/05/05}{新增本宏定义}
% \changes{v3.6}{2021/01/17}{适配 \cs{caption@setmargin} 的当前实现}
% 本宏供某些不需要常规边距设置的宏包支持代码使用，例如 \pkg{sidecap}
% 宏包。
%    \begin{macrocode}
\newcommand*\caption@clearmargin{%
  \caption@setmargin\z@
  \let\caption@minmargin\@undefined}
%    \end{macrocode}
% \end{macro}
%
%    \begin{macrocode}
\caption@setbool{needfreeze}{0}
\caption@AtBeginDocument*{%
  \caption@ifneedfreeze{%
%    \end{macrocode}
%
% \begin{macro}{\caption@freezetype}
% \changes{v3.3}{2013/05/01}{新增本宏}
%  |\caption@freezetype|\marg{类型}\par
%  供 \pkg{fltpage} 与 \pkg{sidecap} 宏包支持使用。
%    \begin{macrocode}
    \newcommand*\caption@freezetype[1]{%
      \caption@settype{#1}%
      \captionsetup*[sub]{hypcap=true}% 注意：这只是一个临时（q&d）变通办法！
      \caption@freeze}%
%    \end{macrocode}
% \end{macro}
%
% \begin{macro}{\caption@freeze}
% \changes{v3.1}{2007/07/22}{新增本宏}
% \changes{v3.1c}{2007/10/06}{修订 \cs{caption} 与 \cs{@caption} 以增强兼容性}
% \changes{v3.1c}{2007/10/06}{缺陷修复：将 \cs{@caption} 改为 long 命令}
% \changes{v3.1f}{2007/11/16}{现使用“空格技巧”而非 \cs{ignorespaces}}
% \changes{v3.1g}{2008/01/29}{新增 \cs{ContinuedFloat} 与 \cs{captionsetup} 的带星号形式}
% \changes{v3.2}{2010/10/26}{缺陷修复 10-09-01a：重新使用 \cs{ignorespaces} 而非“空格技巧”}
% \changes{v3.2a}{2011/08/17}{移除 \cs{caption@freeze} 的带星号形式}
% \changes{v3.2b}{2011/08/18}{现重定义 \cs{caption@setup} 而非 \cs{captionsetup}}
% \changes{v3.2b}{2011/08/18}{修订对 \cs{label} 的重定义}
% \changes{v3.2c}{2011/08/21}{修订}
% \changes{v3.3}{2013/01/06}{适配新开关 \cs{caption@caption} 等}
% |\caption@freeze|\par
% 供 \pkg{bicaption} 宏包使用。
%    \begin{macrocode}
    \newcommand*\caption@freeze{%
      \let\caption@frozen@continuedfloat\continuedfloat
      \def\continuedfloat{%
        \caption@@freeze{\caption@@@continuedfloat\@captype}%
        \caption@frozen@continuedfloat}%
%    \end{macrocode}
%    \begin{macrocode}
      \let\caption@frozen@nextfloat\nextfloat
      \def\nextfloat{\caption@freeze@error}%
%    \end{macrocode}
%    \begin{macrocode}
      \let\caption@frozen@setup\caption@setup
      \def\caption@setup##1{%
        \caption@@freeze{\caption@setup{##1}}%
        \caption@frozen@setup{##1}}%
%    \end{macrocode}
%    \begin{macrocode}
      \let\caption@frozen@caption\caption
      \def\caption{%
        \def\caption{\caption@freeze@error\caption@gobble}%
        \@ifstar
          {\caption@SC@caption*}%
          {\let\@currentlabel\caption@SClabel
           \caption@withoptargs\caption@SC@caption}}%
      \long\def\caption@SC@caption##1##2{%
        \caption@@freeze{\caption##1{##2}}%
        \ignorespaces}%
%    \end{macrocode}
%    \begin{macrocode}
      \let\caption@frozen@label\label
      \def\label{%
        \caption@withoptargs\caption@SC@label}%
      \def\caption@SC@label##1##2{%
        \ifx\@currentlabel\caption@SClabel
          \@bsphack
            \caption@freeze@label{##1}{##2}%
          \@esphack
        \else
          \caption@frozen@label##1{##2}%
        \fi}%
      \def\caption@SClabel{\caption@undefinedlabel}%
      \def\caption@freeze@label##1##2{%
        \caption@@freeze{\label##1{##2}}}%
%    \end{macrocode}
%    \begin{macrocode}
      \global\let\caption@frozen@content\@empty
      \long\def\caption@@freeze{%
        \g@addto@macro\caption@frozen@content}%
%    \end{macrocode}
%    \begin{macrocode}
      \def\caption@warmup{%
        \let\continuedfloat\caption@frozen@continuedfloat
        \let\nextfloat\caption@frozen@nextfloat
        \let\caption@setup\caption@frozen@setup
        \let\caption\caption@frozen@caption
        \let\label\caption@frozen@label}}%
%    \end{macrocode}
%    \begin{macrocode}
    \newcommand*\caption@freeze@error{%
      \caption@Error{Only one \noexpand\caption can be placed in this environment}}
%    \end{macrocode}
% \end{macro}
%
% \begin{macro}{\caption@prepare@defrost}
% \changes{v3.3}{2013/01/06}{新增本宏}
% 防止重置题注标志。
%    \begin{macrocode}
    \newcommand*\caption@prepare@defrost{%
      \def\caption@settype{\caption@@settype{}}}
%    \end{macrocode}
% \end{macro}
%
% \begin{macro}{\caption@defrost}
% \changes{v3.1}{2007/07/22}{新增本宏}
% \changes{v3.1g}{2008/01/30}{加入 \cs{caption@defrost@setup}}
% \changes{v3.2c}{2011/08/21}{适配 \cs{caption@freeze} 的改动}
% \changes{v3.2c}{2011/08/21}{移除 \cs{caption@defrost@setup}}
%  |\caption@defrost|
%    \begin{macrocode}
    \newcommand*\caption@defrost{%
      \ifx\caption@frozen@caption\@undefined
        \caption@frozen@content
      \else
        \caption@Error{Internal Error:\MessageBreak
          \noexpand\caption@defrost in same group as \string\caption@freeze}%
      \fi}%
%    \end{macrocode}
% \end{macro}
%
%    \begin{macrocode}
  }{}%
  \caption@undefbool{needfreeze}}
%    \end{macrocode}
%
% \begin{macro}{\caption@warmup}
% \changes{v3.2a}{2011/08/17}{新增本宏}
%  |\caption@warmup|
%    \begin{macrocode}
\let\caption@warmup\relax
%    \end{macrocode}
% \end{macro}
%
% \subsubsection{changepage 宏包}
% \changes{v3.4k}{2020/05/30}{加入对 \pkg{changepage} 宏包的修补}
%
%    \begin{macrocode}
\caption@IfPackageLoaded{changepage}[2008/07/23 memoir-v3.7k]{%
%    \end{macrocode}
%
% \begin{macro}{\caption@patch@adjustwidth}
% \changes{v3.4k}{2020/05/30}{新增本宏定义}
% \changes{v3.6j}{2023/03/12}{环境重定义中加入 \cs{ignorespaces}}
% \env{adjustwidth} 与 \env{adjustwidth*} 两个环境都开启一个
%  \env{list} \cs{item}。但在 \env{list} 条目内部 \cs{hsize} 未被正确
%  设置，因此文档类提供的 \cs{@makecaption} 在 \env{adjustwidth} 内部
%  使用时不能正常工作。为此我们在 \cs{@makecaption} 内部检测“是否位于
%  \env{adjustwidth} 环境中”这一情形并加以修正。（唉！）
%    \begin{macrocode}
  \providecommand*\caption@patch@adjustwidth[1]{%
    \expandafter\let\csname caption@#1\expandafter\endcsname\csname #1\endcsname
    \expandafter\renewcommand\csname #1\endcsname[2]{%
      \@nameuse{caption@#1}{##1}{##2}%
      \caption@set@adjustwidth
      \ignorespaces}}%
  \caption@patch@adjustwidth{adjustwidth}%
  \caption@patch@adjustwidth{adjustwidth*}%
  \let\caption@patch@adjustwidth\relax
%    \end{macrocode}
% \end{macro}
%
% \begin{macro}{\caption@set@adjustwidth}
% \changes{v3.4k}{2020/05/30}{新增本宏定义}
% \changes{v3.5h}{2022/03/01}{\cs{AtBeginCaption} 的使用改为 \cs{l@addto@macro}\cs{caption@beginex@hook}}
% 此处不能简单设置并检测标志，因为例如在 \env{adjustwidth} 环境内的
%  \cs{parbox} 中不应对 \cs{hsize} 做调整。因此我们保存 \cs{hsize} 与
%  \cs{linewidth}，改为将当前 \cs{hsize}、\cs{linewidth} 与所存值比较。
% 若二者相等，即可（但愿）相当有把握地假定当前位于普通
%  \env{adjustwidth} 环境内部，因而需要在排版题注前修正 \cs{hsize}。
% （长叹！）
%    \begin{macrocode}
  \newdimen\caption@adjustwidth@hsize
  \newdimen\caption@adjustwidth@linewidth
%    \end{macrocode}
%    \begin{macrocode}
  \providecommand*\caption@set@adjustwidth{%
    \caption@adjustwidth@hsize\hsize
    \caption@adjustwidth@linewidth\linewidth
    \l@addto@macro\caption@beginex@hook\caption@test@adjustwidth}%
%    \end{macrocode}
% \end{macro}
%
% \begin{macro}{\caption@test@adjustwidth}
% \changes{v3.4k}{2020/05/30}{新增本宏定义}
% 检测 \cs{hsize} 与 \cs{linewidth} 的值是否与 \env{adjustwidth}
% 环境中所存值相同；若是，则将 \cs{hsize} 设为正确值。
%    \begin{macrocode}
  \providecommand\caption@test@adjustwidth{%
    \ifdim\hsize=\caption@adjustwidth@hsize
      \ifdim\linewidth=\caption@adjustwidth@linewidth
        \hsize=\linewidth
      \fi
    \fi}%
%    \end{macrocode}
% \end{macro}
%
%    \begin{macrocode}
  }{}%
%    \end{macrocode}
%
% \subsubsection{chkfloat 宏包}
% \changes{v3.3}{2018/09/06}{加入对 \pkg{chkfloat} 宏包的支持}
%
%    \begin{macrocode}
\caption@IfPackageLoaded{chkfloat}[2012/08/19 v0.1]{%
%    \end{macrocode}
%
% \begin{macro}{\caption@chkfloat}
% \changes{v3.3}{2018/09/06}{新增本宏}
%    \begin{macrocode}
  \caption@ifdefined\chkfloat@addcontentsline{%
    \renewcommand\caption@chkfloat{\chkfloat@addcontentsline}%
  }{%
    \renewcommand\caption@chkfloat[2]{%
      \caption@ifdefined\chkfloat@page{%
        \addtocontents{fof}{\protect\chkfloat@{\thepage}{\chkfloat@page}{\csname fnum@#1\endcsname: #2}}%
      }{}}%
  }%
%    \end{macrocode}
% \end{macro}
%
%    \begin{macrocode}
  }{}%
%    \end{macrocode}
%
% \subsubsection{float 宏包}
% \changes{v3.0f}{2005/05/31}{重写并改进对 \pkg{float} 宏包的支持}
% \iffalse
% （这是一种更直接的实现，但目前只适用于 \version{1.3}，
%   不再支持 \version{1.2}。）
% \fi
%
% \pkg{float} 宏包通常不使用 \LaTeX\ 内核命令 |\@caption| 排版题注，而是
% 使用 |\float@caption|。（仅当浮动体用 |\restylefloat*| 重新定义样式时
% 才会使用 |\@caption|。）
%
% |\float@caption| 主要有两点不同：
% \begin{itemize}
%   \item 题注排在名为 |\@floatcapt| 的 |\savebox| 内部，以便之后放置到
%         浮动体内容的上方或下方。
%   \item 最终排版题注不使用 |\@makecaption|，而使用 |\@fs@capt|；其定义
%         属于浮动体样式的一部分。（注意：|\@fs@capt| 不会排版题注上下的
%         竖直间距；该间距由浮动体样式代码本身排版。）
% \end{itemize}
%
%    \begin{macrocode}
\caption@IfPackageLoaded{float}[2001/11/08 v1.3d]{%
 \@ifpackageloaded{floatrow}{%
  \caption@If@Package@Loaded{floatrow}[2007/08/24 v0.2a]{}{}%
 }{%
%    \end{macrocode}
%
% \begin{macro}{\@float@setevery}
% \changes{v3.1k}{2009/10/09}{缺陷修复 08-12-05：在 \cs{caption@setfloatcapt} 的重定义中加入 \cs{color@begin/endgroup}}
% \changes{v3.4b}{2019/12/16}{加入 \cs{caption@float@setname} 的使用}
%  |\@float@setevery|\marg{浮动体类型} 由 \pkg{float} 宏包提供；每当用
%  |\newfloat| 或 |\restylefloat| 定义的浮动环境开始时调用。我们利用该
%  钩子做一些适配，并设置合适的题注样式（若已定义）以及通过
%  |\captionsetup|\oarg{浮动体样式} 声明的附加设置。
%    \begin{macrocode}
  \let\caption@ORI@float@setevery\@float@setevery
  \def\@float@setevery#1{%
    \float@ifcaption{#1}{%
%    \end{macrocode}
%
%  首先，通过将 |\@fs@iftopcapt|（浮动体样式的一部分，控制题注排在浮动
%  体内容上方还是下方）转换为本宏包的 |position=| 设置，把题注位置设为
%  正确值。由于题注上下间距由浮动体样式、\emph{而非} 本宏包负责，这看似
%  无用；但实际并非如此，因为基于 \thispackage\ 的某些宏包（如
%  \pkg{subfig}）可能需要该信息，因而会使用本内核提供的
%  |\caption@iftop| 宏。此外，|\captionof| 最终借助 |\@makecaption|
%  排版带间距的题注，本宏包自身也需要该信息。
%    \begin{macrocode}
      \caption@setposition{\@fs@iftopcapt t\else b\fi}%
%    \end{macrocode}
%
%  随后重定义 |\caption@setfloatcapt|（在 |\@caption| 内部使用），使题注
%  设置在盒子 |\@floatcapt| 内部，不带额外竖直间距。
%    \begin{macrocode}
      \renewcommand\caption@setfloatcapt[1]{%
        \let\@makecaption\caption@@make
        \global\setbox\@floatcapt\vbox{%
          \color@begingroup ##1\color@endgroup}}%
%    \end{macrocode}
%
%  为使不同浮动体样式可搭配不同题注样式，还需确定当前浮动体样式（如
%  `ruled'），并选择同名题注样式（及附加设置）——若已定义。
%    \begin{macrocode}
      \float@getstyle\float@style{#1}%
      \caption@setstyle*\float@style
      \caption@setoptions\float@style
%    \end{macrocode}
%
%    \begin{macrocode}
    }{}%
    \caption@float@setname{#1}%
    \caption@freezeHref % 将在 \float@makebox 中解冻
    \caption@ORI@float@setevery{#1}}%
%    \end{macrocode}
% \end{macro}
%
% \begin{macro}{\float@makebox}
% \changes{v3.3}{2013/01/06}{现纠正 hyperref 与 float 的错误加载顺序}
% \changes{v3.3}{2013/02/03}{修订对 hyperref-float 错误加载顺序的纠正}
% 重定义 |\float@makebox|：
%    \begin{macrocode}
  \caption@AtBeginDocument{%
%    \end{macrocode}
% 若该命令已被 \pkg{hyperref} 宏包重定义，需基于原始定义（而非重定义后的
% 版本）构建。
%    \begin{macrocode}
    \caption@ifdefined\HyOrg@float@makebox
      {\let\caption@ORI@float@makebox\HyOrg@float@makebox}%
      {\let\caption@ORI@float@makebox\float@makebox}%
    \renewcommand\float@makebox[1]{%
      \caption@ORI@float@makebox{#1\relax \caption@defrostHref}}%
%    \end{macrocode}
%    \begin{macrocode}
  }%
%    \end{macrocode}
% \end{macro}
%
% \begin{macro}{\caption@float@setname}
% \changes{v3.1}{2007/07/08}{加入 float 宏包钩子}
% \changes{v3.4b}{2019/12/16}{由 \cs{caption@typehook} 改名为 \cs{caption@float@setname}}
% \changes{v3.6i}{2023/02/19}{现使用 \cs{def} 而非 \cs{let}}
%  \LaTeX\ 及几乎所有其他宏包都使用
%    |\|\meta{类型}|name|
%  提供类型或环境名对应的宏；例如命令 |\figure|\-|name| 通常存放浮动
%  环境 |figure| 的名称：
%  \begin{quote}
%    |\newcommand\figurename{Figure}|
%  \end{quote}
%  但 \pkg{float} 宏包不遵循这一通用命名约定：对用 |\new|\-|float|
%  定义的浮动体，它改用 |\fname@|\meta{类型}，这会与本代码（以及
%  |\auto|\-|ref| 等其他功能）冲突。因此此处必须把 \pkg{float} 宏包的
%  名称映射为通用形式。\par
%  \Note{若浮动体并非用 \cs{newfloat}、而是用 \cs{restylefloat} 定义，
%  \cs{fname@}\meta{类型} 并未定义。}
%    \begin{macrocode}
  \newcommand*\caption@float@setname[1]{%
    \expandafter\ifx\csname #1name\endcsname\relax % 未定义或为 \relax
      \@namedef{#1name}{\@nameuse{fname@#1}}%
    \fi}%
%    \end{macrocode}
% \end{macro}
%
% \begin{macro}{\caption@setname}
% \changes{v3.6i}{2023/02/19}{在 \pkg{float} 宏包支持中加入对 \cs{caption@setname} 的重定义}
% 与 |\caption@float@set|\-|name| 存在相同问题，但方向相反：使用选项
%  |name=| 时，也应设置 |\fname@|\meta{类型}。
%    \begin{macrocode}
  \renewcommand*\caption@setname[2]{%
    \@ifundefined{fname@#1}{}{\caption@ifdefined\floatname\floatname\@gobbletwo{#1}{#2}}%
    \@namedef{#1name}{#2}}%
%    \end{macrocode}
% \end{macro}
%
% \changes{v3.0f}{2005/06/26}{纠正 \env{plaintop} 与 \env{boxed} 浮动体的间距}
% \begin{macro}{\fs@plaintop}
% \begin{macro}{\fs@boxed}
%  由于浮动体样式 |plaintop| 与 |boxed| 不使用可由 |skip=| 设置的
%  |\abovecaptionskip|（|plaintop| 改用 |\belowcaptionskip|，|boxed| 使用
%  固定 |2pt| 间距），此处修补相应浮动体样式宏以改变这一点。
%    \begin{macrocode}
  \g@addto@macro\fs@plaintop{\def\@fs@mid{\vspace\abovecaptionskip\relax}}%
  \g@addto@macro\fs@boxed{\def\@fs@mid{\kern\abovecaptionskip\relax}}%
%    \end{macrocode}
% \end{macro}
% \end{macro}
%
% \begin{macro}{\float@getstyle}
% \changes{v3.1}{2007/06/10}{新增本宏}
% |\float@getstyle|\marg{命令}\marg{类型}\par
%  确定浮动体样式并非易事，因为 \pkg{float} 宏包提供的唯一线索是宏
%   |\fst@|\meta{浮动体类型}，它指向代表浮动体样式的宏。例如在
%  \begin{quote}
%    |\floatstyle{ruled}|\\
%    |\newfloat{Program}{tbp}{lop}|
%  \end{quote}
%  |\fst@Program| will be defined as
%  \begin{quote}
%    |\def\fst@Program{\fs@ruled}|\quad.
%  \end{quote}\par
%  做法如下：把 |\fst@|\meta{浮动体类型} 的第一层展开结果做成字符串，
%  从而吞掉前四个记号（即 |\fs@|），只留下浮动体样式的名称。\par
%  \emph{待办：} 此处还需转换 catcode。
%    \begin{macrocode}
  \providecommand*\float@getstyle[2]{%
    \edef#1{%
      \noexpand\expandafter\noexpand\@gobblefour\noexpand\string
        \expandafter\expandafter\expandafter\noexpand
          \csname fst@#2\endcsname}%
    \edef#1{#1}%
    \caption@Debug{floatstyle{#2} = `#1'}}%
%    \end{macrocode}
% \end{macro}
%
% \begin{macro}{\float@ifcaption}
% \changes{v3.1}{2007/06/23}{新增本宏}
% |\float@ifcaption|\marg{类型}\marg{真分支}\marg{假分支}\par
% 此处判断用户使用的是 |\newfloat| 或 |\restylefloat|，还是
%  |\restylefloat*|。方法很简单：若 |\@float@c@|\meta{题注类型} 与
%  |\float@caption| 相同，则使用的是 |\newfloat| 或 |\restylefloat|；
% 否则假定使用了 |\restylefloat*|。（若有宏包重定义了
%  |\float@caption|，该检测会失效，因此只能假定不存在这种情况。）
%    \begin{macrocode}
  \providecommand*\float@ifcaption[1]{%
    \expandafter\ifx\csname @float@c@#1\endcsname\float@caption
      \expandafter\@firstoftwo
    \else
      \expandafter\@secondoftwo
    \fi}%
%    \end{macrocode}
% \end{macro}
%
%    \begin{macrocode}
}}{%
  \providecommand*\float@ifcaption[1]{\@secondoftwo}%
% \clearcaptionsetup{boxed}% 是否被 floatrow 宏包使用？
}
%    \end{macrocode}
%
% `boxed' 浮动体与其题注之间的间距默认为 |2pt|。
%
%    \begin{macrocode}
\captionsetup[boxed]{skip=2pt} % 未使用时不发出警告
%    \end{macrocode}
%
% \changes{v3.0f}{2005/05/26}{向题注样式 \opt{ruled} 加入 \opt{strut=0}}
% \changes{v3.0f}{2005/05/27}{修复选项 \opt{ruled}}
% 为模仿 |\@fs@capt| 的 `ruled' 定义，我们提供带有相应选项的题注样式
%  `ruled'。但若指定了宏包选项 |ruled|，则设置若干题注参数以模仿
%  \thispackage\ \version{1.x} 选项 |ruled| 的行为：沿用当前题注设置，
%  但不使用边距，也不做“单行检测”。
%    \begin{macrocode}
\caption@ifbool{ruled}{%
  \captionsetup[ruled]{margin=0pt,minmargin=0,slc=0}%
}{%
  \DeclareCaptionStyle{ruled}{labelfont=bf,labelsep=space,strut=0}%
}
\caption@undefbool{ruled}
%    \end{macrocode}
%
% \subsubsection{floatflt 宏包}
% \changes{v3.1}{2007/06/10}{加入对 \pkg{floatflt} 宏包的支持}
%
%    \begin{macrocode}
\caption@IfPackageLoaded{floatflt}[1996/02/27 v1.3]{%
%    \end{macrocode}
%
% \begin{macro}{\floatingfigure}
% 修补 |\floatingfigure|，使其使用 |\caption@floatflt|。
%    \begin{macrocode}
  \let\caption@ORI@floatingfigure\floatingfigure
  \def\floatingfigure{%
    \caption@floatflt{figure}%
    \caption@ORI@floatingfigure}%
%    \end{macrocode}
% \end{macro}
%
% \begin{macro}{\floatingtable}
% 对 |\floatingtable| 同样处理……
%    \begin{macrocode}
  \let\caption@ORI@floatingtable\floatingtable
  \def\floatingtable{%
    \caption@floatflt{table}%
%   \caption@setautoposition b%
    \caption@ORI@floatingtable}%
%    \end{macrocode}
% \end{macro}
%
% \begin{macro}{\caption@floatflt}
% \changes{v3.1c}{2007/10/06}{缺陷修复：加入 \cs{caption@clearmargin}}
% 此处做两件事：
%  \begin{enumerate}
%    \item 使用 |\caption@set|\-|options{|\-|floating|\meta{类型}|}|，以支持
%          |\caption|\-|setup[|\-|floating|\meta{类型}|]{|\ldots|}|。
%    \item 必须正确设置 |\line|\-|width|。通常这由 |\@par|\-|box|\-|restore|
%          在 |\@caption| 内部完成；但此处使用 |\@caption@parbox|\-|restore|，
%          因而需将其映射到 |\@par|\-|box|\-|restore|。
%  \end{enumerate}
%    \begin{macrocode}
  \newcommand*\caption@floatflt[1]{%
    \caption@settype{#1}%
    \caption@clearmargin
    \caption@setoptions{floating#1}%
    \caption@setanchor}%
%    \end{macrocode}
% \end{macro}
%
%    \begin{macrocode}
}{}
%    \end{macrocode}
%
% \subsubsection{floatrow 宏包}
%
%    \begin{macrocode}
\caption@IfPackageLoaded{floatrow}[2008/08/02 v0.3b]{%
%    \end{macrocode}
%
% \begin{macro}{\captionlabel}
% \changes{v3.6e}{2022/04/17}{加入本宏的重定义}
% |\caption|\-|label| 的原始定义使用内部宏
%  |\caption@set|\-|sub|\-|type|，该宏已不存在，需替换为
%  |\set|\-|caption|\-|sub|\-|type|。
%    \begin{macrocode}
  \renewcommand\captionlabel[1]{{\def\FR@tmp{\@captype}\ifnum\floatbox@depth>\@ne
       \def\FR@tmp{sub\@captype}\setcaptionsubtype*\stepcounter{\FR@tmp}\fi
   \caption@@@make{\caption@fnum\FR@tmp}{#1}}}
%    \end{macrocode}
% \end{macro}
%
% \begin{environment}{subfloatrow}
% \changes{v3.6d}{2022/04/17}{加入本环境的重定义}
% 环境 |sub|\-|float|\-|row| 的原始定义使用
%  |\caption|\-|setup{sub|\-|type}|，将 |\caption| 等题注命令切换为
%  \pkg{subcaption} 宏包提供的对应命令。自 \pkg{caption} \version{3.6}
%  起，这已无法与 \pkg{subfig} 宏包协同工作；因此此处用一个钩子替代，
%  可由 \pkg{subcaption} 宏包填充。
%    \begin{macrocode}
  \renewenvironment{subfloatrow}{\capsubrowsettings
    \caption@subfloatrow@hook % 作为 \captionsetup{subtype} 的替代
    \@nameuse{subfloatrow*}}{\@nameuse{endsubfloatrow*}}
%    \end{macrocode}
%    \begin{macrocode}
  \providecommand*\caption@subfloatrow@hook{}
%    \end{macrocode}
% \end{environment}
%
% \begin{environment}{subfloatrow*}
% \changes{v3.6d}{2022/04/17}{加入本环境的重定义}
% 某些宏包（如 \pkg{float} 与 \pkg{floatrow}）会为重定义
%  |\@make|\-|caption|，对主题注做特殊处理。由于这通常无法与子题注
%  良好配合，命令 |\set|\-|caption|\-|sub|\-|type| 会恢复
%  |\@make|\-|caption| 的定义。但此处不希望如此，因此需确保子题注也
%  使用 |\flrow@make|\-|caption|。
%    \begin{macrocode}
  \renewenvironment{subfloatrow*}[1][2]{\let\flrow@to\empty
    \let\Xhsize\sXhsize%\let\Zhsize\sZhsize
    \FR@ifunloc{\let\FR@ifcountH\@firstoftwo}\let\c@FRobj\c@FRsobj
    \def\FRleftmargin{}\def\FRrightmargin{}\let\floatrowsep\subfloatrowsep
    \floatrow[#1]\killfloatstyle
    \let\caption@makecaption\flrow@makecaption}% <- 加入本行
    {\ifdim\lastskip>\z@\unskip\fi
     \@ifundefined{adjustsubfloats}\relax\adjustsubfloats\endfloatrow}
%    \end{macrocode}
% \end{environment}
%
%    \begin{macrocode}
}{}
%    \end{macrocode}
%
% \subsubsection{fltpage 宏包}
% \changes{v3.1}{2007/07/22}{加入对 \pkg{fltpage} 宏包的支持}
%
%    \begin{macrocode}
\caption@IfPackageLoaded{fltpage}[1998/10/29 v.0.3]{%
  \caption@setbool{needfreeze}{1}%
%    \end{macrocode}
%
% \begin{macro}{\FP@positionLabel}
% \changes{v3.2c}{2011/08/28}{加入本重定义}
% 原始代码：
% \begin{verbatim}
% \newcommand{\FP@positionLabel}{%
%   FP\@captype-\number\value{FP@\@captype C}-pos}
% \end{verbatim}%^^A
%    \begin{macrocode}
  \renewcommand\FP@positionLabel{%
    FP\FP@captype-\number\value{FP@\FP@captype C}-pos}%
%    \end{macrocode}
% \end{macro}
%
% \begin{macro}{\FP@helpNote}
% 原始代码：
% \begin{verbatim}
% \newcommand{\FP@helpNote}[2]{%
%   \typeout{FP#1 is inserted on page \pageref{#2}!}}%
% \end{verbatim}%^^A
%    \begin{macrocode}
  \renewcommand\FP@helpNote[2]{%
    \begingroup % save \caption@thepage
      \caption@pageref{#2}%
      \typeout{FP#1 is inserted on page \caption@thepage!}%
    \endgroup}%
%    \end{macrocode}
% \end{macro}
%
% \begin{macro}{\FP@floatBegin}
% \changes{v3.1f}{2007/11/16}{现改用“空格技巧”代替 \cs{ignorespaces}}
% \changes{v3.1k}{2009/10/09}{缺陷修复 09-09-11：补上遗漏的 \cs{ignorespaces}}
% \changes{v3.2c}{2011/08/28}{适配 \cs{caption@freeze} 的当前版本}
% \changes{v3.2f}{2012/02/19}{缺陷修复 12-02-16：加入 \texttt{minipage} 环境}
% \changes{v3.3}{2013/05/01}{以 \cs{caption@freezetype} 取代 \cs{caption@freeze} 的使用}
% 原始代码：
% \begin{verbatim}
% \newcommand{\FP@floatBegin}[1]{%
%   \gdef\@captype{#1}%
%   \global\let\FP@savedCaptionCommand\caption%
%   \global\let\FP@savedLabelCommand\label%
%   \ifthenelse{\equal{\@captype}{figure}}
%      {\global\let\old@Fnum\fnum@figure}%
%      {\global\let\old@Fnum\fnum@table}%
%   \let\FP@LabelText\@empty%
%   \let\FP@CaptionText\@empty%
%   \let\FP@optionalCaptionText\@empty%
%   \renewcommand\label[1]{\gdef\FP@LabelText{##1}}%
%   \renewcommand\caption[2][]{%
%     \gdef\FP@optionalCaptionText{##1}\gdef\FP@CaptionText{##2}}%
%   \begin{lrbox}{\FP@floatCorpusBOX}%
% }%
% \end{verbatim}%^^A
%    \begin{macrocode}
  \renewcommand*\FP@floatBegin[1]{%
    \def\FP@captype{#1}%
    \begin{lrbox}{\FP@floatCorpusBOX}%
    \minipage\hsize % changes from LR mode to vertical mode
    \caption@freezetype{#1}%
    \ignorespaces}%
%    \end{macrocode}
% \end{macro}
%
% \begin{macro}{\FP@floatEnd}
% \changes{v3.1g}{2008/01/30}{若干小缺陷修复}
% \changes{v3.2c}{2011/08/28}{适配 \cs{caption@freeze} 的当前版本}
% \changes{v3.2f}{2012/02/19}{缺陷修复 12-02-16：加入 \texttt{minipage} 环境}
% 原始代码：
% \begin{verbatim}
% \newcommand{\FP@floatEnd}{%
%   \end{lrbox}%
%   \global\setbox\FP@floatCorpusBOX=\box\FP@floatCorpusBOX
%   \stepcounter{FP@\@captype C}%
%   \FP@savedLabelCommand{\FP@positionLabel}%
%   \FP@helpNote{\@captype}{\FP@positionLabel}%
%   \FP@float
%     {\FP@positionLabel}% location label test
%     {\begin{\@captype}[p!]
%        \usebox{\FP@floatCorpusBOX}%
%        \refstepcounter{\@captype}%
%        \ifthenelse{\equal{\FP@LabelText}{\@empty}}
%          {}{\FP@savedLabelCommand{\expandafter\protect\FP@LabelText}}%
%      \end{\@captype}}
%     {\addtocounter{\@captype}{-1}}
%     {\begin{\@captype}[b!]%
%        \ifthenelse{\equal{\FP@guide}{\@empty}}%
%          {}{\ifthenelse{\equal{\@captype}{figure}}%
%              {\renewcommand{\fnum@figure}{\old@Fnum\ {\FP@guide}}}%
%              {\renewcommand{\fnum@table}{\old@Fnum\ {\FP@guide}}}}%
%        \setlength{\abovecaptionskip}{2pt plus2pt minus 1pt} % length above caption
%        \setlength{\belowcaptionskip}{2pt plus2pt minus 1pt} % length above caption
%        \FP@separatorCaption%
%        \ifthenelse{\equal{\FP@optionalCaptionText}{\@empty}}%
%          {\FP@savedCaptionCommand{\expandafter\protect\FP@CaptionText}}%
%          {\FP@savedCaptionCommand[\expandafter\protect\FP@optionalCaptionText]%
%                                  {\expandafter\protect\FP@CaptionText}}%
%      \end{\@captype}}%
% }%
% \end{verbatim}%^^A
%    \begin{macrocode}
  \renewcommand*\FP@floatEnd{%
    \endminipage
    \end{lrbox}%
%    \end{macrocode}
%    \begin{macrocode}
    \stepcounter{FP@\FP@captype C}%
    \caption@label\FP@positionLabel
    \FP@helpNote\FP@captype\FP@positionLabel
%    \end{macrocode}
%    \begin{macrocode}
    \FP@float
      {\FP@positionLabel}% 位置标签测试
      {\caption@prepare@defrost
       \begin\FP@captype[p!]%
         \usebox\FP@floatCorpusBOX
       \end\FP@captype}%
      {\@ifundefined{theH\FP@captype}{}{%
         \expandafter\l@addto@macro\csname theH\FP@captype\endcsname{.FP}}}%
      {\caption@prepare@defrost
       \begin\FP@captype[b!]%
         \let\FP@savedSetfnumCommand\caption@setfnum
         \def\caption@setfnum##1{%
           \FP@savedSetfnumCommand{##1}%
           \ifx\FP@guide\@empty \else
             \expandafter\l@addto@macro\csname fnum@##1\endcsname{\ {\FP@guide}}%
           \fi}%
         \setlength\abovecaptionskip{2pt plus 2pt minus 1pt}% 题注上方间距
         \setlength\belowcaptionskip{2pt plus 2pt minus 1pt}% 题注下方间距
         \caption@setoptions{FP\@captype}%
         \FP@separatorCaption
         \caption@defrost
       \end\FP@captype}%
  }%
%    \end{macrocode}
% \end{macro}
%
%    \begin{macrocode}
}{%
  \let\caption@ifFPlistcap\@undefined
  \let\caption@ifFPrefcap\@undefined
}
%    \end{macrocode}
%
% \subsubsection{hyperref 宏包}
% \changes{v3.0f}{2005/06/11}{加入对 \pkg{hyperref} 宏包的支持}
% \changes{v3.1}{2007/03/09}{以 \cs{hyper@makecurrent} 代替 \cs{theHfigure} 作为 \pkg{hyperref} 宏包的判据}
% \changes{v3.2}{2011/07/30}{修订对 \pkg{hyperref} 是否提前停止的检测}
% \changes{v3.2c}{2011/08/24}{把对 \pkg{hyperref} 是否提前停止的检测适配到其当前版本：加入 \cs{IfHyperBoolean} 的使用}
%
%    \begin{macrocode}
\caption@IfPackageLoaded{hyperref}[2003/11/30 v6.74m]{%
  % 检测 hyperref 是否提前停止
  \caption@ifundefined\IfHyperBoolean{%
    \caption@set@bool\caption@ifhyp@stoppedearly0%
    \caption@ifundefined\H@refstepcounter
      {\caption@set@bool\caption@ifhyp@stoppedearly1}{%
    \caption@ifundefined\hyper@makecurrent
      {\caption@set@bool\caption@ifhyp@stoppedearly1}{%
    \caption@ifundefined\measuring@true
      {\caption@set@bool\caption@ifhyp@stoppedearly1}{}}}%
  }{%
    \def\caption@ifhyp@stoppedearly{\IfHyperBoolean{stoppedearly}}%
  }%
  \caption@ifhyp@stoppedearly{% hyperref 已提前停止
    \caption@InfoNoLine{%
      Hyperref support is turned off\MessageBreak
      because hyperref has stopped early}%
  }{%
    \AtCaptionSingleLineCheck{\measuring@true}%
%    \end{macrocode}
%
% \begin{macro}{\caption@refstepcounter@}
% 重定义 |\caption@refstepcounter@|，使 |\caption| 与 |\captionlistentry|
% 内部使用 |\H@refstepcounter|，而非 |\refstepcounter|。
%    \begin{macrocode}
    \renewcommand*\caption@refstepcounter@{\H@refstepcounter}%
%    \end{macrocode}
% \end{macro}
%
% \begin{macro}{\caption@makecurrent}
% \changes{v3.1l}{2010/01/09}{适配 nameref 宏包的当前版本}
% 重定义 |\caption@makecurrent|，使 \pkg{hyperref} 标签在 |\@caption|
% 内部定义。
% \Note{将由 \cs{caption@start} 重定义。}
%    \begin{macrocode}
    \renewcommand*\caption@makecurrent[2]{%
      \caption@makecurrentHref{#1}%
      \caption@Debug{hyperref current=\@currentHref}%
      \caption@gettitle{#2}}%
%    \end{macrocode}
%    \begin{macrocode}
    \newcommand*\caption@makecurrentHref{\hyper@makecurrent}%
%    \end{macrocode}
% \end{macro}
%
% \begin{macro}{\caption@makeanchor}
% \changes{v3.1l}{2010/01/09}{适配 hyperref 宏包当前版本}
% 重定义 |\caption@makeanchor|，使 \pkg{hyperref} 锚点在 |\@caption|
% 内部设置。
% \Note{将由 \cs{caption@start} 重定义。}
%    \begin{macrocode}
    \renewcommand\caption@makeanchor[1]{%
      \caption@Debug{hyperref anchor: \@currentHref}%
      % 若不能嵌套，锚点为空。
      \ifHy@nesting
        \expandafter\hyper@@anchor\expandafter{\@currentHref}{#1}%
      \else
        \Hy@raisedlink{%
          \expandafter\hyper@@anchor\expandafter{\@currentHref}{\relax}%
        }#1%
      \fi}%
%    \end{macrocode}
%    \begin{macrocode}
    \AtCaptionSingleLineCheck{\let\caption@makeanchor\@firstofone}%
%    \end{macrocode}
% \end{macro}
%
% \paragraph{hypcap 选项}
%
% \begin{macro}{\if@capstart}
% 与 \pkg{hypcap} 宏包一样，我们也定义开关 |\if@capstart|。
%    \begin{macrocode}
    \newif\if@capstart
%    \end{macrocode}
% \end{macro}
%
% \begin{macro}{\caption@start}
% \changes{v3.0l}{2007/02/20}{%
%         \cs{caption@currentHref} 改名为 \cs{hc@currentHref}}
% \changes{v3.1}{2007/06/21}{%
%         本宏由 \cs{capstart} 改名为 \cs{caption@start} 并作修订}
% \pkg{hypcap} 宏包定义的宏名为 |\capstart|；本变体称为
%  |\caption@start|，由选项 |hypcap=||false|/|true| 控制。
%    \begin{macrocode}
    \def\caption@start{\caption@ifhypcap\caption@start@\relax}%
    \def\caption@start@{%
%    \end{macrocode}
% 生成 \pkg{hyperref} 标签并设置 \pkg{hyperref} 锚点；通常（若
%  |hypcap=false|）两者都在 |\@caption| 内部完成。
%    \begin{macrocode}
      \caption@makestart\@captype
      \caption@startanchor\@currentHref
%    \end{macrocode}
% 阻止 |\@caption| 生成新的 \pkg{hyperref} 标签，改用保存在
%  |\hc@currentHref| 中的标签。（同时支持 \pkg{hypcap} 宏包的
%  |@capstart| 标志。）
%    \begin{macrocode}
      \global\@capstarttrue
      \let\hc@currentHref\@currentHref
      \def\caption@makecurrentHref##1{%
        \global\@capstartfalse
        \global\let\@currentHref\hc@currentHref}%
%    \end{macrocode}
% 由于锚点已经设置，阻止 |\@caption| 再生成 \pkg{hyperref} 锚点。
%    \begin{macrocode}
      \let\caption@makeanchor\@firstofone
%    \end{macrocode}
%    \begin{macrocode}
    }%
%    \end{macrocode}
% \end{macro}
%
% \begin{macro}{\caption@makestart}
% \changes{v3.1c}{2007/09/30}{现以 \texttt{caption} 代替 \texttt{xxx} 扩展 \cs{@currentHref}}
%  |\caption@makestart|\marg{类型}
%  在 |\caption@start| 内部定义一个 \pkg{hyperref} 锚点。\par
%  由于提供了 |\continuedfloat|，浮动体计数器可能在“此刻”与
%  |\caption| 之间发生变化；也就是说，我们尚不知道图或表的计数器值，
%  因而无法生成“正确”的 \pkg{hyperref} 标签。
%  想到两种解决办法：
%  \begin{enumerate}
%  \item 可借助 aux 文件实现。\par
%  ——或者——
%  \item 在局部设置 |hypertexnames=false|。
%  此外，为避免与用 |hypertexnames=true| 生成的其他超链接冲突，
%  |\@currentHref| 的命名方案采用 |#1.caption.|\meta{计数器}
%  （而非 |#1.|\meta{计数器}）。
%  \end{enumerate}
%  第一种方案的优点是能生成“正确”的锚点名，但在插入或删除图、表时
%  需要额外再跑一遍 \LaTeX。\par
%  第二种方案实现非常简单，但会带来一些副作用，例如锚点名不再跟随图或
%  表的题注编号。\par
%  由于图省事，这里采用第二种方案；日后或许会再作修订。
%    \begin{macrocode}
    \newcommand*\caption@makestart[1]{%
      \begingroup
        \Hy@hypertexnamesfalse
%       \gdef\@currentHlabel{}%
        \hyper@makecurrent{#1.caption}%
      \endgroup
      \caption@Debug{hypcap start=\@currentHref}}%
%    \end{macrocode}
% \end{macro}
%
% \begin{macro}{\caption@startanchor}
% \changes{v3.1b}{2007/09/17}{缺陷修复 06-09-17：改用 \cs{caption@anchor} 代替 \cs{hyper@@anchor}}
% \changes{v3.1g}{2008/03/16}{把关于 TeX 模式的警告改为调试信息}
% |\caption@startanchor|\marg{超链接锚点}
% 在 |\caption@start| 内部设置一个 \pkg{hyperref} 锚点。\par
% 本代码取自 \pkg{hypcap} 宏包\cite{hypcap}并作了适配。
% \Note{由于 \cs{hyper@@anchor}\marg{超链接锚点}\csmarg{\string\relax} 可能导致
% 由竖直模式进入水平模式（疑为 \pkg{hyperref} 宏包的设计缺陷！），而
% \pkg{hypcap} 宏包中的变通写法 \cs{let}\cs{leavevmode}\cs{relax}
% 并非总是有效（例如在“直接 pdfmark 支持”与
% \texttt{breaklinks=true} 情形下），因此此处改用 \cs{caption@anchor}
% 代替 \cs{hyper@@anchor}。}
%    \begin{macrocode}
    \newcommand*\caption@startanchor[1]{%
      \ifvmode\begingroup
        \caption@Debug{hypcap anchor: #1 (vertical mode)}%
        \@tempdima\prevdepth
        \nointerlineskip
        \vspace*{-\caption@hypcapspace}%
        \caption@anchor{#1}%
        \vspace*{\caption@hypcapspace}%
        \prevdepth\@tempdima
      \endgroup\else
        \caption@Debug{hypcap anchor: #1 (horizontal mode)}%
        \caption@anchor{#1}%
      \fi}%
%    \end{macrocode}
% \end{macro}
%
% \begin{macro}{\caption@anchor}
% \changes{v3.1c}{2007/10/17}{变通写法 07-10-17：现改用 \cs{caption@raisedlink} 代替 \cs{Hy@raisedlink}}
% |\caption@anchor|\marg{超链接锚点} 设置一个 \pkg{hyperref} 锚点。
%    \begin{macrocode}
    \newcommand*\caption@anchor[1]{%
      \ifmeasuring@ \else
        \caption@raisedlink{\hyper@anchorstart{#1}\hyper@anchorend}%
      \fi}%
%    \end{macrocode}
% \Note{由于 \cs{Hy@raisedlink} 会改动 \cs{@tempdima}，我们用 \cs{ifvmode}
% 将其包裹起来，以消除 \texttt{sideways} 浮动体中
% “\texttt{LaTeX Warning: Float too large for page by 1.0pt}” 的警告。
% （自 \pkg{hyperref} \version{6.77} 起已无此必要。）}
%    \begin{macrocode}
    \ifx\HyperRaiseLinkLength\@tempdima
      \def\caption@raisedlink#1{\ifvmode#1\else\Hy@raisedlink{#1}\fi}%
    \else
      \let\caption@raisedlink\Hy@raisedlink
    \fi
%    \end{macrocode}
% \end{macro}
%
% \begin{macro}{\caption@@start}
% 由 |\caption@freezeHref| 使用。此外，若预期来自 |\caption@start| 的已保存
%  \pkg{hyperref} 标签并不存在，则发出警告。
%    \begin{macrocode}
    \def\caption@@start{%
      \caption@ifundefined\hc@currentHref{%
        \caption@Warning{%
          The option `hypcap=true' will be ignored for this\MessageBreak
          particular \string\caption}}{}}%
%    \end{macrocode}
% \end{macro}
%
% \begin{macro}{\caption@freezeHref}
%  阻止 |\caption@start| 生成 \pkg{hyperref} 标签并设置 \pkg{hyperref}
%  锚点；改为：若 |\@caption| 生成 \pkg{hyperref} 标签，则将其存入
%  |\caption@currentHref|。此外还需重定义 |\caption@setfloatcapt|，使
%  |\@caption| 内不放置 \pkg{hyperref} 锚点。
% \iffalse
% （由于 |\caption@setfloatcapt| 在 |\subcaption| 内部会被设为
%  |\@firstofone|，无需把它的恢复加入 |\caption@subtypehook|。）
% \fi
%    \begin{macrocode}
    \def\caption@freezeHref{%
      \let\caption@ORI@start\caption@start
      \def\caption@start{\let\caption@start\caption@ORI@start}%
%    \end{macrocode}
%    \begin{macrocode}
%     \let\caption@ORI@@start\caption@@start
%     \l@addto@macro\caption@subtypehook{%
%       \let\caption@@start\caption@ORI@@start}%
%    \end{macrocode}
%    \begin{macrocode}
      \global\let\caption@currentHref\@undefined
      \def\caption@@start{\global\let\caption@currentHref\@currentHref}%
%    \end{macrocode}
%    \begin{macrocode}
      \let\caption@ORI@setfloatcapt\caption@setfloatcapt
      \renewcommand*\caption@setfloatcapt{%
        \ifx\caption@currentHref\@undefined \else
          \let\caption@makeanchor\@firstofone
        \fi
        \caption@ORI@setfloatcapt}}%
%    \end{macrocode}
% \end{macro}
%
% \begin{macro}{\caption@defrostHref}
%  若存在已冻结的 |\@currentHref|，在此设置 \pkg{hyperref} 锚点。
%    \begin{macrocode}
    \def\caption@defrostHref{%
      \ifx\caption@currentHref\@undefined \else
        \caption@startanchor\caption@currentHref
        \global\let\caption@currentHref\@undefined
      \fi}%
%    \end{macrocode}
% \end{macro}
%
%    \begin{macrocode}
  }}{}
%    \end{macrocode}
%
% \subsubsection{hypcap 宏包}
% \changes{v3.0f}{2005/06/22}{加入对 \pkg{hypcap} 宏包的支持}
% \changes{v3.0l}{2007/02/20}{适配 \pkg{hypcap} \version{1.6}}
% \changes{v3.1e}{2007/11/04}{把 \pkg{hypcap} 支持适配到 \pkg{hyperref} \version{6.77}}
%
%    \begin{macrocode}
\caption@IfPackageLoaded{hypcap}{% v1.0
  \ifx\caption@start\relax \else % hyperref 未提前停止
%    \end{macrocode}
%
% 若已加载 \pkg{hypcap} 宏包，则放弃本宏包自身的超链接定位算法，改由
%  \pkg{hypcap} 宏包控制定位。
%
% \begin{macro}{\capstart}
% \changes{v3.1k}{2009/10/09}{加入对 \cs{ifcapstart}（\pkg{hypcap} 宏包 \version{1.10}）的支持}
% 做法很简单：把 |\capstart| 映射到 |\caption@start@|。不过本代码的行为
% 与原始版本并不完全一致：原始 |\capstart| 只作用于下一个 |\caption|，
% 而本版本会影响同一环境内\emph{所有} |\caption|，直到放置新的
%  |\capstart| 为止。
%    \begin{macrocode}
    \caption@ifundefined\capstarttrue % check for v1.10 of hypcap package
      {\def\capstart{\caption@start@}}%
      {\def\capstart{\ifcapstart\caption@start@\fi}}%
    \let\caption@start\relax
    \let\caption@@start\relax
%    \end{macrocode}
% \end{macro}
%
% \begin{macro}{\caption@hypcapspace}
% 此外，把本宏包的 |\caption@hypcapspace| 映射到 \pkg{hypcap} 宏包提供的
%  |\hypcapspace|。
%    \begin{macrocode}
    \caption@set@bool\caption@ifhypcap 1%
    \renewcommand*\caption@hypcapspace{\hypcapspace}%
%    \end{macrocode}
% \end{macro}
%
%    \begin{macrocode}
  \fi}{}
%    \end{macrocode}
%
% \subsubsection{listings 宏包}
% \changes{v3.0b}{2004/05/16}{加入对 \pkg{listings} 宏包的支持}
%
%    \begin{macrocode}
\caption@IfPackageLoaded{listings}[2004/02/13 v1.2]{%
%    \end{macrocode}
%
% \begin{macro}{\lst@MakeCaption}
% \changes{v3.1g}{2008/01/20}{缺陷修复：现正确处理间距}
% \changes{v3.1k}{2008/03/29}{缺陷修复 08-10-16-dctt：加入 \texttt{rule=0}}
% \changes{v3.2}{2010/10/26}{缺陷修复 10-07-20a：必要时把 \cs{hsize} 设为 \cs{linewidth}}
% \changes{v3.4c}{2020/01/01}{加入用于列表环境中的粗糙变通写法}
% \changes{v3.4d}{2020/01/01}{修订用于列表环境中的粗糙变通写法}
% \changes{v3.6}{2021/01/04}{加入对 \cs{addcontentsline} 的重定义}
% 为支持 \pkg{listings} 宏包，需重定义 |\lst@Make|\-|Caption|，使原有
% 内容与 |\caption@begin|、|\caption@end| 等嵌套起来。\par
% 在 \pkg{listings} 宏包的“Init”钩子内部设置一个标志（|\line|\-|width|
% 等的适配就在此处进行）。
%    \begin{macrocode}
  \caption@setbool{lst@Init}{0}%
  \lst@AddToHook{Init}{\caption@setbool{lst@Init}{1}}%
%    \end{macrocode}
%    \begin{macrocode}
  \let\caption@ORI@lst@MakeCaption\lst@MakeCaption
  \renewcommand*\lst@MakeCaption[1]{% #1 is `t' or `b'
%    \end{macrocode}
% \pkg{listings} 宏包总会两次调用本宏（参数为 `t' 或 `b'），
% 因此此处需要一个额外的分组。
%    \begin{macrocode}
    \begingroup
%    \end{macrocode}
% 在列表环境中，\pkg{listings} 宏包会忽略题注的左边距
% （|\@total|\-|left|\-|margin|）。我们通过在 |\caption@box| 之前加入一个
% 负间距来模仿该行为。\par
% 注意：仅当 \pkg{listings} 宏包的“Init”钩子尚未调用时才需要这样做。
%    \begin{macrocode}
      \caption@iflst@Init{}{%
        \edef\caption@lst@leftskip{\noexpand\hskip-\the\@totalleftmargin\noexpand\relax}%
        \l@addto@macro\caption@beginhook\caption@lst@leftskip}%
%    \end{macrocode}
% 设置 |position=#1|；若为 `top'，则交换上下间距，以不改变 \pkg{listings}
% 宏包的默认行为。（注意：在调用 \cs{lst@MakeCaption} 之前，\pkg{listings}
% 宏包已设置了自身的 |\abovecaptionskip| 与 |\belowcaptionskip| 取值。）
%    \begin{macrocode}
      \caption@setposition{#1}%
      \caption@iftop{%
        \@tempdima\belowcaptionskip
        \belowcaptionskip\abovecaptionskip
        \abovecaptionskip\@tempdima}{}%
%    \end{macrocode}
% 针对间距错误问题的变通写法（待定：尚需进一步研究）
%    \begin{macrocode}
      \caption@setup{rule=0}%
%    \end{macrocode}
% 随后设置局部 `lstlisting' 选项。
%    \begin{macrocode}
      \caption@setoptions{lstlisting}%
%    \end{macrocode}
% 若此时 |position=| 设为 |auto|，则沿用 \pkg{listings} 宏包的
%  |captionpos=| 设置。
%    \begin{macrocode}
      \caption@setautoposition{#1}%
%    \end{macrocode}
%    \begin{macrocode}
% 重定义 |\addcontentsline|，使 lol 条目等使用本宏包的目录格式。
      \ifx\lst@@caption\@empty \else
        \let\caption@ORI@addcontentsline\addcontentsline
        \renewcommand*\addcontentsline[3]{%
          \let\addcontentsline\caption@ORI@addcontentsline
          \caption@addcontentsline{lstlisting}{\lst@@caption}}%
      \fi
%    \end{macrocode}
% 最后，与本宏包 \cs{@caption} 代码中做类似处理。
%    \begin{macrocode}
      \caption@begin{lstlisting}%
        \caption@ORI@lst@MakeCaption{#1}%
      \caption@end
%    \end{macrocode}
%    \begin{macrocode}
    \endgroup}%
%    \end{macrocode}
% \end{macro}
%
% \begin{macro}{\lst@makecaption}
% \begin{macro}{\lst@maketitle}
% 用于排版 |caption=| 或 |title=| 取值的包装宏。
%    \begin{macrocode}
  \def\lst@makecaption{\caption@starfalse\@makecaption}%
  \def\lst@maketitle{\caption@startrue\@makecaption\@empty}%
%    \end{macrocode}
% \end{macro}
% \end{macro}
%
% \begin{macro}{\ext@lstlisting}
% \changes{v3.1}{2007/03/03}{新增本宏}
% 由于 \pkg{listings} 宏包未定义 |\ext@lstlisting|，而终端用户使用
%  |\captionof{lstlisting}| 时需要它，故在此定义。
%    \begin{macrocode}
  \providecommand*\ext@lstlisting{lol}%
%    \end{macrocode}
% \end{macro}
%
% \begin{macro}{\lst@@caption}
% \changes{v3.1}{2007/03/03}{新增本宏}
% \changes{v3.1a}{2007/07/13}{缺陷修复 07-09-13：不再重定义 \cs{lst@@caption}}
% \changes{v3.2}{2010/10/26}{缺陷修复 09-05-15：把 \cs{lst@@caption} 的设置加入 \cs{caption@beginex}}
% \changes{v3.5h}{2022/03/01}{把 \cs{lst@@caption} 的重定义由 \cs{caption@beginex} 移至 \cs{caption@beginex@hook}}
% 将 |\lst@@caption| 定义为非 |\@empty|，使 |\fnum@lst|\-|listing|
% 包含编号。
%    \begin{macrocode}
  \l@addto@macro\caption@beginex@hook{\let\lst@@caption\relax}%
%    \end{macrocode}
% \end{macro}
%
%    \begin{macrocode}
}{}
%    \end{macrocode}
%
% \subsubsection{longtable 宏包}
%
% \begin{macro}{\LTcaptype}
% \changes{v3.1}{2007/07/06}{新增本宏及其支持}
%  |\LTcaptype| 预置为 |table|。
%    \begin{macrocode}
\providecommand*\LTcaptype{table}
%    \end{macrocode}
% \end{macro}
%
% \begin{macro}{\caption@AtBeginLongtable}
% \changes{v3.5}{2018/04/15}{新增本宏}
%    \begin{macrocode}
\providecommand*\caption@LT@array{}
\newcommand\caption@AtBeginLongtable{%
  \g@addto@macro\caption@LT@array}
%    \end{macrocode}
% \end{macro}
%
%    \begin{macrocode}
\caption@IfPackageLoaded{longtable}[1995/05/24 v3.14]{%
%    \end{macrocode}
%    \begin{macrocode}
  \RequirePackage{ltcaption}[2018/08/26]%
  \let\LT@@makecaption\@undefined
%    \end{macrocode}
%
% \begin{macro}{\LT@array}
% \changes{v3.1}{2007/07/22}{加入本重定义}
% \changes{v3.1a}{2007/09/16}{缺陷修复：本重定义总是在 \cs{AtBeginDocument} 时执行}
% \changes{v3.1g}{2008/01/20}{修复 \cs{captionlistentry} 缺陷：不再递增表格计数器，\cs{nameref} 可正常工作}
% \changes{v3.3}{2016/01/31}{加入对 \cs{bicaption} 的支持}
% \changes{v3.3}{2018/10/05}{修复 \cs{captionlistentry} 的定义}
% \changes{v3.5}{2018/04/15}{加入 \cs{caption@AtBeginLongtable} 的使用}
% \changes{v3.6}{2022/01/05}{以 \cs{caption@OutsideFloat} 取代 \cs{caption@Error} 的使用}
% \changes{v3.6g}{2022/04/18}{适配 \pkg{caption3} 宏包的最新版本}
% 在此重定义 |\LT@array|，使 |\captionsetup|\marg{选项} 能在
%  |longtable| 内生效。
% \Note{由于 \pkg{hyperref} 宏包也会修补 \cs{LT@array}，而该修补仅基于
%  \cs{LT@array} 的原始定义有效，因此本操作必须在 \pkg{hyperref} 宏包之后，
%  即 \cs{AtBeginDocument} 时进行。}
%    \begin{macrocode}
  \caption@AtBeginDocument{%
    \let\caption@ORI@LT@array\LT@array
    \renewcommand*\LT@array{%
      \caption@LT@array
      \caption@ORI@LT@array}}%
%    \end{macrocode}
%    \begin{macrocode}
  \caption@AtBeginLongtable{%
% longtable 用的 |\captionsetup|：
%    \begin{macrocode}
    \global\let\caption@opt@@longtable\@undefined
    \def\captionsetup{%
      \noalign\bgroup
      \@ifstar\@captionsetup\@captionsetup}% 去掉 *
    \def\@captionsetup#1{\LT@captionsetup{#1}\egroup}%
    \def\LT@captionsetup#1{%
      \caption@setup@options\@gobble[@longtable]{#1}%
      \global\let\caption@opt@@longtable\caption@opt@@longtable}%
%    \end{macrocode}
% longtable 用的 |\captionlistentry|：
%    \begin{macrocode}
    \def\captionlistentry{%
      \noalign\bgroup
      \@ifstar{\egroup\LT@listentry\LTcaptype}% 去掉 *
              {\egroup\LT@listentry\LTcaptype}}%
%    \end{macrocode}
% longtable 用的 |\continuedfloat|：\\
% {\small（已注释掉，因为未经过深入测试，而且本来也相当不实用）}
% \Note{\pkg{hyperref} 版本 $<$ v6.76j 会使用两次 \cs{hyper@makecurrent}}
%    \begin{macrocode}
%   \caption@ifhypcap{%
%     \let\caption@ORI@hyper@makecurrent\hyper@makecurrent
%     \def\hyper@makecurrent##1{%
%       \let\hyper@makecurrent\caption@ORI@hyper@makecurrent
%       \caption@makestart{##1}%
%%      \let\Hy@LT@currentHlabel\@currentHlabel
%       \let\Hy@LT@currentHref\@currentHref
%       \def\hyper@makecurrent####1{%
%%        \let\@currentHlabel\Hy@LT@currentHlabel
%         \let\@currentHref\Hy@LT@currentHref}}%
%     \let\caption@ORI@continuedfloat\continuedfloat
%     \def\continuedfloat{\noalign{%
%       \gdef\caption@setContinuedFloat{%
%         \let\caption@reset@continuedfloat\@gobble}%
%       \def\caption@setoptions##1{%
%         \g@addto@macro\caption@setContinuedFloat{%
%           \caption@setoptions{##1}}}%
%       \let\@captype\LTcaptype
%       \caption@ORI@ContinuedFloat}}%
%   }{%
%     \def\continuedfloat{\noalign{%
%       \caption@Error{%
%         \noexpand\continuedfloat inside longtables\MessageBreak
%         is only available with `hypcap=true'}}}%
%   }%
%   \global\let\caption@setContinuedFloat\@empty
    \def\continuedfloat{\noalign{%
      \caption@OutsideFloat\continuedfloat}}%
%    \end{macrocode}
% longtable 用的 |\bicaption|：
%    \begin{macrocode}
    \let\bicaption\LT@bicaption
%    \end{macrocode}
%    \begin{macrocode}
  }%
%    \end{macrocode}
% \end{macro}
%
% \begin{macro}{\LT@capti@n}
% \changes{v3.3}{2016/02/01}{加入本重定义}
%  原始实现：
%  \begin{verbatim}
%  \def\LT@capti@n{%
%    \@ifstar
%      {\egroup\LT@c@ption\@gobble[]}%
%      {\egroup\@xdblarg{\LT@c@ption\@firstofone}}}
%  \end{verbatim}%^^A
%  本实现以 |\caption@xdblarg| 代替 |\@xdblarg|：
%    \begin{macrocode}
  \def\LT@capti@n{%
    \@ifstar
      {\egroup\LT@c@ption\@gobble[]}%
      {\egroup\caption@xdblarg{\LT@c@ption\@firstofone}}}%
%    \end{macrocode}
% \end{macro}
%
% \begin{macro}{\LT@makecaption}
% \changes{v3.0d}{2004/08/10}{缺陷修复 04-08-04：现使用 \cs{abovecaptionskip} 与 \cs{belowcaptionskip}}
% \changes{v3.0e}{2005/05/05}{缺陷修复：\cs{captionsetup[longtable]} 现覆盖 \cs{LTcapwidth}}
% \changes{v3.0h}{2005/10/07}{引入 \cs{caption@LT@make}}
% \changes{v3.3}{2016/02/01}{经 \cs{caption@LT@setup} 加入对 \cs{bicaption} 的支持}
% \changes{v3.3a}{2019/04/02}{加入针对 \pkg{floatrow} 问题的变通写法}
% \changes{v3.6}{2021/01/09}{适配 \cs{caption@prepare@stepcounter} 的当前接口}
%  |\LT@makecaption|\marg{命令}\marg{标签}\marg{文本}\par
%  \smallskip
%  原始定义：
%  \begin{verbatim}
%  \def\LT@makecaption#1#2#3{%
%    \LT@mcol\LT@cols c{\hbox to\z@{\hss\parbox[t]\LTcapwidth{%
%      % 基于 article 文档类的 "\@makecaption"：#1 在带星形式下为 "\@gobble"，
%      % 否则为 "\@firstofone"。
%      \sbox\@tempboxa{#1{#2: }#3}%
%      \ifdim\wd\@tempboxa>\hsize
%        #1{#2: }#3%
%      \else
%        \hbox to\hsize{\hfil\box\@tempboxa\hfil}%
%      \fi
%      \endgraf\vskip\baselineskip}%
%    \hss}}}
%  \end{verbatim}%^^A
%  本宏包的定义：
%    \begin{macrocode}
  \renewcommand\LT@makecaption[3]{%
    \caption@LT@make{%
      \caption@LT@settype\LTcaptype
%    \end{macrocode}
% |\caption@LT@setup| 在 \pkg{bicaption} 宏包内被重定义，包含双语题注
% 排版前的准备工作。
%    \begin{macrocode}
      \caption@LT@setup
%    \end{macrocode}
% \iffalse
%  longtable 的默认 |position=| 设置为 |top|。
% （这模仿了 \pkg{longtable} 宏包的标准行为：题注上方无间距，下方有间距。）
% \fi
%  对 longtable 而言 |position=auto| 并非良策，但仍尽力处理。该做法对
%  longtable 正文内的题注效果尚可，对 longtable（尾）注脚内的题注则
%  不尽理想。
%  \Note{若不明确，应取 `top'！}
%    \begin{macrocode}
      \caption@setautoposition{\ifcase\LT@rows t\else b\fi}%
%    \end{macrocode}
%  根据第一个参数设置 |\ifcaption@star|。
%    \begin{macrocode}
      \caption@startrue#1\caption@starfalse
%    \end{macrocode}
%    \begin{macrocode}
      \caption@prepare@stepcounter\LTcaptype
%    \end{macrocode}
% \changes{v3.1}{2007/05/05}{加入 \cs{caption@font@normal}}
%    \begin{macrocode}
      \caption@begin\LTcaptype
        \caption@normalsize
%    \end{macrocode}
%  下面这个间距用于校正 |\parbox[t]| 的高度。通常该高度是第一行的高度，
%  但由于附加了间距（|\abovecaptionskip| 与 |\belowcaptionskip|），它
%  始终为 |0pt|。\par
% （另一种思路是用 |\noalign{\vskip|\ldots|}| 在 longtable 列之外排版
%  第一个间距；但这意味着必须把 |\caption@begin| 移到别处，因为它不能在
%  表格模式下工作。目前还想不出优雅的做法……）
%    \begin{macrocode}
        \vskip-\ht\strutbox
%    \end{macrocode}
%  针对 \pkg{floatrow} 问题的变通写法：
%  |\below|\-|caption|\-|skip| 会扰乱实际排版题注的环境（产生多余的、
%  不需要的空白），因此若已加载 \pkg{floatrow} 宏包，则抑制该间距。
% （此修复对应 \issue{50}）
%    \begin{macrocode}
        \caption@ifdefined\FBifcaptop{%
          \let\caption@belowskip\@empty}{}%
%    \end{macrocode}
%  下面这段代码应当很眼熟：设置本宏包的间距，并用 |\caption@@make|
%  排版题注本身。
%    \begin{macrocode}
        \caption@make@above
        \caption@@make{#2}{#3}\endgraf
        \caption@make@below
      \caption@end}}%
%    \end{macrocode}
% \end{macro}
%
% \begin{macro}{\LT@listentry}
% \changes{v3.3}{2018/08/26}{新增本宏}
% \changes{v3.3}{2018/10/05}{缺陷修复：加入 \cs{ignorespaces}}
% \changes{v3.3}{2018/12/26}{缺陷修复：补上遗漏的花括号}
%    \begin{macrocode}
  \renewcommand*\LT@listentry[2]{%
    \begingroup
      \caption@LT@settype{#1}%
      \caption@listentry\@firstoftwo[{#1}]{#2}%
    \endgroup\ignorespaces}%
%    \end{macrocode}
% \end{macro}
%
% \begin{macro}{\caption@LT@settype}
% \changes{v3.3}{2018/08/26}{新增本宏}
% \changes{v3.3}{2018/10/05}{改用第 1 个参数，而非 \cs{LTcaptype}}
% \changes{v3.6}{2021/01/17}{适配 \cs{caption@setwidth} 的当前实现}
%    \begin{macrocode}
  \newcommand*\caption@LT@settype[1]{%
    \caption@settype{#1}%
%    \end{macrocode}
% 若 |\LTcapwidth| 未取默认值 |4in|，则认为它要覆盖本宏包的设置。
% （但 |\captionsetup[longtable]{width=|\ldots|}| 会覆盖 |\LTcapwidth|。）
%    \begin{macrocode}
    \ifdim\LTcapwidth=4in \else
      \caption@setwidth\LTcapwidth
    \fi
%    \end{macrocode}
%    \begin{macrocode}
    \caption@setoptions{longtable}%
%   \caption@setContinuedFloat
%    \end{macrocode}
% 最后设置在 longtable 内部由 |\captionsetup| 应用的选项。
%    \begin{macrocode}
    \caption@setoptions{@longtable}}%
%    \end{macrocode}
% \end{macro}
%
% \begin{macro}{\flrow@setlist}
% \changes{v3.4g}{2020/01/03}{加入本宏的重定义}
% \pkg{floatrow} 宏包在其 |\LT@array| 重定义内部使用
%  |\flrow@setlist{{table}{longtable}}|，却未先设置 |\@captype|；这导致
%  longtable 中的 |\captionsetup[table]{name=...}| 必然失败
%  （“Undefined control sequence”）。遗憾的是，我们无法再次方便地修补
%  |\LT@array| 来修复 \pkg{floatrow} 中的这个缺陷，因为它会在
%  |\AtBeginDocument| 时被修补。因此改而修补 |\flrow@setlist|，就地捕获并
%  修复该缺陷。
%    \begin{macrocode}
  \caption@AtBeginDocument{\caption@ifdefined\flrow@setlist@{%
    \caption@InfoNoLine{Fixing longtable code of `floatrow' package}%
    \let\caption@flrow@setlist@\flrow@setlist@
    \renewcommand\flrow@setlist@[1]{%
      \def\caption@tempa{{table}{longtable}}%
      \def\caption@tempb{#1}%
      \ifx\caption@tempa\caption@tempb
        \def\@captype{table}%
      \fi
      \caption@flrow@setlist@{#1}}%
  }{}}%
%    \end{macrocode}
% \end{macro}
%
%    \begin{macrocode}
}{%
  \let\caption@AtBeginLongtable\@gobble
  \let\caption@LT@array\@undefined
}
%    \end{macrocode}
%
% \begin{macro}{\caption@LT@setup}
% \changes{v3.3}{2016/02/01}{新增本宏}
% 为 \pkg{longtable} 题注排版做准备的钩子。
%    \begin{macrocode}
\providecommand*\caption@LT@setup{}
%    \end{macrocode}
% \end{macro}
%
% \subsubsection{picinpar 宏包}
% \changes{v3.1}{2007/06/13}{加入对 \pkg{picinpar} 宏包的支持}
%
%    \begin{macrocode}
\caption@IfPackageLoaded{picinpar}{%
%    \end{macrocode}
%
% \begin{macro}{\figwindow}
% \begin{macro}{\tabwindow}
% \pkg{picinpar} 宏包自带题注代码（|\wincaption|、|\@wincaption|、
%  |\@makewincaption| 等），因此此处重定义 |\figwindow| 与
%  |\tabwindow|，改用 |\caption|。
%    \begin{macrocode}
  \long\def\figwindow[#1,#2,#3,#4] {%
    \caption@window{figure}%
    \caption@setoptions{figwindow}%
    \begin{window}[#1,#2,{#3},\caption@wincaption{#4}] }%
%    \end{macrocode}
%    \begin{macrocode}
  \long\def\tabwindow[#1,#2,#3,#4] {%
    \caption@window{table}%
    \caption@setoptions{tabwindow}%
    \begin{window}[#1,#2,{#3},\caption@wincaption{#4}] }%
%    \end{macrocode}
% \end{macro}
% \end{macro}
%
% \begin{macro}{\caption@window}
% \changes{v3.1c}{2007/10/06}{缺陷修复：加入 \cs{caption@clearmargin}}
% 除调用 |\caption@settype| 外，此处还重定义 |\caption@par|\-|box|\-|restore|
% （与 \pkg{floatflt}、\pkg{picins} 宏包支持中的做法相同）以及
%  |\@makecaption|（与 \pkg{float} 宏包支持中的做法相同）。
%    \begin{macrocode}
  \newcommand*\caption@window[1]{%
    \let\@makecaption\caption@@make
    \caption@setautoposition b%
    \caption@settype{#1}%
    \caption@clearmargin
    \caption@setanchor}%
%    \end{macrocode}
% \end{macro}
%
% \begin{macro}{\caption@wincaption}
% \changes{v3.1d}{2007/10/25}{补上遗漏的 \%}
% \changes{v3.3}{2013/04/14}{加入 \cs{caption@ContinuedFloattrue}}
% \changes{v3.3}{2018/10/06}{以 \cs{caption@setflag1} 取代 \cs{caption@ContinuedFloattrue}}
% \changes{v3.6}{2020/12/28}{以 \cs{caption@clrflags}\cs{caption@setcontinued} 取代 \cs{caption@setflag1}}
% 最终借助 |\caption| 排版题注。
%    \begin{macrocode}
  \newcommand\caption@wincaption[1]{%
%    \end{macrocode}
% 每个 |figwindow| 与 |tabwindow| 题注都会执行两次：第一遍 |\picwd| 为
%  |0pt|，第二遍 |\picwd| 为 |\hsize|。
%    \begin{macrocode}
    \ifdim\picwd=\z@
      \let\caption@makecurrent\@gobbletwo
      \let\caption@@start\relax
      \caption@prepareslc
    \else
      % 不再递增图或表的计数器
      \caption@clrflags
      \caption@setcontinued
    \fi
%    \end{macrocode}
% 参数 |#1| 可以只是题注文本（如 |一则插图题注|），也可以包含可选参数
%  \meta{目录条目}（如 |[一条目录条目]||{一则插图题注}|）。因此需检测
%  |#1| 是否以~|[| 开头；此外还像 |\caption*| 一样支持带星形式，所以也要
%  检测~|*|。
%    \begin{macrocode}
    \edef\@tempa{\expandafter\noexpand\@car#1\@nil}%
    \if\@tempa*%
      \let\@tempa\@firstofone
    \else\if\@tempa[%]
      \let\@tempa\@firstofone
    \else
      \let\@tempa\@empty
    \fi\fi
    \expandafter\caption\@tempa{#1}}%
%    \end{macrocode}
% \end{macro}
%
%    \begin{macrocode}
}{}
%    \end{macrocode}
%
% \subsubsection{picins 宏包}
% \changes{v3.0j}{2006/01/26}{加入对 \pkg{picins} 宏包的支持}
%
% \begin{macro}{\piccaptiontype}
% |\piccaptiontype|\marg{类型}\par
% 提供本宏用于更改题注的 \meta{类型}，用户无需像 \pkg{picins} 文档所建议
% 那样重定义 |\@captype|。
% \Note{此处定义本宏，使其可在文档导言区使用，即便 \thispackage\ 先于
%  \pkg{picins} 宏包加载。}
%    \begin{macrocode}
\newcommand*\piccaptiontype[1]{\def\@piccaptype{#1}}
%    \end{macrocode}
% \end{macro}
%
%    \begin{macrocode}
\caption@IfPackageLoaded{picins}{%
%    \end{macrocode}
%
% 初始设置 |\@piccaptype|，并取消由 \pkg{picins} 宏包设为 |figure| 的
%  |\@captype|。
%    \begin{macrocode}
  \caption@ifundefined\@piccaptype{%
    \caption@iftype{%
      \let\@piccaptype\@captype
    }{%
      \def\@piccaptype{figure}%
    }%
  }{}%
  \let\@captype\@undefined
%    \end{macrocode}
%
% \begin{macro}{\piccaption}
%  原始代码：
%  \begin{verbatim}
%  \def\piccaption{\@ifnextchar [{\@piccaption}{\@piccaption[]}}
%  \end{verbatim}%^^A
%  本代码使用 |\caption@star|，使 |\piccaption*| 可用；并使用
%  |\caption@dblarg|，使 |\piccaption{}| 正确工作。
%    \begin{macrocode}
  \def\piccaption{\caption@star\relax{\caption@dblarg\@piccaption}}%
%    \end{macrocode}
% \end{macro}
%
% \begin{macro}{\make@piccaption}
% \changes{v3.1}{2007/12/06}{缺陷修复：现总是使用 \cs{caption@clearmargin}}
%  原始代码：
%  \begin{verbatim}
%  \def\make@piccaption{%
%   [...]
%   \setbox\@TEXT=\vbox{\hsize\hsiz@\caption[\sh@rtf@rm]{\capti@nt@xt}}%
%  }
%  \end{verbatim}%^^A
%  本代码需修正以下几点：
%  \begin{enumerate}
%    \item 必须定义 |\@captype|，因为已移除其全局定义。
%    \item 使用 |\caption@setoptions{parpic}|，以支持
%          |\captionsetup[parpic]{|\ldots|}|。
%    \item 必须正确设置 |\linewidth|。通常这由 |\@parbox|\-|restore| 在
%          |\@caption| 内部完成；但此处使用 |\@caption@parbox|\-|restore|，
%          因而需映射到 |\@parbox|\-|restore|。
%    \item |\caption| 的两个参数（|\sh@rtf@rm| 与 |\capti@nt@xt|）应在
%          第一层展开，使 |\caption[]{|\ldots|}| 与
%          |\caption[|\ldots|]{}| 正确工作。
%  \end{enumerate}
%    \begin{macrocode}
  \let\caption@ORI@make@piccaption\make@piccaption
  \def\make@piccaption{%
    \let\caption@ORI\caption
%    \end{macrocode}
%    \begin{macrocode}
    \long\def\caption[##1]##2{%
      \caption@freezeHref % 将在 \ivparpic 中解冻
      \caption@settype\@piccaptype
%     \ifnum\c@piccaptionpos>2\relax
        \caption@clearmargin
%     \else
%       \caption@width\z@ % 不使用 "width=" 设置
%     \fi
      \caption@setoptions{parpic}%
      \caption@setautoposition b%
      \caption@setanchor
%    \end{macrocode}
%    \begin{macrocode}
      \expandafter\expandafter\expandafter\caption@ORI
        \expandafter\expandafter\expandafter[%
        \expandafter\expandafter\expandafter{%
        \expandafter##1\expandafter}\expandafter]\expandafter{##2}}%
%    \end{macrocode}
% {\footnotesize\begin{quote}
%   \leavevmode\llap{-or-\quad}%^^A
%   |\begingroup|\\
%   |  \toks0\expandafter{##1}| |\toks2\expandafter{##2}|\\
%   |  \edef\x{\endgroup|\\
%   |    \noexpand\caption@ORI[{\the\toks0}]{\the\toks2}}|\\
%   |  \x|
% \end{quote}
% \begin{quote}
%   \leavevmode\llap{-or-\quad}%^^A
%   |\edef\x{%|\\
%   |  \noexpand\caption@ORI[{\unexpanded\expandafter{##1}}]%|\\
%   |                        {\unexpanded\expandafter{##2}}}|\\
%   |\-|
% \end{quote}}
%    \begin{macrocode}
    \caption@ORI@make@piccaption
    \let\caption\caption@ORI}%
%    \end{macrocode}
% \end{macro}
%

% \begin{macro}{\ivparpic}
% \changes{v3.3}{2018/12/26}{缺陷修复：补上遗漏的花括号}
% 需在此设置本宏包的 \pkg{hyperref} 锚点。由于在此重定义了 |\noindent|，
% 该做法并非万无一失！
%    \begin{macrocode}
  \let\caption@ORI@ivparpic\ivparpic
  \def\ivparpic(#1,#2)(#3,#4)[#5][#6]#7{%
    \let\caption@ORI@noindent\noindent
    \def\noindent{%
      \caption@defrostHref
      \let\noindent\caption@ORI@noindent
      \noindent}%
    \caption@ORI@ivparpic(#1,#2)(#3,#4)[{#5}][{#6}]{#7}%
    \let\noindent\caption@ORI@noindent}%
%    \end{macrocode}
% \end{macro}
%
%    \begin{macrocode}
}{%
  \let\piccaptiontype\@undefined
}
%    \end{macrocode}
%
% \subsubsection{rotating 宏包}
%
%    \begin{macrocode}
\caption@IfPackageLoaded{rotating}[1995/08/22 v2.10]{%
%    \end{macrocode}
%
% \begin{macro}{\rotcaption}
% \changes{v3.0c}{2004/07/16}{缺陷修复：移除对 \cs{caption@star} 的检测}
% \changes{v3.0i}{2005/12/07}{重写，现可与 \pkg{hyperref} 协同工作}
% 使 |\rotcaption*| 可用。
%    \begin{macrocode}
  \def\rotcaption{\let\@makecaption\@makerotcaption\caption}%
% \let\@rotcaption\@undefined
%    \end{macrocode}
% \end{macro}
%
% \begin{macro}{\rotcaptionof}
% \changes{v3.0i}{2005/12/07}{新增}
% 使 |\rotcaptionof(*)| 可用。
%    \begin{macrocode}
  \def\rotcaptionof{%
    \caption@teststar\caption@of{\rotcaption*}\rotcaption}%
%    \end{macrocode}
% \end{macro}
%
% \begin{macro}{\@makerotcaption}
% \changes{v3.0i}{2005/12/06}{重写，现在应当终于能正确工作}
% \changes{v3.0m}{2007/03/04}{使用 \cs{caption@parbox}，而非 \cs{caption@start/endbox}}
% \changes{v3.3}{2012/12/26}{（再次）重写}
% \changes{v3.6}{2021/01/17}{适配 \cs{caption@setmargin} 的当前实现}
% 原始（已修复缺陷的）代码：
%  \begin{verbatim}
%  \long\def\@makerotcaption#1#2{%
%    \setbox\@tempboxa\hbox{#1: #2}%
%    \ifdim \wd\@tempboxa > .8\vsize
%      \rotatebox{90}{%
%      \begin{minipage}{.8\textheight}#1: #2\end{minipage}%
%      }%\par   % <== 移除 \par（AR）
%    \else%
%      \rotatebox{90}{\box\@tempboxa}%
%    \fi
%    \nobreak\hspace{12pt}% <== 加入 \nobreak（AR）
%  }
%  \end{verbatim}%^^A
%  本版本模仿该行为；但若设置了 |width=|，旋转后的题注总是按
%  |\parbox| 排版。（注意：此处不支持 |margin=|。）
%    \begin{macrocode}
  \long\def\@makerotcaption#1#2{%
    \rotatebox{90}{%
      \ifdim\caption@width=\z@
        \caption@setwidth{.8\vsize}%
        \l@addto@macro\caption@singleline{%
          \caption@setup{parbox=none}}%
      \fi
      \let\caption@calcmargin\relax
      \caption@@make{#1}{#2}}%
    \nobreak\hspace{12pt}}%
%    \end{macrocode}
% \end{macro}
%
%    \begin{macrocode}
}{}
%    \end{macrocode}
%
% \subsubsection{scrextend 宏包}
% \changes{v3.6j}{2023/03/12}{加入对 \pkg{scrextend} 宏包的修补}
%
% \begin{macro}{\caption@IfKomaPackageLoaded}
% \changes{v3.6j}{2023/03/12}{加入本宏定义}
% \changes{v3.6k}{2023/03/14}{以 \cs{caption@ifdocumentclass} 代替对 \cs{scr@caption} 的检测}
%    \begin{macrocode}
\caption@ifdocumentclass{koma}%
  {\long\def\caption@IfKomaPackageLoaded#1[#2]#3#4{%
     \caption@InfoNoLine{KOMA-Script #1 package detected}%
     #3}}%
  {\let\caption@IfKomaPackageLoaded\caption@IfPackageLoaded}
%    \end{macrocode}
% \end{macro}
%
%    \begin{macrocode}
\caption@IfKomaPackageLoaded{scrextend}[2002/07/01 v3.0]{%
%    \end{macrocode}
%
% \begin{macro}{\caption@patch@addmargin}
% \changes{v3.6j}{2023/03/12}{加入本宏定义}
% \env{addmargin} 与 \env{addmargin*} 两个环境都会开启一个 \env{list}
%  \cs{item}。但在 \env{list} 项内部 \cs{hsize} 未被正确设置，因此文档类
%  提供的 \cs{@makecaption} 在 \env{addmargin} 内部使用时无法正常工作。
%  为此我们在 \cs{@makecaption} 内部尝试判断“当前是否处于
%  \env{addmargin} 环境中”，并加以修正。（唉！）
%    \begin{macrocode}
  \let\caption@addmargin@env\@addmargin
  \renewcommand*\@addmargin{%
    \@testopt\caption@@addmargin\@tempa}%
  \def\caption@@addmargin[#1]#2{%
    \caption@addmargin@env[{#1}]{#2}%
    \caption@set@addmargin
    \ignorespaces}%
%    \end{macrocode}
% \end{macro}
%
% \begin{macro}{\caption@set@addmargin}
% \changes{v3.6j}{2023/03/12}{加入本宏定义}
% 此处不能简单设置并检测一个标志，因为例如在 \env{addmargin} 环境内的
%  \cs{parbox} 中，本宏包不应调整 \cs{hsize}；因而改为保存 \cs{hsize} 与
%  \cs{linewidth}，并将当前 \cs{hsize}、\cs{linewidth} 与保存值比较。
%  若二者相等，则（有望）可以相当有把握地假定当前处于普通
%  \env{addmargin} 环境中，需在排版题注前修正 \cs{hsize}。（长叹！）
%    \begin{macrocode}
  \newdimen\caption@addmargin@hsize
  \newdimen\caption@addmargin@linewidth
%    \end{macrocode}
%    \begin{macrocode}
  \providecommand*\caption@set@addmargin{%
    \caption@addmargin@hsize\hsize
    \caption@addmargin@linewidth\linewidth
    \l@addto@macro\caption@beginex@hook\caption@test@addmargin}%
%    \end{macrocode}
% \end{macro}
%
% \begin{macro}{\caption@test@addmargin}
% \changes{v3.6j}{2023/03/12}{加入本宏定义}
% 检测 \cs{hsize} 与 \cs{linewidth} 的当前值是否与 \env{addmargin}
% 环境中保存的值相同；若是，则把 \cs{hsize} 设为正确值。
%    \begin{macrocode}
  \providecommand\caption@test@addmargin{%
    \ifdim\hsize=\caption@addmargin@hsize
      \ifdim\linewidth=\caption@addmargin@linewidth
        \hsize=\linewidth
      \fi
    \fi}%
%    \end{macrocode}
% \end{macro}
%
%    \begin{macrocode}
  }{}%
%    \end{macrocode}
%
% \subsubsection{sidecap 宏包}
%
%    \begin{macrocode}
\caption@IfPackageLoaded{sidecap}[2003/06/06 v1.6f]{%
  \caption@setbool{needfreeze}{1}%
%    \end{macrocode}
%
% \begin{macro}{\SC@zfloat}
% \changes{v3.0b}{2004/05/16}{加入 \cs{captionsetup} 的局部定义}
% \changes{v3.0c}{2004/07/16}{缺陷修复 04-07-15：移除对 \cs{caption@star} 的检测}
% \changes{v3.1c}{2007/10/06}{现保存并恢复 \cs{caption}}
% \changes{v3.2a}{2011/08/17}{缺陷修复：现使用 \cs{caption@freeze}，而非 \cs{caption@freeze*}}
% \changes{v3.3}{2013/05/01}{以 \cs{caption@freezetype} 取代 \cs{caption@freeze} 的使用}
% \changes{v3.6i}{2022/07/10}{追加 \cs{ignorespaces}}
% 本宏在环境开始时调用，正好可借此对 |\caption| 与 |\captionsetup| 做若干
% 适配。
%    \begin{macrocode}
  \let\caption@ORI@SC@zfloat\SC@zfloat
  \def\SC@zfloat#1#2#3[#4]{%
%    \end{macrocode}
% 首先使用原始定义，但恢复 \cs{caption} 与 \cs{label}，使
%  \cs{caption@freeze} 与 \cs{caption@warmup} 能正确工作。
%    \begin{macrocode}
    \caption@ORI@SC@zfloat{#1}{#2}{#3}[#4]%
    \SC@RestoreCommands
%    \end{macrocode}
% 由于 sidecap 宏包在环境外部使用本宏包的 |\caption| 代码，常规的
%  |\caption|\-|setup| 无法生效。因此此处需要一个特殊版本，保存所给参数
%  列表，留待之后执行；此外还需使 |\caption*| 可用。
%    \begin{macrocode}
    \caption@freezetype{#2}%
%    \end{macrocode}
% sidecap 宏包用 |\ifx\label\SC@label| 检测是否正处于 |SC|\-|figure|
% 内部，因此此处重定义 |\SC@label|，使该检测仍然有效。
%    \begin{macrocode}
    \let\SC@label\label
%    \end{macrocode}
% 由于 |\SC@zfloat| 的原始定义以 |\begin{lrbox}| 结尾，而后者以
%  |\ignore|\-|spaces| 结尾，本重定义的 |\SC@zfloat| 也必须以
%  |\ignore|\-|spaces| 结尾。
%    \begin{macrocode}
    \ignorespaces}%
%    \end{macrocode}
%    \begin{macrocode}
  \providecommand*\SC@RestoreCommands{%
    \let\caption=\SC@orig@caption \let\label=\SC@orig@label}%
%    \end{macrocode}
% \end{macro}
%
% \begin{macro}{\endSC@FLOAT}
% \changes{v3.0i}{2006/01/12}{加入 \cs{@listdepth}\cs{z@}}
% \changes{v3.1}{2007/05/05}{加入 \cs{caption@setoptions}\arg{SCfloat} 的调用}
% 本宏在环境结束时调用；需在 \pkg{sidecap} 宏包实际排版题注之前完成本
% 宏包相关设置。
%    \begin{macrocode}
  \let\caption@ORI@endSC@FLOAT\endSC@FLOAT
  \def\endSC@FLOAT{%
%    \end{macrocode}
% 排版题注前需把边距设为零，任何额外边距在此都只会造成干扰。\par
% （无需关心题注位置，因为 sidecap 宏包已把 |\abovecaptionskip| 与
%  |\belowcaptionskip| 都设为零间距。）\par
% 此外，若设置了 |\SC@justify|，它会覆盖题注的对齐方式。|\SC@justify|
%  的用法在不同 \pkg{sidecap} 版本间存在差异：\par
%  \begin{tabular}{ll}
%    版本 1.4：& |\SC@justify| 未定义\\
%    版本 1.5：& 未设置时 |\SC@justify| 为 |\relax|\\
%    版本 1.6：& 未设置时 |\SC@justify| 为 |\@empty|\\
%  \end{tabular}
%    \begin{macrocode}
    \def\caption@setSC@justify{%
      \caption@clearmargin
        \ifx\SC@justify\@empty \else
          \let\caption@justification\SC@justify
          \let\SC@justify\@empty
        \fi}%
%    \end{macrocode}
% 让 |\endSC@FLOAT| 的原始定义改用本宏包的题注机制，而非其自身的实现。
%  \Note{此刻生效的是 \pkg{sidecap} 对 \cs{caption} 的定义，而非常规
%  定义！}
%    \begin{macrocode}
    \let\caption\SC@orig@caption
    \def\SC@orig@caption[##1]##2{%
      \caption@setSC@justify
%%%   \caption@setoptions{SC}%
      \caption@setoptions{SC\@captype}%
      \caption@defrost}%
%    \end{macrocode}
% 最后调用 |\endSC@FLOAT| 的原始定义。
%    \begin{macrocode}
    \caption@setSC@justify % 兼容模式
    \caption@prepare@defrost
    \caption@ORI@endSC@FLOAT}%
%    \end{macrocode}
% \end{macro}
%
%    \begin{macrocode}
}{}
%    \end{macrocode}
%
% \subsubsection{subfigure 宏包}
% \changes{v3.1}{2007/04/06}{加入对 subfigure 宏包的支持}
%
%    \begin{macrocode}
\caption@IfPackageLoaded{subfigure}[2002/01/23 v2.1]{%
%    \end{macrocode}
%
% \begin{macro}{\sf@ifpositiontop}
% \changes{v3.1a}{2007/09/15}{缺陷修复 07-09-15：加入对 \cs{@captype} 的检测}
% 若已加载 \pkg{subfigure} 宏包，把 |\sf@ifpositiontop| 映射到
%  |\iffiguretopcap| 或 |\iftabletopcap|，使 \pkg{subfigure}
%  \version{2.1} 的 \opt{figbotcap} 等选项仍能工作。
%    \begin{macrocode}
  \def\sf@ifpositiontop{%
    \ifx\@captype\@undefined
      \expandafter\@gobbletwo
    \else\ifx\@captype\relax
      \expandafter\expandafter\expandafter\@gobbletwo
    \else
      \expandafter\expandafter\expandafter\sf@if@position@top
    \fi\fi}
%    \end{macrocode}
%    \begin{macrocode}
  \def\sf@if@position@top{%
    \@ifundefined{if\@captype topcap}%
      {\@gobbletwo}%
      {\@nameuse{if\@captype topcap}%
         \expandafter\@firstoftwo
       \else
         \expandafter\@secondoftwo
       \fi}}
%    \end{macrocode}
% \end{macro}
%
%    \begin{macrocode}
}{}
%    \end{macrocode}
%
% \subsubsection{supertabular 与 xtab 宏包}
%
%    \begin{macrocode}
\caption@IfPackageLoaded{supertabular}[2002/07/19 v4.1e]{%
%    \end{macrocode}
%
% \begin{macro}{\tablecaption}
% \changes{v3.0e}{2005/05/05}{使 \cs{topcaption*} 与 \cs{bottomcaption*} 可用}
% \changes{v3.1h}{2008/04/01}{缺陷修复：带星形式不再递增表格计数器}
% 使 |\topcaption*| 与 |\bottomcaption*| 可用。
%    \begin{macrocode}
  \renewcommand*\tablecaption{%
    \caption@star
      {\refstepcounter{table}}%
      {\caption@dblarg{\@xtablecaption}}}%
%    \end{macrocode}
% \end{macro}
%
% \begin{macro}{\@xtablecaption}
% \changes{v3.1h}{2008/04/01}{使 \cs{nameref} 与 \cs{autoref} 可用}
% \changes{v3.1l}{2010/01/09}{适配 nameref 宏包当前版本}
% \changes{v3.3}{2018/12/26}{缺陷修复：补上遗漏的花括号}
% 使 |\nameref| 与 |\autoref| 可用。
%    \begin{macrocode}
  \let\caption@ORI@xtablecaption\@xtablecaption
  \long\def\@xtablecaption[#1]#2{%
    \caption@gettitle{#2}%
    \caption@ORI@xtablecaption[{#1}]{#2}}%
%    \end{macrocode}
% \end{macro}
%
% \begin{macro}{\ST@caption}
% \changes{v3.0a}{2004/01/23}{缺陷修复：补上遗漏的 \cs{par}}
% \changes{v3.1h}{2008/04/01}{缺陷修复：现设置 \cs{@currentlabelname}}
% 原始代码：
%  \begin{verbatim}
%  \long\def\ST@caption#1[#2]#3{\par%
%    \addcontentsline{\csname ext@#1\endcsname}{#1}%
%                    {\protect\numberline{%
%                        \csname the#1\endcsname}{\ignorespaces #2}}
%    \begingroup
%      \@parboxrestore
%      \normalsize
%      \if@topcaption \vskip -10\p@ \fi
%      \@makecaption{\csname fnum@#1\endcsname}{\ignorespaces #3}\par
%      \if@topcaption \vskip 10\p@ \fi
%    \endgroup}
%  \end{verbatim}%^^A
%    \begin{macrocode}
  \long\def\ST@caption#1[#2]#3{\par%
    \caption@settype{#1}%
    \caption@setoptions{supertabular}%
%    \end{macrocode}
%  |position=| 设置会被 \pkg{supertabular} 宏包覆盖：若使用了
%  |\topcaption|，位置自动为 |top|，否则为 |bottom|。
%    \begin{macrocode}
    \caption@setposition{\if@topcaption t\else b\fi}%
%    \end{macrocode}
%    \begin{macrocode}
    \caption@beginex{#1}{#2}{#3}%
      \caption@parboxrestore
      \caption@normalsize
      \@makecaption{\csname fnum@#1\endcsname}{\ignorespaces #3}\par
    \caption@end}%
%    \end{macrocode}
% \end{macro}
%
%    \begin{macrocode}
}{}
%    \end{macrocode}
%
% \changes{v3.1h}{2008/04/01}{加入对 xtab 宏包的支持}
%    \begin{macrocode}
\caption@IfPackageLoaded{xtab}[2000/04/09 v2.3]{%
%    \end{macrocode}
%
% \begin{macro}{\tablecaption}
% 使 |\topcaption*| 与 |\bottomcaption*| 可用。
%    \begin{macrocode}
  \renewcommand*\tablecaption{%
    \caption@star
      {\refstepcounter{table}}%
      {\caption@dblarg{\@xtablecaption}}}%
%    \end{macrocode}
% \end{macro}
%
% \begin{macro}{\@xtablecaption}
% \changes{v3.1l}{2010/01/09}{适配 nameref 宏包当前版本}
% \changes{v3.3}{2018/12/26}{缺陷修复：补上遗漏的花括号}
% 使 |\nameref| 与 |\autoref| 可用。
%    \begin{macrocode}
  \let\caption@ORI@xtablecaption\@xtablecaption
  \long\def\@xtablecaption[#1]#2{%
    \caption@gettitle{#2}%
    \caption@ORI@xtablecaption[{#1}]{#2}}%
%    \end{macrocode}
% \end{macro}
%
% \begin{macro}{\ST@caption}
% 原始代码：
%  \begin{verbatim}
%  \long\def\ST@caption#1[#2]#3{\par%
%    \@initisotab
%    \addcontentsline{\csname ext@#1\endcsname}{#1}%
%                    {\protect\numberline{%
%                      \csname the#1\endcsname}{\ignorespaces #2}}%
%    \begingroup
%      \@parboxrestore
%      \normalsize
%  %%  \if@topcaption \vskip -10\p@ \fi
%      \@makecaption{\csname fnum@#1\endcsname}{\ignorespaces #3}\par
%  %%  \if@topcaption \vskip 10\p@ \fi
%    \endgroup
%    \global\advance\ST@pageleft -\PWSTcapht
%    \ST@trace\tw@{Added caption. Space left for xtabular: \the\ST@pageleft}}
%  \end{verbatim}%^^A
%    \begin{macrocode}
  \long\def\ST@caption#1[#2]#3{\par%
    \caption@settype{#1}%
    \caption@setoptions{xtabular}%
%    \end{macrocode}
%    \begin{macrocode}
    \caption@set{position}{\if@topcaption t\else b\fi}%
%    \end{macrocode}
%    \begin{macrocode}
    \@initisotab
    \caption@beginex{#1}{#2}{#3}%
      \caption@parboxrestore
      \caption@normalsize
      \@makecaption{\csname fnum@#1\endcsname}{\ignorespaces #3}\par
    \caption@end
    \global\advance\ST@pageleft -\PWSTcapht
    \ST@trace\tw@{Added caption. Space left for xtabular: \the\ST@pageleft}}%
%    \end{macrocode}
% \end{macro}
%
%    \begin{macrocode}
}{}
%    \end{macrocode}
%
% \subsubsection{threeparttable 宏包}
% \changes{v3.1}{2007/05/08}{加入对 \pkg{threeparttable} 宏包的支持}
%
%    \begin{macrocode}
\caption@IfPackageLoaded{threeparttable}[2003/06/13 v3.0]{%
%    \end{macrocode}
%
% \begin{macro}{\threeparttable}
% \changes{v3.3}{2016/01/31}{缺陷修复：若 \cs{@captype} 已定义则使用它}
% \changes{v3.6e}{2022/04/17}{主体代码抽出至 \cs{caption@setthreeparttable}}
% 遗憾的是，使用 |\TPT@common| 时 |\@captype| 尚未设置，因此需改而重定义
%  |\threeparttable| 与 |\measuredfigure|。
%    \begin{macrocode}
  \let\caption@ORI@threeparttable\threeparttable
  \renewcommand*\threeparttable{%
    \caption@setthreeparttabletype{threepart}{table}%
    \caption@ORI@threeparttable}%
%    \end{macrocode}
% \end{macro}
%
% \begin{macro}{\measuredfigure}
% \changes{v3.3}{2016/01/31}{缺陷修复：若 \cs{@captype} 已定义则使用它}
% \changes{v3.6e}{2022/04/17}{主体代码抽出至 \cs{caption@setthreeparttable}}
% 此处同理……
%    \begin{macrocode}
  \let\caption@ORI@measuredfigure\measuredfigure
  \renewcommand*\measuredfigure{%
    \caption@setthreeparttabletype{measured}{figure}%
    \caption@ORI@measuredfigure}%
%    \end{macrocode}
% \end{macro}
%
% \begin{macro}{\caption@setthreeparttabletype}
% \changes{v3.6e}{2022/04/17}{由 \cs{threeparttable} 与 \cs{measuredfigure} 中抽出}
% \changes{v3.6e}{2022/04/17}{加入在 \env{subtable} 或 \env{subfigure} 内部使用的支持}
%  |\caption@setthreeparttabletype|\marg{附注类型}\marg{主类型}
%    \begin{macrocode}
  \newcommand*\caption@setthreeparttabletype[2]{%
    \caption@ifsubtype
      {\def\caption@threeparttabletype{#1sub#2}%
       \let\caption@setanchor\relax}%
      {\caption@settype{\@ifundefined{@captype}{#2}{\@captype}}%
       \def\caption@threeparttabletype{#1#2}}%
%%% \caption@setposition{auto}% ?
    \caption@clearmargin
    \caption@setoptions\caption@threeparttabletype
    \caption@setanchor}%
%    \end{macrocode}
% \end{macro}
%
% \begin{macro}{\TPT@caption}
% 原始代码：
%  \begin{verbatim}
%  \def\TPT@caption#1[#2]#3{\gdef\TPT@docapt
%   {\par\global\let\TPT@docapt\@undefined \TPT@LA@caption{#1}[{#2}]%
%     {\strut\ignorespaces#3\ifhmode\unskip\@finalstrut\strutbox\fi}}%
%   \ifx\TPT@hsize\@empty \let\label\TPT@gatherlabel \abovecaptionskip\z@skip
%   \else \TPT@docapt \fi \ignorespaces}
%  \end{verbatim}%^^A
%    \begin{macrocode}
  \def\TPT@caption#1[#2]#3{%
    \gdef\TPT@docapt{%
      \global\let\TPT@docapt\@undefined
      \caption@setautoposition\caption@TPT@position
      \TPT@LA@caption{#1}[{#2}]{#3}}%
    \ifx\TPT@hsize\@empty
      \let\label\TPT@gatherlabel % 缺陷：对 measuredfigure 无效
      \gdef\caption@TPT@position{t}%
      \g@addto@macro\TPT@docapt\caption@TPT@eatvskip
    \else
      \def\caption@TPT@position{b}%
      \TPT@docapt
    \fi
    \ignorespaces}%
%    \end{macrocode}
%    \begin{macrocode}
  %\newcommand*\caption@TPT@eatvskip{\vskip-.2\baselineskip}%
  \def\caption@TPT@eatvskip#1\vskip{#1\@tempdima=}%
%    \end{macrocode}
% \end{macro}
%
%    \begin{macrocode}
}{}
%    \end{macrocode}
%
% \subsubsection{wrapfig 宏包}
% \changes{v3.1}{2007/05/05}{加入对 \pkg{wrapfig} 宏包的支持}
% \changes{v3.2d}{2011/11/02}{重写对 \pkg{wrapfig} 宏包的支持}
%
%    \begin{macrocode}
\caption@IfPackageLoaded{wrapfig}[2003/01/31 v3.6]{%
%    \end{macrocode}
%
% \begin{macro}{\wrapfloat}
% \changes{v3.2d}{2011/11/02}{加入本宏的重定义}
% \changes{v3.6}{2020/12/28}{加入 \cs{@testopt} 的使用}
% 首先使 \pkg{wrapfig} 宏包在与 \pkg{float} 宏包共处时不再受二者加载
% 顺序影响。由于 \pkg{wrapfig} 代码中未使用 |\@float@set|\-|every|
% （它本应出现在 |\float@re|\-|style| 的重定义中，紧跟
%  |\@nameuse{fst@#1}| 之后），此处也不使用；尤其在搭配重新定义样式的
%  浮动体时，通常不会使用 |\wrap|\-|float|。
%    \begin{macrocode}
  \renewcommand*\wrapfloat[1]{%
    \def\@captype{#1}%
    \@ifundefined{fst@#1}{}{%
      \@nameuse{fst@#1}%
%     \@float@setevery{#1}%
      \def\WF@floatstyhook{\let\@currbox\WF@box 
         \global\setbox\WF@box\float@makebox{\wd\WF@box}}}%
    \@testopt\WF@wr{}}
%    \end{macrocode}
% \end{macro}
%
% \begin{macro}{\WF@rapt}
% \changes{v3.2d}{2011/11/02}{加入 \cs{WF@captionstyhook} 的使用}
% 原始代码：
% \begin{verbatim}
%   \def\WF@rapt[#1]#2{% 最后两个参数：#1 = 缩进量，#2 = 宽度
%     \gdef\WF@ovh{#1}% 保存缩进量，待得知 \width 后再用
%     \global\setbox\WF@box\vtop\bgroup \setlength\hsize{#2}%
%     \ifdim\hsize>\z@ \@parboxrestore \else
%     \setbox\z@\hbox\bgroup \let\wf@@caption\caption \let\caption\wf@caption
%     \ignorespaces \fi}
% \end{verbatim}%^^A
% 本代码额外加入 |\WF@captionstyhook|：
%    \begin{macrocode}
  \def\WF@rapt[#1]#2{% 最后两个参数：#1 = 缩进量，#2 = 宽度
    \gdef\WF@ovh{#1}% 保存缩进量，待得知 \width 后再用
    \global\setbox\WF@box\vtop\bgroup \setlength\hsize{#2}%
    \expandafter\WF@captionstyhook\expandafter{\@captype}% <= 新增
    \ifdim\hsize>\z@ \@parboxrestore \else
    \setbox\z@\hbox\bgroup \let\wf@@caption\caption \let\caption\wf@caption
    \ignorespaces \fi}%
%    \end{macrocode}
% \end{macro}
%
% \begin{macro}{\WF@captionstyhook}
% \changes{v3.2d}{2011/11/02}{新增本宏}
% \changes{v3.2e}{2011/11/10}{缺陷修复 11-11-09：修复对 float 宏包的支持}
% 在此放置本宏包的 \pkg{hyperref} 锚点、应用 `wrap' 选项等。由于
%  \pkg{wrapfig} 宏包未使用 |\@float@set|\-|every|，在此补上对
%  \pkg{float} 宏包的必要适配。
%    \begin{macrocode}
  \def\WF@captionstyhook#1{%
    \let\@captype\@undefined
    \@ifundefined{fst@#1}{}{\@float@setevery{#1}}%
    \caption@settype{#1}%
    \caption@clearmargin
%%% \caption@setoptions{wrap}%
    \caption@setoptions{wrap#1}%
    \caption@setanchor}%
%    \end{macrocode}
% \end{macro}
%
%    \begin{macrocode}
}{}
%    \end{macrocode}
%
% \subsection{调试相关内容}
%
% \begin{macro}{\caption@checkcounter}
% \changes{v3.6}{2020/12/28}{新增本宏}
% |\caption@checkref|\marg{计数器}\marg{预期值}\\
% 检测 |\value|\marg{计数器} 是否会得到 \meta{预期值}。
%    \begin{macrocode}
\newcommand*\caption@checkcounter[1]{%
  \caption@check{\the#1}{\@nameuse{the#1}}}
%    \end{macrocode}
%    \begin{macrocode}
\newcommand*\caption@check[3]{%
  \edef\caption@tempa{#2}%
  \edef\caption@tempb{#3}%
  \ifx\caption@tempa\caption@tempb
    \caption@Info{\string#1 = \caption@tempa}%
  \else
    \caption@Error{\string#1 = \caption@tempa\space <> \caption@tempb}%
  \fi}
%    \end{macrocode}
% \end{macro}
%
% \begin{macro}{\caption@checkref}
% \changes{v3.6}{2020/12/28}{新增本宏}
% |\caption@checkref|\marg{标记}\marg{预期值}\\
% 检测 |\ref|\marg{标记} 是否会得到 \meta{预期值}。
%    \begin{macrocode}
\newcommand*\caption@checkref[2]{%
  \@ifundefined{r@#1}{}%
    {\expandafter\expandafter\expandafter
       \caption@@checkref\csname r@#1\endcsname\@nil{#1}{#2}}}
%    \end{macrocode}
%    \begin{macrocode}
\def\caption@@checkref#1#2\@nil#3#4{%
  \caption@check{\ref{#3}}{#1}{#4}}
%    \end{macrocode}
% \end{macro}
%
% \subsection{就到这里，各位！}
%
% 执行用 |\After|\-|Caption|\-|Package| 声明的代码。
%
%    \begin{macrocode}
\caption@endtitles@code
\let\caption@endtitles@code\@undefined
\let\AfterCaptionPackage\@firstofone
%    \end{macrocode}
%
% \iffalse
%</package>
% \fi
%
% \iffalse
% --------------------------------------------------------------------------- %
% \fi
%
% \clearpage
% \Finale
%
\endinput

